% L’Etat et le pouvoir — Philosophie Tle Tech — Cours signé OnBuch+
% Programme officiel MINESEC. Compiler avec : tectonic N19-etat-pouvoir.tex
\newcommand{\DOCMATIERE}{Philosophie}
\newcommand{\DOCNIVEAU}{Tle Tech}
\newcommand{\DOCMODULE}{Philosophie – classe de Terminale technique}
\newcommand{\DOCLECON}{N19}
\newcommand{\DOCTITRE}{L’Etat et le pouvoir}
% OnBuch+ — préambule « cours signés » ENRICHI (compilé avec tectonic / XeLaTeX)
% Système visuel « Pop » d'OnBuch+ : fond crème (jamais blanc), bordures encre
% épaisses, ombres « dures » décalées, titres Archivo Black en capitales.
% Version MATHÉMATIQUES : graphes (pgfplots/tikz), boîtes pédagogiques
% (définition, propriété, méthode, attention, le savais-tu, exercice résolu,
% exercice, TP, bilan), page de garde et signature OnBuch+ ancrée en bas.
%
% Macros attendues dans main.tex AVANT \input{preamble} :
%   \DOCMATIERE  \DOCNIVEAU  \DOCMODULE  \DOCLECON (numéro)  \DOCTITRE
\documentclass[11pt]{article}

\usepackage{fontspec}
\usepackage[french]{babel}
\usepackage[a4paper,margin=2cm,top=3.4cm,bottom=2.8cm,headheight=1.4cm,footskip=1.3cm]{geometry}
\usepackage{amsmath,amssymb,amsfonts}
\usepackage{mathtools} % \xmapsto, \coloneqq... souvent utilisés par les modèles
\usepackage{cancel} % simplifications barrées (bilans de forces)
% \abs{z} (module, valeur absolue) et \norm{u} (norme), utilisés par les modèles
\providecommand{\abs}{}\let\abs\relax\DeclarePairedDelimiter{\abs}{\lvert}{\rvert}
\providecommand{\norm}{}\let\norm\relax\DeclarePairedDelimiter{\norm}{\lVert}{\rVert}
\usepackage[table]{xcolor} % [table] : fournit aussi \rowcolors (alternance de lignes)
\usepackage{pagecolor}
\usepackage{tikz}
\usetikzlibrary{arrows.meta,positioning,calc,shapes.geometric,shapes.misc,shapes.arrows,decorations.pathmorphing,decorations.pathreplacing,decorations.markings,patterns,shadows,mindmap,trees,fit,backgrounds,matrix,intersections,angles,quotes}
\usepackage{pgfplots}
\usepackage[version=4]{mhchem} % équations nucléaires \ce{^{235}_{92}U}
\usepackage[european,siunitx]{circuitikz} % schémas électriques
\pgfplotsset{compat=1.18}
\usepgfplotslibrary{fillbetween,groupplots} % aires entre courbes (calcul intégral)
\usepackage{enumitem}
\usepackage{fancyhdr}
% [most] charge toutes les bibliothèques tcolorbox (listings, minted, raster,
% documentation...) et gonfle inutilement la mémoire TeX sur un document long
% (~300 boîtes) : on ne charge que ce qui est réellement utilisé.
\usepackage[skins,breakable]{tcolorbox}
\usepackage{graphicx}
\usepackage{colortbl}
\usepackage{booktabs}
\usepackage{tabularx}
\usepackage{multirow}
\usepackage{diagbox} % case d'en-tête diagonale des tableaux à double entrée
% Colonnes extensibles centrée / gauche / droite (C, L, R), souvent utilisées par les modèles
\newcolumntype{C}{>{\centering\arraybackslash}X}
\newcolumntype{L}{>{\raggedright\arraybackslash}X}
\newcolumntype{R}{>{\raggedleft\arraybackslash}X}
\usepackage{array}
\usepackage{multicol}
\usepackage{siunitx}
% parse-numbers=false : accepte aussi \SI{2,5 \times 10^{-3}}{...}
\sisetup{locale=FR,output-decimal-marker={,},group-minimum-digits=4,parse-numbers=false}
\DeclareSIUnit\ohm{\ensuremath{\symup{\Omega}}} % U+2126 absent de la police
\DeclareSIUnit\division{div} % réglages d'oscilloscope (ms/div, V/div)
\DeclareSIUnit\tr{tr} % tours (vitesse de rotation, ex. \SI{1500}{\tr\per\minute})
\DeclareSIUnit\molar{M} % concentration molaire (ex. \SI{3}{\molar}, \SI{1}{\micro\molar})
\DeclareSIUnit\an{an} % année (le modèle confond parfois avec \year, un compteur TeX) — ex. \SI{5}{\kilo\meter\per\an}
\DeclareSIUnit\FCFA{FCFA} % franc CFA (ex. \SI{500}{\FCFA\per\kilo\gram})
\DeclareSIUnit\degreeBrix{\textdegree Brix} % teneur en sucre d'un sirop/jus (réfractométrie)
\DeclareSIUnit\month{mois} % le modèle utilise parfois \month, un compteur TeX, comme unité de durée
\DeclareSIUnit\microliter{\micro\liter} % le modèle écrit parfois \microliter en un seul token
\DeclareSIUnit\million{millions} % dénombrement (ex. \SI{15}{\million\per\mL} spermatozoïdes)
\usepackage{qrcode}
% \degree : commande fréquemment utilisée par les modèles pour un angle ou une
% température en dehors de \SI{}{} (non fournie par les paquets chargés).
\providecommand{\degree}{^{\circ}}
% \qed (fin de démonstration), utilisé par les modèles sans amsthm
\providecommand{\qed}{\hfill$\square$}
% \barre : double barre des valeurs interdites dans un tableau de variations (tkz-tab)
\providecommand{\barre}{\ensuremath{\|}}
% environnement proof (amsthm non chargé) utilisé par les modèles
\ifdefined\proof\else\newenvironment{proof}[1][Démonstration]{\par\noindent\textit{#1.}\ }{\qed\par}\fi
\usepackage{lastpage}
\usepackage{hyperref}
\hypersetup{hidelinks}

% ============================================================
% Palette « Pop » OnBuch+ — mêmes hex que la classe P (plus_ui.dart)
% ============================================================
\definecolor{poporange}{HTML}{F59321}
\definecolor{popdark}{HTML}{E07A0C}
\definecolor{popcream}{HTML}{FFF6E8}
\definecolor{popink}{HTML}{1C1714}
\definecolor{poppurple}{HTML}{7A5AE0}
\definecolor{popgreen}{HTML}{2E9E6A}
\definecolor{poppink}{HTML}{E0457B}
\definecolor{popblue}{HTML}{2F7DE1}
\definecolor{popgold}{HTML}{C98A12}
\definecolor{popmuted}{HTML}{8A7F76}
\definecolor{popline}{HTML}{EADFD2}
\definecolor{popgray}{HTML}{8A7F76}
\colorlet{popgrey}{popgray}
% Alias tolérants (le modèle nomme parfois les couleurs différemment)
\colorlet{popwhite}{white}
\colorlet{popblack}{popink}
\colorlet{popred}{poppink}
\colorlet{popyellow}{popgold}
% Teintes claires (fonds de tableaux/encadrés) — le modèle nomme parfois
% les couleurs avec un suffixe L (ex. popblueL) sans les avoir définies.
\colorlet{poporangeL}{poporange!20}
\colorlet{popdarkL}{popdark!20}
\colorlet{poppurpleL}{poppurple!20}
\colorlet{popgreenL}{popgreen!20}
\colorlet{poppinkL}{poppink!20}
\colorlet{popblueL}{popblue!20}
\colorlet{popgoldL}{popgold!20}
\colorlet{popredL}{popred!20}
\colorlet{popgrayL}{popgray!20}
% Teintes claires pour fonds de tableaux / zones
\colorlet{popblueL}{popblue!10!white}
\colorlet{poporangeL}{poporange!14!white}
\colorlet{popgreenL}{popgreen!12!white}
\colorlet{poppurpleL}{poppurple!12!white}
\colorlet{poppinkL}{poppink!12!white}
\colorlet{popgoldL}{popgold!14!white}

\pagecolor{popcream}

% ============================================================
% Typographie
% ============================================================
\defaultfontfeatures{Ligatures=TeX}
\setmainfont{PlusJakartaSans-Medium.ttf}[Path=../../../fonts/,BoldFont=PlusJakartaSans-Bold.ttf]
% Maths en sans-serif (Fira Math) avec chiffres et lettres droites de Plus
% Jakarta, pour que les formules se fondent dans le texte.
\usepackage[math-style=ISO,bold-style=ISO]{unicode-math}
\setmathfont{FiraMath-Regular.otf}[Path=../../../fonts/]
\setmathfont{PlusJakartaSans-Medium.ttf}[Path=../../../fonts/,range={up/num,up/{Latin,latin},\mathup}]
\usepackage{icomma} % virgule décimale sans espace en mode math
% Caractères mathématiques Unicode absents des polices de texte (Plus Jakarta,
% Archivo Black) : rendus par leur équivalent mathématique, en texte comme en
% formule — sinon ils s'affichent en carré vide (ex. « Arithmétique dans ℤ »).
\usepackage{newunicodechar}
\newunicodechar{ℝ}{\ensuremath{\mathbb{R}}}
\newunicodechar{ℤ}{\ensuremath{\mathbb{Z}}}
\newunicodechar{ℂ}{\ensuremath{\mathbb{C}}}
\newunicodechar{ℕ}{\ensuremath{\mathbb{N}}}
\newunicodechar{ℚ}{\ensuremath{\mathbb{Q}}}
\newunicodechar{𝔻}{\ensuremath{\mathbb{D}}}
\newunicodechar{α}{\ensuremath{\alpha}}
\newunicodechar{β}{\ensuremath{\beta}}
\newunicodechar{γ}{\ensuremath{\gamma}}
\newunicodechar{δ}{\ensuremath{\delta}}
\newunicodechar{ε}{\ensuremath{\varepsilon}}
\newunicodechar{θ}{\ensuremath{\theta}}
\newunicodechar{λ}{\ensuremath{\lambda}}
\newunicodechar{μ}{\ensuremath{\mu}}
\newunicodechar{σ}{\ensuremath{\sigma}}
\newunicodechar{τ}{\ensuremath{\tau}}
\newunicodechar{φ}{\ensuremath{\varphi}}
\newunicodechar{ψ}{\ensuremath{\psi}}
\newunicodechar{ω}{\ensuremath{\omega}}
\newunicodechar{Σ}{\ensuremath{\Sigma}}
% majuscules produites par \MakeUppercase dans les titres (α-aminés -> Α)
\newunicodechar{Α}{\ensuremath{\alpha}}
\newunicodechar{Β}{\ensuremath{\beta}}
\newunicodechar{Γ}{\ensuremath{\Gamma}}
\newunicodechar{Δ}{\ensuremath{\Delta}}
\newunicodechar{Ε}{\ensuremath{\varepsilon}}
\newunicodechar{Θ}{\ensuremath{\Theta}}
\newunicodechar{Λ}{\ensuremath{\Lambda}}
\newunicodechar{Μ}{\ensuremath{\mu}}
\newunicodechar{Τ}{\ensuremath{\tau}}
\newunicodechar{Φ}{\ensuremath{\Phi}}
\newunicodechar{Ψ}{\ensuremath{\Psi}}
\newunicodechar{Ω}{\ensuremath{\Omega}}
\newunicodechar{↦}{\ensuremath{\mapsto}}
\newunicodechar{∈}{\ensuremath{\in}}
\newunicodechar{∉}{\ensuremath{\notin}}
\newunicodechar{∧}{\ensuremath{\wedge}}
\newunicodechar{⇔}{\ensuremath{\Leftrightarrow}}
\newunicodechar{⟺}{\ensuremath{\Longleftrightarrow}}
\newunicodechar{⇒}{\ensuremath{\Rightarrow}}
\newunicodechar{⟹}{\ensuremath{\Longrightarrow}}
\newunicodechar{∘}{\ensuremath{\circ}}
\newunicodechar{∪}{\ensuremath{\cup}}
\newunicodechar{∩}{\ensuremath{\cap}}
\newunicodechar{≡}{\ensuremath{\equiv}}
\newunicodechar{⊂}{\ensuremath{\subset}}
\newunicodechar{⊥}{\ensuremath{\perp}}
\newunicodechar{∃}{\ensuremath{\exists}}
\newunicodechar{∀}{\ensuremath{\forall}}
\newunicodechar{∛}{\ensuremath{\sqrt[3]{\,}}}
\newunicodechar{∜}{\ensuremath{\sqrt[4]{\,}}}
\newunicodechar{⃗}{\ensuremath{{}^{\to}}}
\newunicodechar{ⁿ}{\ifmmode^{n}\else\textsuperscript{\ensuremath{n}}\fi}
\newunicodechar{ˣ}{\ifmmode^{x}\else\textsuperscript{\ensuremath{x}}\fi}
\newunicodechar{ᵇ}{\ifmmode^{b}\else\textsuperscript{\ensuremath{b}}\fi}
\newunicodechar{ʳ}{\ifmmode^{r}\else\textsuperscript{\ensuremath{r}}\fi}
\newunicodechar{ᵘ}{\ifmmode^{u}\else\textsuperscript{\ensuremath{u}}\fi}
\newunicodechar{ʸ}{\ifmmode^{y}\else\textsuperscript{\ensuremath{y}}\fi}
\newunicodechar{ᵃ}{\ifmmode^{a}\else\textsuperscript{\ensuremath{a}}\fi}
\newunicodechar{ᵉ}{\ifmmode^{e}\else\textsuperscript{\ensuremath{e}}\fi}
\newunicodechar{ᵏ}{\ifmmode^{k}\else\textsuperscript{\ensuremath{k}}\fi}
\newunicodechar{ᵖ}{\ifmmode^{p}\else\textsuperscript{\ensuremath{p}}\fi}
\newunicodechar{ᴺ}{\ifmmode^{N}\else\textsuperscript{\ensuremath{N}}\fi}
\newunicodechar{ᶜ}{\ifmmode^{c}\else\textsuperscript{\ensuremath{c}}\fi}
\newunicodechar{ʰ}{\ifmmode^{h}\else\textsuperscript{\ensuremath{h}}\fi}
\newunicodechar{ᵠ}{\ifmmode^{q}\else\textsuperscript{\ensuremath{q}}\fi}
\newunicodechar{ₙ}{\ifmmode_{n}\else\textsubscript{\ensuremath{n}}\fi}
\newunicodechar{ₐ}{\ifmmode_{a}\else\textsubscript{\ensuremath{a}}\fi}
\newunicodechar{ᵢ}{\ifmmode_{i}\else\textsubscript{\ensuremath{i}}\fi}
\newunicodechar{ᵣ}{\ifmmode_{r}\else\textsubscript{\ensuremath{r}}\fi}
\newunicodechar{ₖ}{\ifmmode_{k}\else\textsubscript{\ensuremath{k}}\fi}
\newunicodechar{ᵦ}{\ifmmode_{\beta}\else\textsubscript{\ensuremath{\beta}}\fi}
\newunicodechar{ₓ}{\ifmmode_{x}\else\textsubscript{\ensuremath{x}}\fi}
\newfontfamily\popdisplay{ArchivoBlack-Regular.ttf}[Path=../../../fonts/]
\newfontfamily\poplogo{SpaceGrotesk-Bold.ttf}[Path=../../../fonts/]
\newfontfamily\popsemi{PlusJakartaSans-SemiBold.ttf}[Path=../../../fonts/]
\newfontfamily\popxbold{PlusJakartaSans-ExtraBold.ttf}[Path=../../../fonts/]
\newfontfamily\popxboldls{PlusJakartaSans-ExtraBold.ttf}[Path=../../../fonts/,LetterSpace=3.0]

\setlist[enumerate]{leftmargin=1.7em,itemsep=5pt,topsep=4pt}
\setlist[itemize]{leftmargin=1.7em,itemsep=4pt,topsep=4pt,label={\color{poporange}\textbullet}}
\setlength{\parindent}{0pt}
\setlength{\parskip}{5pt}
\renewcommand{\arraystretch}{1.25}

% pgfplots : style Pop par défaut
\pgfplotsset{
  every axis/.append style={axis line style={popink,line width=1pt},tick style={popink},
    grid style={popline,line width=0.5pt},label style={font=\small},tick label style={font=\footnotesize}},
  popcurve/.style={poporange,line width=1.8pt,no markers,smooth},
}
% Même style disponible dans un \draw TikZ simple (hors pgfplots)
\tikzset{popcurve/.style={poporange,line width=1.8pt}}
\tikzset{popgrid/.style={popline,very thin}}
\colorlet{popcurve}{poporange} % le nom du style est parfois utilisé comme couleur

% ============================================================
% Pill / squiggle
% ============================================================
\newcommand{\poppill}[3]{%
  \tikz[baseline=-0.55ex]\node[draw=popink,line width=1.2pt,rounded corners=2.6pt,%
    fill=#2,inner xsep=6.5pt,inner ysep=3pt]{%
    \popxboldls\fontsize{7}{7}\selectfont\color{#3}\MakeUppercase{#1}};%
}
\newcommand{\popsquiggle}[1]{%
  \tikz[baseline=-0.4ex]{
    \draw[#1,line width=2.6pt,line cap=round]
      plot [smooth] coordinates {(0,0.09) (0.34,0.01) (0.68,0.21) (1.02,0.01) (1.36,0.10)};
  }%
}
% Mot-clé surligné (marqueur orange sous le texte)
\newcommand{\cle}[1]{{\popxbold\color{popdark}#1}}

% ============================================================
% En-tête / pied
% ============================================================
\pagestyle{fancy}
\fancyhf{}
\renewcommand{\headrulewidth}{0pt}
\renewcommand{\footrulewidth}{0pt}
\lhead{\raisebox{-0.15ex}{\poplogo\fontsize{13}{13}\selectfont\color{popink}OnBuch\color{poporange}+}}
\rhead{\poppill{\DOCMATIERE\ \textbullet\ \DOCNIVEAU\ \textbullet\ Leçon \DOCLECON}{white}{popink}}
\lfoot{\popxbold\fontsize{7.6}{7.6}\selectfont\color{popmuted}\MakeUppercase{Cours signé OnBuch+}\ \textbullet\ {\popsemi\DOCTITRE}}
\rfoot{\tikz[baseline=-0.6ex]\node[rounded corners=8.5pt,fill=popink,minimum height=17pt,minimum width=17pt,inner xsep=5pt,inner sep=0]{\popxbold\fontsize{8}{8}\selectfont\color{popcream}\thepage\,/\,\pageref*{LastPage}};}

% ============================================================
% Titres
% ============================================================
\newcommand{\chaptitre}[2]{%
  \begin{tcolorbox}[enhanced,colback=poporange,colframe=popink,boxrule=2.6pt,arc=11pt,
      drop shadow=popink,left=15pt,right=15pt,top=12pt,bottom=11pt,before skip=3pt,after skip=18pt]
    \poppill{#2}{popink}{popcream}\par\vspace{9pt}
    {\raggedright\hyphenpenalty=10000\exhyphenpenalty=10000\popdisplay\fontsize{22}{24}\selectfont\color{popcream}\MakeUppercase{#1}\par}
  \end{tcolorbox}
}
% Section numérotée : pastille orange avec le numéro + titre Archivo
\newcounter{popsec}
\newcommand{\coursec}[1]{%
  \stepcounter{popsec}\setcounter{popsub}{0}%
  \par\vspace{16pt}\noindent
  \begin{minipage}{\linewidth}
  \tikz[baseline=-0.7ex]\node[draw=popink,line width=1.6pt,fill=poporange,rounded corners=4pt,minimum size=22pt,inner sep=0]{\popdisplay\fontsize{12}{12}\selectfont\color{popcream}\thepopsec};%
  \hspace{8pt}{\popdisplay\fontsize{15}{17}\selectfont\color{popink}\MakeUppercase{#1}}\par
  \vspace{4pt}\noindent{\color{poporange}\rule{36pt}{3.6pt}}
  \end{minipage}\par\nopagebreak\vspace{8pt}
}
\newcounter{popsub}
\newcommand{\courssub}[1]{%
  \stepcounter{popsub}%
  \par\vspace{9pt}\noindent{\popxbold\fontsize{11.5}{13}\selectfont\color{popink}{\color{poporange}\thepopsec.\thepopsub}\hspace{6pt}#1}\par\nopagebreak\vspace{4pt}
}

% ============================================================
% Boîtes pédagogiques — même recette : fond blanc, bordure épaisse colorée,
% ombre dure encre, badge plein. Titre optionnel [#1].
% ============================================================
\tcbset{popbase/.style={breakable,enhanced,colback=white,boxrule=2.3pt,arc=9pt,
  shadow={3pt}{-3pt}{0pt}{popink},left=14pt,right=14pt,top=11pt,bottom=11pt,before skip=11pt,after skip=15pt,
  fontupper=\popsemi\fontsize{10.6}{14.5}\selectfont\color{popink},
  fontlower=\popsemi\fontsize{10.6}{14.5}\selectfont\color{popink}}}
% Coche dessinée (le glyphe ✓ n'existe pas dans les polices du document)
\newcommand{\popcheck}[1]{\tikz[baseline=-0.5ex]\draw[#1,line width=1.6pt,line cap=round,line join=round](-0.13,0.0)--(-0.04,-0.09)--(0.13,0.1);}
\providecommand{\coche}{\popcheck{popgreen}}
\providecommand{\resultat}[1]{\par\smallskip\noindent\fcolorbox{popgreen}{popgreenL}{\parbox{\dimexpr\linewidth-2\fboxsep-2\fboxrule\relax}{\textbf{Résultat.} #1}}\par\smallskip}
\newcommand{\popboxhead}[3]{\poppill{#1}{#2}{white}\ifx\relax#3\relax\else\hspace{6pt}{\popxbold\fontsize{10.6}{12}\selectfont\color{popink}#3}\fi\par\vspace{7pt}}

\newtcolorbox{aretenir}[1][]{popbase,colframe=popgreen,before upper={\popboxhead{À retenir}{popgreen}{#1}}}
\newtcolorbox{exemplebox}[1][]{popbase,colframe=poporange,before upper={\popboxhead{Exemple}{poporange}{#1}}}
\newtcolorbox{definition}[1][]{popbase,colframe=popblue,before upper={\popboxhead{Définition}{popblue}{#1}}}
\newtcolorbox{propriete}[1][]{popbase,colframe=poppurple,before upper={\popboxhead{Propriété}{poppurple}{#1}}}
\newtcolorbox{methode}[1][]{popbase,colframe=popgold,before upper={\popboxhead{Méthode}{popgold}{#1}}}
\newtcolorbox{attention}[1][]{popbase,colframe=poppink,before upper={\popboxhead{Attention}{poppink}{#1}}}
\newtcolorbox{experience}[1][]{popbase,colframe=popblue,colback=popblueL,before upper={\popboxhead{Expérience}{popblue}{#1}}}
% « Le savais-tu ? » : carte violette pleine (contexte, culture, Cameroun)
\newtcolorbox{savaistu}[1][]{popbase,colback=poppurple,colframe=popink,
  fontupper=\popsemi\fontsize{10.4}{14.2}\selectfont\color{white},
  before upper={\poppill{Le savais-tu\,?}{white}{poppurple}\ifx\relax#1\relax\else\hspace{6pt}{\popxbold\color{white}#1}\fi\par\vspace{7pt}}}
% Exercice résolu : énoncé en haut, \tcblower puis la solution sur fond crème
\newcounter{popexo}\newcounter{popexr}
\newtcolorbox{exoresolu}[1][]{popbase,colframe=poporange,colbacklower=popcream,
  segmentation style={popink,line width=1.2pt,dashed},
  before upper={\stepcounter{popexr}\popboxhead{Exercice résolu \thepopexr}{poporange}{#1}},
  before lower={\poppill{Solution}{popink}{popcream}\par\vspace{6pt}}}
% Exercice (non résolu en place) — la correction est dans la partie Corrigés
\newtcolorbox{exercice}[1][]{popbase,colframe=popink,
  before upper={\stepcounter{popexo}\popboxhead{Exercice \thepopexo}{popink}{#1}}}
% Corrigé
\newtcolorbox{corrige}[1][]{popbase,colframe=popgreen,colback=popgreenL,
  before upper={\popboxhead{Corrigé}{popgreen}{#1}}}
% Objectifs / prérequis / bilan
\newtcolorbox{objectifs}{popbase,colframe=popink,colback=poporangeL,
  before upper={\popboxhead{Objectifs de la leçon}{popink}{}}}
\newtcolorbox{prerequis}{popbase,colframe=popink,colback=popblueL,
  before upper={\popboxhead{Prérequis}{popblue}{}}}
\newtcolorbox{bilan}{popbase,colframe=popink,colback=white,boxrule=2.8pt,
  before upper={\popboxhead{Fiche bilan}{poporange}{L'essentiel à connaître pour l'examen}}}
% Figure encadrée avec légende
\newtcolorbox{popfigure}[1][]{enhanced,breakable=false,colback=white,colframe=popline,boxrule=1.4pt,arc=8pt,
  left=10pt,right=10pt,top=10pt,bottom=8pt,before skip=10pt,after skip=12pt,halign=center,
  lower separated=false,halign lower=center,
  fontlower=\popsemi\fontsize{9}{11.5}\selectfont\color{popmuted},
  #1}
\newcommand{\legende}[1]{\par\vspace{6pt}{\popsemi\fontsize{9}{11.5}\selectfont\color{popmuted}#1}}

% ============================================================
% Signature OnBuch+
% ============================================================
\newcommand{\poplogoinline}[1]{{\poplogo\fontsize{#1}{#1}\selectfont\color{popink}OnBuch\color{poporange}+}}
\newcommand{\signatureonbuch}{%
  \begin{tcolorbox}[enhanced,colback=popink,colframe=popink,boxrule=2.2pt,arc=10pt,
      drop shadow=poporange,left=14pt,right=14pt,top=10pt,bottom=10pt,before skip=0pt,after skip=0pt]
    \tikz[baseline=-0.7ex]\node[circle,fill=popgreen,draw=popcream,line width=1.3pt,minimum size=22pt,inner sep=0]{\popcheck{white}};%
    \hspace{9pt}\begin{minipage}[c]{0.62\linewidth}
      {\popxbold\fontsize{12}{13}\selectfont\color{popcream}Signé\ \textbullet\ {\poplogo OnBuch\color{poporange}+}}\par\vspace{2pt}
      {\raggedright\popsemi\fontsize{8}{10.5}\selectfont\color{popline}\MakeUppercase{Équipe pédagogique OnBuch+}\\\MakeUppercase{Conforme au programme officiel MINESEC}\par}
    \end{minipage}\hfill
    {\poplogo\fontsize{22}{22}\selectfont\color{popcream}OnBuch\color{poporange}+}
  \end{tcolorbox}%
}

% ============================================================
% Page de garde
% #1 titre · #2 matière · #3 niveau · #4 module · #5 n° leçon · #6 durée ·
% #7 description / objectifs (texte ou itemize)
% ============================================================
\newcommand{\pagegarde}[7]{%
  \thispagestyle{empty}
  \newgeometry{a4paper,margin=1.7cm,top=1.6cm,bottom=1.4cm}%
  \begin{tikzpicture}[remember picture,overlay]
    \fill[poporange!18] ([xshift=-3.2cm,yshift=-3.6cm]current page.north east) circle (4.6cm);
    \fill[poppurple!14] ([xshift=2.4cm,yshift=7.6cm]current page.south west) circle (3.8cm);
  \end{tikzpicture}%
  {\poplogo\fontsize{30}{30}\selectfont\color{popink}OnBuch\color{poporange}+}\hfill
  \raisebox{0.6ex}{\poppill{Cours signé \textbullet\ Premium}{popink}{popcream}}\par
  \vspace{1pt}\popsquiggle{poppurple}\par
  \vspace{8pt}
  \begin{tcolorbox}[enhanced,colback=poporange,colframe=popink,boxrule=2.8pt,arc=13pt,
      drop shadow=popink,left=18pt,right=18pt,top=12pt,bottom=13pt,after skip=10pt]
    \poppill{#2 \textbullet\ #3}{popink}{popcream}\hspace{5pt}\poppill{#4}{white}{popink}\par\vspace{8pt}
    {\popdisplay\fontsize{14}{14}\selectfont\color{popink}LEÇON #5}\par\vspace{4pt}
    {\raggedright\hyphenpenalty=10000\exhyphenpenalty=10000\popdisplay\fontsize{25}{27}\selectfont\color{popcream}\MakeUppercase{#1}\par}
    \vspace{8pt}
    {\popxbold\fontsize{9.5}{9.5}\selectfont\color{popink}\MakeUppercase{Durée conseillée : #6}}
  \end{tcolorbox}
  \begin{tcolorbox}[enhanced,colback=white,colframe=popink,boxrule=2.2pt,arc=10pt,
      drop shadow=popink,left=16pt,right=16pt,top=10pt,bottom=10pt,after skip=8pt]
    \poppill{Dans cette leçon}{poppurple}{white}\par\vspace{5pt}
    \popsemi\fontsize{9.3}{12.6}\selectfont\color{popink}#7
  \end{tcolorbox}
  \noindent\begin{minipage}[t]{0.487\linewidth}
    \begin{tcolorbox}[enhanced,colback=popgreen,colframe=popink,boxrule=2.2pt,arc=10pt,
        drop shadow=popink,left=11pt,right=11pt,top=10pt,bottom=10pt,height=3.1cm,valign=center]
      \tcbox[colback=white,colframe=popink,boxrule=1.3pt,arc=5pt,boxsep=0pt,nobeforeafter,
        left=3pt,right=3pt,top=3pt,bottom=3pt,valign=center]{\qrcode[height=1.9cm]{https://play.google.com/store/apps/details?id=com.luvvix.onbuch&hl=fr}}%
      \hspace{8pt}\begin{minipage}[c]{0.5\linewidth}
        {\popdisplay\fontsize{11}{12.5}\selectfont\color{white}\MakeUppercase{Télécharge OnBuch}}\par\vspace{4pt}
        {\raggedright\popsemi\fontsize{8.4}{11}\selectfont\color{white}Gratuit \textbullet\ Google Play\\Tuteur IA, annales, concours, résultats\par}
      \end{minipage}
    \end{tcolorbox}
  \end{minipage}\hfill
  \begin{minipage}[t]{0.487\linewidth}
    \begin{tcolorbox}[enhanced,colback=poppurple,colframe=popink,boxrule=2.2pt,arc=10pt,
        drop shadow=popink,left=11pt,right=11pt,top=10pt,bottom=10pt,height=3.1cm,valign=center]
      \tikz[baseline=-0.7ex]\node[circle,fill=white,draw=popink,line width=1.3pt,minimum size=1.6cm,inner sep=0]{\popdisplay\fontsize{24}{24}\selectfont\color{poppurple}+};%
      \hspace{8pt}\begin{minipage}[c]{0.6\linewidth}
        {\popdisplay\fontsize{11}{12.5}\selectfont\color{white}\MakeUppercase{Passe à OnBuch+}}\par\vspace{4pt}
        {\raggedright\popsemi\fontsize{8.4}{11}\selectfont\color{white}Tous les cours signés, Tuteur IA illimité, fiches PDF — onglet OnBuch+ de l'app\par}
      \end{minipage}
    \end{tcolorbox}
  \end{minipage}
  \vspace{8pt}
  \popcontenu
  \vfill
  \signatureonbuch
  \restoregeometry
  \newpage
}

% Rangée « Ce cours contient » (page de garde)
\newcommand{\poptile}[3]{% #1 couleur #2 gros texte #3 légende
  \begin{tcolorbox}[enhanced,colback=#1,colframe=popink,boxrule=2pt,arc=9pt,shadow={2.5pt}{-2.5pt}{0pt}{popink},
      left=6pt,right=6pt,top=9pt,bottom=9pt,halign=center,valign=center,height=1.75cm,width=\linewidth,nobeforeafter]
    {\popdisplay\fontsize{12.5}{13}\selectfont\color{white}\MakeUppercase{#2}}\par\vspace{3pt}
    {\centering\popsemi\fontsize{7.8}{9.5}\selectfont\color{white}#3\par}
  \end{tcolorbox}}
\newcommand{\popcontenu}{%
  \par\noindent{\popxboldls\fontsize{7.5}{7.5}\selectfont\color{popmuted}CE COURS SIGNÉ CONTIENT}\par\vspace{7pt}
  \noindent\begin{minipage}[t]{0.235\linewidth}\poptile{popblue}{Cours}{complet, illustré, pas à pas}\end{minipage}\hfill
  \begin{minipage}[t]{0.235\linewidth}\poptile{poporange}{Méthodes}{et exercices résolus}\end{minipage}\hfill
  \begin{minipage}[t]{0.235\linewidth}\poptile{poppink}{Exercices}{type examen + corrigés}\end{minipage}\hfill
  \begin{minipage}[t]{0.235\linewidth}\poptile{popgreen}{Fiche bilan}{l'essentiel à retenir}\end{minipage}\par}

% Sommaire
\newtcolorbox{sommaire}{enhanced,colback=white,colframe=popink,boxrule=2.2pt,arc=10pt,
  drop shadow=popink,left=18pt,right=18pt,top=13pt,bottom=13pt,before skip=6pt,after skip=20pt,
  before upper={\poppill{Sommaire}{poppurple}{white}\par\vspace{9pt}\popsemi\fontsize{10.6}{14.5}\selectfont\color{popink}}}

% Signature de fin de cours — poussée tout en bas de la dernière page
\newcommand{\findecours}{%
  \par\vfill\nopagebreak
  \begin{center}\popsquiggle{poporange}\end{center}\vspace{4pt}
  \signatureonbuch
}
\AtBeginDocument{\def\llbracket{\mathopen{[\mkern-3.5mu[}}\def\rrbracket{\mathclose{]\mkern-3.5mu]}}} % intervalles d'entiers (glyphe ⟦ absent de la police)
% Glyphes absents de Fira Math : redéfinis à partir de glyphes disponibles
\AtBeginDocument{%
  \def\setminus{\mathbin{\backslash}}\let\smallsetminus\setminus
  \def\vdots{\mathord{\vbox{\baselineskip3.2pt\lineskiplimit0pt\kern4pt\hbox{.}\hbox{.}\hbox{.}}}}%
  \def\ddots{\mathinner{\mkern1mu\raise6.5pt\hbox{.}\mkern2mu\raise3.6pt\hbox{.}\mkern2mu\raise0.7pt\hbox{.}\mkern1mu}}%
  \def\triangle{\mathord{\scalebox{0.9}{$\symup{\Delta}$}}}%
  \def\nabla{\mathord{\raisebox{\depth}{\scalebox{1}[-1]{$\symup{\Delta}$}}}}%
  \def\checkmark{\popcheck{popdark}}%
  \def\star{\mathbin{\ast}}\def\bigstar{\mathord{\ast}}%
  \def\diamond{\mathbin{\scalebox{0.6}{$\Diamond$}}}%
  \def\bigcup{\mathop{\mathchoice{\raisebox{-0.3em}{\scalebox{1.7}{$\cup$}}}{\scalebox{1.25}{$\cup$}}{\cup}{\cup}}\displaylimits}%
  \def\bigcap{\mathop{\mathchoice{\raisebox{-0.3em}{\scalebox{1.7}{$\cap$}}}{\scalebox{1.25}{$\cap$}}{\cap}{\cap}}\displaylimits}%
  \def\lesseqqgtr{\mathrel{\lessgtr}}\def\bigoplus{\mathop{\oplus}\displaylimits}%
}
\providecommand{\ssi}{si et seulement si} % abréviation fréquente
\AtBeginDocument{\providecommand{\wideparen}{\overparen}} % arc de cercle

%% FIGFIX-BEGIN v5 — correctifs globaux des figures (docs/onbuch-generation/tools/figfix)
\usepackage{adjustbox}\usepackage{etoolbox}\tcbuselibrary{raster}
% toute figure TikZ/pgfplots plus large que la ligne est réduite à la largeur disponible
\BeforeBeginEnvironment{tikzpicture}{\begin{adjustbox}{max width=\linewidth}}
\AfterEndEnvironment{tikzpicture}{\end{adjustbox}}
% légendes pgfplots lisibles par-dessus les courbes
\pgfplotsset{every axis/.append style={legend style={fill=white,fill opacity=0.92,text opacity=1,draw=black!15,inner sep=3pt}}}
% étiquettes d'arêtes (syntaxe edge["..."]) avec fond blanc
\tikzset{every edge quotes/.append style={fill=white,inner sep=1.2pt}}
%% FIGFIX-END
\begin{document}

\pagegarde{L’Etat et le pouvoir}{Philosophie}{Tle Tech}{Philosophie – classe de Terminale technique}{N19}{3 séances (fiche de progression harmonisée)}{Cette leçon analyse la nature, les fondements et les formes de l'État moderne en s'appuyant sur la réalité institutionnelle camerounaise. Elle outille l'élève pour comprendre le contrat social, distinguer les régimes politiques et évaluer l'action publique à travers des documents officiels et des situations vécues.
\vspace{4pt}
\begin{multicols}{2}\small
\begin{itemize}
\item À la fin de cette leçon, je sais définir l'État par ses quatre éléments constitutifs (population, territoire, souveraineté, gouvernement) et les illustrer par l'exemple du Cameroun.
\item Je sais expliquer les trois théories majeures du fondement de l'obligation politique (force, droit divin, contrat social) et les relier à l'histoire politique camerounaise.
\item Je sais classifier les formes de gouvernement (monarchie, république ; parlementaire, présidentiel, semi-présidentiel) et situer le régime camerounais dans cette typologie.
\item Je sais analyser les fonctions régaliennes et sociales de l'État à partir du budget général et des politiques publiques (éducation, santé, infrastructure).
\item Je sais appliquer la distinction weberienne du monopole de la violence physique légitime à des situations concrètes (maintien d'ordre, justice, défense).
\item Je sais rédiger une dissertation philosophique structurée sur la légitimité du pouvoir ou la justice sociale en mobilisant des références précises (Hobbes, Locke, Rousseau, Kant, Rawls).
\end{itemize}\end{multicols}}

\begin{sommaire}
\begin{multicols}{2}
\begin{enumerate}
\item 1. Qu'est-ce que l'État ? Définition et éléments constitutifs
\item 2. Fondements et légitimité du pouvoir politique
\item 3. Les formes de gouvernement et le régime politique camerounais
\item 4. Rôle, finalité et fonctions de l'État : L'État-providence en question
\item 5. TP Intégré : Cartographie institutionnelle et Analyse de l'action publique locale
\item 6. Dissertation type Bac : Méthodologie et Sujets probables
\item Activité d'intégration
\item Méthodes et exercices résolus
\item Exercices
\item Corrigés des exercices
\item Fiche bilan
\end{enumerate}
\end{multicols}
\end{sommaire}

\begin{objectifs}
\begin{itemize}
\item À la fin de cette leçon, je sais définir l'État par ses quatre éléments constitutifs (population, territoire, souveraineté, gouvernement) et les illustrer par l'exemple du Cameroun.
\item Je sais expliquer les trois théories majeures du fondement de l'obligation politique (force, droit divin, contrat social) et les relier à l'histoire politique camerounaise.
\item Je sais classifier les formes de gouvernement (monarchie, république ; parlementaire, présidentiel, semi-présidentiel) et situer le régime camerounais dans cette typologie.
\item Je sais analyser les fonctions régaliennes et sociales de l'État à partir du budget général et des politiques publiques (éducation, santé, infrastructure).
\item Je sais appliquer la distinction weberienne du monopole de la violence physique légitime à des situations concrètes (maintien d'ordre, justice, défense).
\item Je sais rédiger une dissertation philosophique structurée sur la légitimité du pouvoir ou la justice sociale en mobilisant des références précises (Hobbes, Locke, Rousseau, Kant, Rawls).
\item Je sais construire un organigramme institutionnel du Cameroun et un croquis de la répartition des pouvoirs (exécutif, législatif, judiciaire).
\item Je sais évaluer critique ment l'efficacité de l'État providence face aux défis du développement (corruption, décentralisation, émergence 2035).
\end{itemize}
\end{objectifs}

\begin{prerequis}
\begin{itemize}
\item Notion de société, de droit et de morale (programme de 1ère) : distinction fait / valeur, norme juridique vs norme morale.
\item Connaissance de l'histoire politique du Cameroun : indépendance, réunification, passage au multipartisme (1990), révision constitutionnelle de 1996/2008.
\item Notion de séparation des pouvoirs (Montesquieu) et de souveraineté nationale (programme d'Histoire-Géographie).
\item Maîtrise de la méthode de la dissertation philosophique (thèse, antithèse, synthèse) et de l'explication de texte.
\item Repères géographiques : carte administrative du Cameroun (10 régions, chef-lieux, capitale politique vs économique).
\end{itemize}
\end{prerequis}

\newpage

\coursec{Qu'est-ce que l'État ? Définition et éléments constitutifs}

Au carrefour Nlongkak à Yaoundé, un agent de police arrête un taxi pour un contrôle routier ; le chauffeur s'exécute sans discuter, bien qu'il soit pressé. Plus loin, des riverains bloquent la route avec des pneus brûlés pour réclamer l'électricité, forçant l'administration à négocier. Pourquoi acceptons-nous l'autorité de l'un et contestons-nous l'autre ? Qu'est-ce qui fonde le pouvoir de l'État au Cameroun et quelles sont ses limites face aux attentes des populations ?

Cette scène ordinaire de la vie camerounaise pose le problème central de cette leçon : \textbf{qu'est-ce que l'État, et de quoi est-il fait pour que son autorité s'impose à nous ?} Nous répondrons d'abord en définissant rigoureusement l'État et en le distinguant des notions voisines (Nation, société civile, gouvernement, administration), puis en analysant ses quatre éléments constitutifs, avant d'examiner sa nature juridique et ses différentes formes d'organisation.

\courssub{Distinguer l'État des notions voisines}

\textbf{L'intuition d'abord.} Dans le langage courant, on dit indifféremment « l'État a décidé », « le gouvernement a décidé », « le pays a décidé ». Cette confusion est un obstacle philosophique majeur. Or, lorsque le président change après une élection, le Cameroun ne cesse pas d'exister : les frontières restent, les fonctionnaires restent, les dettes et les traités restent. Ce qui change, c'est l'équipe dirigeante, pas la structure. Il faut donc séparer soigneusement cinq réalités.

\begin{definition}[L'État]
L'\cle{État} est une \textbf{personne morale de droit public}, dotée de la \cle{souveraineté}, s'exerçant sur une \textbf{population} et un \textbf{territoire} déterminés, pourvue d'un \textbf{gouvernement} et d'une \textbf{administration}.
\end{definition}

Précisons chaque notion voisine :

\begin{itemize}
\item La \cle{Nation} est une \emph{communauté humaine} unie par la langue, la culture, l'histoire, la religion ou la volonté de vivre ensemble (Renan la définit comme un « plébiscite de tous les jours »). C'est une réalité \textbf{culturelle et subjective}.
\item La \cle{société civile} désigne l'ensemble des organisations et des rapports sociaux qui ne relèvent pas de l'État : familles, entreprises, syndicats, associations, Églises, médias, tontines. C'est la sphère de la liberté privée et de l'initiative collective.
\item Le \cle{gouvernement} est l'organe dirigeant \emph{temporaire} qui prend les grandes décisions politiques au nom de l'État.
\item L'\cle{administration} est l'ensemble des services et des agents (fonctionnaires, préfets, policiers, enseignants) qui exécutent les décisions et assurent la continuité des services publics.
\end{itemize}

\begin{attention}[Ne pas confondre Nation et État]
Erreur fréquente : confondre la \textbf{Nation} (communauté humaine, culturelle, subjective) et l'\textbf{État} (structure juridique, politique, objective). Une nation peut exister sans État (les Kurdes), et un État peut rassembler plusieurs nations ou groupes culturels. Le Cameroun est de fait un État \textbf{pluriculturel} (plus de 250 groupes ethniques, deux traditions linguistiques officielles : français et anglais) au sein d'une même structure étatique unitaire.
\end{attention}

\begin{methode}[Distinguer État et Gouvernement en trois questions]
\begin{enumerate}
\item \textbf{Qu'est-ce qui est permanent ?} L'État : institutions, territoire, population, administration.
\item \textbf{Qu'est-ce qui est temporaire ?} Le gouvernement : il est nommé, remanié, remplacé.
\item \textbf{Qui gère quoi ?} Le gouvernement \emph{gère} l'État, il ne s'y confond pas.
\end{enumerate}
Exemple : le Gouvernement dirigé par le Premier ministre Joseph Dion Ngute gère l'État du Cameroun ; s'il est remplacé demain, l'État du Cameroun demeure.
\end{methode}

\begin{popfigure}
\begin{center}
\begin{tikzpicture}[scale=1]
\draw[fill=popblueL,draw=popblue,thick,rounded corners=8pt] (-4.6,-3.1) rectangle (4.6,3.1);
\node[popdark,font=\bfseries] at (0,2.7) {ÉTAT (structure juridique et politique)};
\draw[fill=popgreenL,draw=popgreen,thick] (-2.1,0.4) circle (1.5);
\node[popdark,font=\small\bfseries] at (-2.1,1.15) {Nation};
\node[popdark,font=\scriptsize,text width=2.4cm,align=center] at (-2.1,0.15) {communauté culturelle et historique};
\draw[fill=poporangeL,draw=poporange,thick] (2.1,0.4) circle (1.5);
\node[popdark,font=\small\bfseries] at (2.1,1.15) {Société civile};
\node[popdark,font=\scriptsize,text width=2.4cm,align=center] at (2.1,0.15) {associations, entreprises, familles};
\draw[fill=poppinkL,draw=poppink,thick] (0,-1.9) circle (1.05);
\node[popdark,font=\small\bfseries] at (0,-1.65) {Gouvernement};
\node[popdark,font=\scriptsize] at (0,-2.2) {organe temporaire};
\end{tikzpicture}
\end{center}
\legende{Diagramme d'imbrication : la Nation, la société civile et le gouvernement relèvent de l'État, mais aucun ne s'y confond. Le gouvernement est la partie la plus petite car il est temporaire, alors que l'État englobe tout.}
\end{popfigure}

\courssub{Les quatre éléments constitutifs de l'État}

\textbf{L'intuition d'abord.} Une foule dans un stade n'est pas un État. Un territoire vide non plus. Un chef sans population obéissante non plus. Il faut que \emph{quatre conditions} se rencontrent. Le droit international les a fixées dans la \textbf{Convention de Montevideo (1933)}.

\begin{propriete}[Les quatre éléments constitutifs de l'État]
Selon la Convention de Montevideo (1933), un État suppose :
\begin{enumerate}
\item une \cle{population} : une communauté humaine stable, quelle que soit sa taille ;
\item un \cle{territoire} : un espace délimité par des frontières (terrestres, maritimes, aériennes) ;
\item un \cle{gouvernement} : une autorité organisée capable de faire respecter l'ordre juridique ;
\item la \cle{souveraineté} : la suprématie à l'intérieur (aucune autorité ne lui est supérieure sur son territoire) et l'indépendance à l'extérieur (il n'obéit à aucun autre État et entre en relation d'égal à égal avec eux).
\end{enumerate}
\end{propriete}

\begin{exemplebox}[Application au Cameroun]
\begin{itemize}
\item \textbf{Population} : environ \num{28} millions d'habitants (4\textsuperscript{e} Recensement général de la population et de l'habitat, 2023), jeune et diversifiée.
\item \textbf{Territoire} : \SI{475442}{km^{2}}, des frontières terrestres avec le Nigeria, le Tchad, la Centrafrique, le Congo, le Gabon et la Guinée équatoriale, et une façade maritime sur le golfe de Guinée (environ 400 km de côtes).
\item \textbf{Gouvernement} : un Président de la République, un Premier ministre, des ministres, relayés par l'administration déconcentrée (gouverneurs, préfets, sous-préfets).
\item \textbf{Souveraineté} : le Cameroun siège à l'ONU, à l'Union africaine et à la CEMAC ; aucune puissance étrangère ne légifère sur son sol.
\end{itemize}
\end{exemplebox}

\begin{popfigure}
\begin{center}
\begin{tikzpicture}[scale=0.92]
% Contour très simplifié du Cameroun
\draw[fill=popcream,draw=popdark,thick] (0,0) -- (1.2,2.2) -- (2.6,3.4) -- (3.4,4.6) -- (4.4,5.4) -- (5.2,4.4) -- (5.0,2.6) -- (4.2,1.4) -- (4.6,0.2) -- (3.2,-0.8) -- (1.6,-0.6) -- (0.4,-1.0) -- cycle;
% Nord (flèche)
\draw[->,thick,popdark] (6.0,4.6) -- (6.0,5.6) node[above]{\scriptsize N};
% Villes
\fill[popink] (2.6,1.0) circle (2.2pt) node[right,font=\scriptsize\bfseries]{Yaoundé (capitale politique)};
\fill[poporange] (1.5,0.0) circle (2.2pt) node[left,font=\scriptsize\bfseries]{Douala (capitale économique)};
% Régions : flèches indicatives
\foreach \a/\t in {90/{Nord}, 150/{Ouest}, 210/{Sud-Ouest}, 270/{Sud}, 330/{Est}, 30/{Adamaoua}} {
  \draw[->,popblue,thin] (2.6,2.2) ++(\a:1.1) -- ++(\a:0.55);
}
\node[font=\scriptsize,popmuted] at (2.6,2.2) {10 régions};
% Symboles de souveraineté
\draw[fill=popgreen,draw=popdark] (5.6,-0.4) rectangle (6.6,0.2);
\fill[popgold] (6.1,-0.1) circle (1.6pt);
\node[font=\scriptsize,popdark] at (6.1,-0.7) {Drapeau};
\draw[popdark,thick] (5.6,-1.3) circle (0.35);
\node[font=\scriptsize,popdark] at (6.1,-1.3) {Sceau};
\end{tikzpicture}
\end{center}
\legende{Croquis schématique du Cameroun : contour simplifié, capitale politique (Yaoundé), capitale économique (Douala), les dix régions rayonnant autour du centre, et deux symboles de la souveraineté (drapeau et sceau de la République). Échelle et formes indicatives.}
\end{popfigure}

\begin{savaistu}[Le Cameroun, acteur souverain de la scène internationale]
Le Cameroun est membre de l'ONU depuis 1960, année de son indépendance, et membre fondateur de l'OUA (1963), devenue Union africaine. Sa souveraineté externe s'exerce concrètement à travers une centaine de missions diplomatiques et consulaires dans le monde : ambassades, consulats, représentations permanentes. Chaque passeport délivré, chaque traité signé, chaque vote à l'Assemblée générale des Nations unies est un acte de souveraineté.
\end{savaistu}

\courssub{L'État, personne morale de droit public}

\textbf{L'intuition d'abord.} Une entreprise comme la SABC ou la CAMTEL peut signer des contrats, posséder des immeubles, être condamnée en justice : le droit la traite comme une « personne », bien qu'elle ne soit pas un être humain. C'est une \emph{personne morale}. L'État aussi, mais d'un type supérieur : une personne morale de \textbf{droit public}, dotée de prérogatives de puissance publique (il peut exproprier, lever l'impôt, emprisonner selon la loi).

Cette personnalité juridique a trois conséquences majeures :

\begin{itemize}
\item \textbf{Un patrimoine.} L'État possède un \emph{domaine public} (routes, ports, forêts domaniales, fleuves, littoral) inaliénable, et un \emph{domaine privé} (bâtiments administratifs, terrains non affectés) qu'il peut gérer ou vendre.
\item \textbf{Une capacité juridique.} L'État peut conclure des traités internationaux (accords de la CEMAC, traités de coopération), contracter des emprunts (obligations, dette publique), ester en justice.
\item \textbf{Une responsabilité.} L'État répond des dommages causés par ses services : un élève blessé par la chute d'un toit d'école publique mal entretenu peut engager la responsabilité de l'État.
\end{itemize}

\begin{exemplebox}[La capacité financière de l'État en chiffres]
Le budget de l'État du Cameroun pour l'exercice 2024 s'élève à environ \textbf{\num{6740} milliards de FCFA} (Loi de finances n°2023/017). Sa structure approximative : près de 30 \% pour les dépenses de personnel (salaires des fonctionnaires), environ 22 \% pour l'investissement (routes, écoles, hôpitaux) et autour de 15 \% pour le service de la dette (remboursements et intérêts) ; le solde couvrant les biens, services, subventions et transferts. Ces chiffres montrent que l'État est le premier acteur économique du pays : aucune entreprise privée ne dispose d'une telle capacité d'action financière.
\end{exemplebox}

\begin{propriete}[La permanence de l'État]
\textbf{L'État est permanent, les gouvernements sont temporaires.} L'État ne meurt pas (sauf disparition par annexion ou éclatement) : il transcende les alternances politiques, les remaniements ministériels, les changements de régime. Les traités signés, les dettes contractées et les frontières établies engagent le Cameroun au-delà de toute équipe gouvernementale. C'est le principe de la \emph{continuité de l'État}.
\end{propriete}

\courssub{État unitaire, État fédéral, État-nation : les formes de l'État}

\textbf{L'intuition d'abord.} Tous les États ne distribuent pas le pouvoir de la même manière sur leur territoire. Au Nigeria, l'État de Lagos adopte ses propres lois dans de nombreux domaines ; au Cameroun, toutes les grandes lois viennent de Yaoundé. C'est la différence entre État fédéral et État unitaire.

\begin{definition}[État unitaire et État fédéral]
Un \cle{État unitaire} concentre le pouvoir législatif et politique au centre ; les collectivités locales n'ont que des compétences déléguées (exemple : le Cameroun, la France). Un \cle{État fédéral} partage la souveraineté entre l'État fédéral et des États fédérés dotés de leurs propres constitutions et parlements (exemples : le Nigeria, les États-Unis, l'Allemagne).
\end{definition}

\begin{attention}[Décentralisation n'est pas fédéralisme]
La \cle{décentralisation} camerounaise (régions, communes, conseils régionaux installés en 2020) est un \textbf{transfert de compétences}, non un transfert de souveraineté. Les régions gèrent des affaires locales (développement économique, équipements) mais ne peuvent ni légiférer souverainement, ni se séparer, ni signer des traités. Le Cameroun reste, selon sa Constitution, un État unitaire décentralisé.
\end{attention}

Le concept d'\cle{État-nation} désigne l'idéal d'une adéquation parfaite entre les frontières de l'État et l'identité nationale : un seul peuple, un seul État. En Afrique, les frontières héritées de la colonisation ont souvent découpé les groupes humains sans tenir compte de cette adéquation. Le Cameroun illustre cette difficulté : la question des minorités, et singulièrement la \emph{question anglophone} (régions du Nord-Ouest et du Sud-Ouest, héritées de l'administration britannique), montre que l'unité étatique ne supprime pas automatiquement la diversité des identités. Le défi philosophique est de construire une nation politique \emph{par} l'État, sans nier les cultures qui la composent.

\begin{exoresolu}[État failli ou État fort ? Appliquer les critères]
Énoncé : Le \emph{Fragile States Index} classe les États selon douze indicateurs (cohésion sociale, économie, légitimité politique, services publics, sécurité...). Un État « failli » est un État qui ne contrôle plus son territoire ni ne fournit les services essentiels (exemple : la Somalie dans les années 1990). Le Cameroun se situe dans la catégorie intermédiaire des États « en alerte ». En vous appuyant sur les quatre éléments constitutifs, montrez que le Cameroun, malgré ses fragilités, demeure pleinement un État.
\tcblower
\textbf{Solution détaillée.} Vérifions chaque élément de Montevideo. (1) \emph{Population} : environ \num{28} millions d'habitants recensés en 2023, rattachés à l'État par la nationalité camerounaise. (2) \emph{Territoire} : \SI{475442}{km^{2}} aux frontières internationalement reconnues, y compris après l'arrêt de la Cour internationale de justice (2002) attribuant la presqu'île de Bakassi au Cameroun. (3) \emph{Gouvernement} : des institutions fonctionnelles (exécutif, parlement, justice, administration déconcentrée jusqu'au sous-préfet). (4) \emph{Souveraineté} : siège à l'ONU, à l'UA, à la CEMAC ; monnaie et diplomatie propres en coopération régionale. Certes, des indicateurs révèlent des fragilités (conflits dans certaines régions, inégalités de services publics), mais une fragilité n'est pas une faillite : l'État camerounais continue d'exercer son autorité et de fournir des services sur l'essentiel de son territoire. Conclusion : le Cameroun est un État \emph{fragile} au sens des indicateurs, mais un État \emph{pleinement constitué} au sens du droit international.
\end{exoresolu}

\begin{popfigure}
\begin{center}
\begin{tikzpicture}[
  node distance=0.55cm,
  box/.style={draw=popblue,fill=popblueL,rounded corners=4pt,thick,align=center,font=\scriptsize,text width=3.4cm,minimum height=0.9cm},
  big/.style={draw=popink,fill=popink!15,rounded corners=4pt,thick,align=center,font=\scriptsize\bfseries,text width=5.2cm,minimum height=1cm}]
\node[big, align=center] (pr){Président de la République\\(Chef de l'État, chef de l'exécutif)};
\node[box,below left=0.8cm and -0.4cm of pr, align=center] (pm){Premier ministre\\Chef du gouvernement};
\node[box,below=0.5cm of pm] (min) {Ministères et administration};
\node[box,right=2.6cm of pm, align=center] (parl){Parlement\\Assemblée nationale + Sénat};
\node[box,below=0.5cm of parl, align=center] (jud){Justice et contrôle de constitutionnalité\\Cour suprême, Conseil constitutionnel};
\draw[->,thick,popdark] (pr) -- (pm);
\draw[->,thick,popdark] (pm) -- (min);
\draw[<->,thick,popdark,dashed] (pr.east) -- (parl.west) node[midway,above,font=\tiny]{collaboration / contrôle};
\draw[<->,thick,popdark,dashed] (parl) -- (jud);
\end{tikzpicture}
\end{center}
\legende{Architecture institutionnelle simplifiée du Cameroun : l'exécutif (Président, Premier ministre, ministères), le législatif (Parlement bicaméral) et le judiciaire/constitutionnel. Les flèches pleines indiquent la subordination hiérarchique, les pointillés la collaboration et le contrôle entre pouvoirs.}
\end{popfigure}

\begin{aretenir}[L'essentiel de la section]
\begin{itemize}
\item L'État est une personne morale de droit public, souveraine, reposant sur quatre éléments : population, territoire, gouvernement, souveraineté (Convention de Montevideo, 1933).
\item Il ne se confond ni avec la Nation (communauté culturelle), ni avec la société civile (sphère privée et associative), ni avec le gouvernement (organe temporaire).
\item L'État est permanent ; les gouvernements passent. Il possède un patrimoine, une capacité juridique et une responsabilité.
\item Le Cameroun est un État unitaire décentralisé, pluriculturel de fait, membre souverain des organisations internationales.
\end{itemize}
\end{aretenir}

Ainsi, l'État n'est ni une abstraction lointaine ni un simple synonyme de « gouvernement » : c'est la structure juridique permanente qui organise notre vie commune. Mais dire ce qu'est l'État ne suffit pas : il reste à comprendre \emph{pourquoi} son pouvoir est accepté comme légitime, et selon quelles formes de gouvernement il peut s'exercer — ce sera l'objet des sections suivantes.

\coursec{Fondements et légitimité du pouvoir politique}

\courssub{La question fondamentale : pourquoi obéir ? La distinction fait / valeur}

Après avoir identifié \emph{ce qu'est} l'État (population, territoire, gouvernement, souveraineté), nous devons maintenant interroger \emph{ce qui fonde son autorité}. Pourquoi les citoyens acceptent-ils de se soumettre aux décisions d'un pouvoir qui, in fine, détient le monopole de la violence physique légitime (Weber) ? La réponse ne va pas de soi.

\courssub{Le piège naturaliste : la force ne fait pas le droit}

L'observation immédiate montre que le pouvoir s'exerce souvent par la contrainte : police, armée, prisons. Le sophiste \textbf{Thrasymaque} (dans \emph{La République} de Platon) ose la thèse cynique : « La justice n'est rien d'autre que l'intérêt du plus fort. » Pour lui, le droit n'est que le masque de la force. \textbf{Machiavel} (\emph{Le Prince}, 1532) prolonge cette analyse réaliste : le prince doit savoir être « lion et renard », user de la force et de la ruse pour conserver le pouvoir. La légitimité se confond alors avec l'efficacité.

\begin{attention}[Erreur fréquente : Confondre fait et droit]
Ne pas confondre \textbf{puissance} (capacité physique de contraindre, fait brut) et \textbf{autorité} (droit d'être obéi, valeur morale). Un bandit qui vous vole sous la menace a la puissance, mais pas l'autorité. Obéir par peur n'est pas obéir par devoir. Comme l'écrit \textbf{Pascal} (\emph{Pensées}, Lafuma 298) : \og Justice sans force est impuissante ; force sans justice est tyrannique \fg. La force assure l'ordre, la justice assure la légitimité. Un pouvoir seulement fort est une tyrannie ; un pouvoir seulement juste sans force est une utopie sans prise sur le réel.
\end{attention}

\courssub{Le saut normatif : du « est » au « doit »}

La philosophie politique moderne (Kant, Kelsen, Hart) insiste sur la distinction \textbf{fait / valeur} (ou \emph{Sein / Sollen}). Le fait que l'État \emph{exerce} le pouvoir ne prouve pas qu'il \emph{doit} l'exercer. La légitimité est cette propriété normative qui transforme la contrainte en obligation morale. Elle répond à la question : \og Qu'est-ce qui rend l'obéissance non pas seulement prudente, mais \textbf{juste} ? \fg

\begin{definition}[Légitimité (Max Weber)]
Croyance en la validité de l'ordre établi, fondant le \textbf{devoir d'obéissance} des gouvernés. Elle n'est pas une qualité objective du pouvoir, mais une \textbf{représentation subjective} partagée par les gouvernants et les gouvernés. Weber distingue trois \textbf{types purs} de légitimité (voir § 2.3).
\end{definition}

\courssub{Les grandes théories classiques de la légitimité}

\courssub{La légitimité traditionnelle : le droit divin et la chefferie coutumière}

Historiquement, la première forme de légitimité repose sur la \textbf{tradition} : « C'est ainsi depuis toujours. » Le pouvoir est sacré, transmis par hérédité ou désignation coutumière. En Europe, c'est le \textbf{droit divin} : le roi tient son pouvoir de Dieu (Bossuet : \og L'État, c'est moi \fg, attribution à Louis XIV). Au Cameroun, cette forme persiste vivacement à travers les \textbf{chefferies traditionnelles} (Fon dans le Nord-Ouest, Lamido dans le Nord, Sultan au Sud, Chef de groupement/village partout).

\begin{savaistu}[Chefferie traditionnelle et Constitution camerounaise]
L'article 37 de la Constitution de 1996 (révisée 2008) dispose : \og Les chefferies traditionnelles sont reconnues. La loi fixe leur statut. \fg L'État moderne délègue ainsi une part de son monopole de la violence (police administrative, justice coutumière) aux chefs traditionnels, créant un \textbf{pluralisme juridique} où la légitimité traditionnelle côtoie la légalité rationnelle. Le Lamido de Rey Bouba ou le Fon de Bafut incarnent cette continuité sacrée.
\end{savaistu}

\courssub{La légitimité charismatique : le leader providentiel}

Weber définit le charisme comme une \og qualité extraordinaire \fg conférant au leader une mission personnelle (prophète, héros de guerre, père de la nation). L'obéissance est due à la \textbf{personne}, non à la fonction. Au Cameroun, \textbf{Ahmadou Ahidjo} (premier président, 1960-1982) a bénéficié d'une légitimité charismatique forte : artisan de l'indépendance, de la réunification (1961) et de la stabilité. Les débuts de \textbf{Paul Biya} (1982) ont aussi reposé sur un charisme de « rénovation » et de rigueur morale. Mais le charisme est \textbf{instable} : il doit se \emph{routiniser} (se transformer en institution) pour durer.

\courssub{La théorie du contrat social : le fondement moderne de la légalité-rationnelle}

C'est la rupture majeure des Lumières (XVIIe-XVIIIe siècle). Le pouvoir ne vient ni de Dieu, ni de la force, ni de la tradition, mais d'un \textbf{accord libre et rationnel} entre les hommes pour sortir de l'état de nature. Trois versions majeures structurent notre modernité politique.

\begin{popfigure}
\centering
\begin{tikzpicture}[
    node distance=1.2cm and 2.5cm,
    every node/.style={font=\small},
    theory/.style={rectangle, draw=popblue, thick, fill=popblueL, rounded corners, minimum width=3.5cm, text centered, text width=3.2cm},
    state/.style={rectangle, draw=popdark, fill=popcream, rounded corners, minimum width=3cm, text centered, text width=2.8cm},
    arrow/.style={-Stealth, thick, popdark},
    labelarrow/.style={midway, fill=white, inner sep=2pt, font=\scriptsize, text=poppurple}
]
% Colonnes : Hobbes, Locke, Rousseau
% Ligne 1 : État de nature
\node[state, align=center] (hobbes_nat){ \textbf{Hobbes}\\\og Guerre de tous contre tous \fg\\Insécurité permanente\\Vie \og solitaire, pauvre, méchante, brutale et courte \fg };
\node[state, right=of hobbes_nat, align=center] (locke_nat){ \textbf{Locke}\\\og État de paix, de liberté, de propriété \fg\\Loi de nature (raison)\\Mais pas de juge impartial $\rightarrow$ insécurité des biens };
\node[state, right=of locke_nat, align=center] (rousseau_nat){ \textbf{Rousseau}\\\og Homme naturel \fg : bon, libre, égal\\Pitié, amour de soi\\Corruption par la société (propriété, inégalités) };

% Ligne 2 : Contrat
\node[theory, below=of hobbes_nat, align=center] (hobbes_cont){ \textbf{Pacte de soumission}\\\textit{Léviathan} (1651)\\\textit{Aliénation totale} des droits au Souverain (sauf droit à la vie) };
\node[theory, below=of locke_nat, align=center] (locke_cont){ \textbf{Pacte de confiance}\\\textit{Traité du gouvernement civil} (1690)\\\textit{Délégation} du pouvoir législatif/exécutif\\Conservation des droits naturels (Vie, Liberté, Propriété) };
\node[theory, below=of rousseau_nat, align=center] (rousseau_cont){ \textbf{Pacte d'association}\\\textit{Du contrat social} (1762)\\\textit{Aliénation totale à la communauté} $\rightarrow$ Souveraineté du Peuple\\Volonté Générale $\neq$ Somme des volontés particulières };

% Ligne 3 : Type de pouvoir
\node[state, below=of hobbes_cont, align=center] (hobbes_pow){ \textbf{Pouvoir absolu}\\\textit{Monarchie} ou \textit{Assemblée}\\\textit{Indivisible, illimité, irreprésentable}\\\og Le roi ne peut pas faire injustice \fg };
\node[state, below=of locke_cont, align=center] (locke_pow){ \textbf{Pouvoir limité / Séparé}\\\textit{Gouvernement de confiance}\\\textit{Séparation des pouvoirs} (Législatif > Exécutif)\\\textit{Droit de résistance} à la tyrannie };
\node[state, below=of rousseau_cont, align=center] (rousseau_pow){ \textbf{Pouvoir démocratique}\\\textit{Souveraineté populaire inaliénable}\\\textit{Gouvernement = simple mandataire}\\\textit{Lois = expression Volonté Générale} };

% Ligne 4 : Résistance
\node[state, below=of hobbes_pow, fill=poporangeL, draw=poporange, align=center] (hobbes_res){ \textbf{Pas de droit de résistance}\\\textit{Sauf si le Souverain ne protège plus la vie}\\\og L'obéissance est le prix de la paix \fg };
\node[state, below=of locke_pow, fill=popgreenL, draw=popgreen, align=center] (locke_res){ \textbf{Droit de résistance / Insurrection}\\\textit{Si gouvernement viole la confiance (biens, liberté)}\\\textit{Retour à l'état de nature pour juger} };
\node[state, below=of rousseau_pow, fill=poppurpleL, draw=poppurple, align=center] (rousseau_res){ \textbf{Résistance = fidélité au contrat}\\\textit{Contre le gouvernement (mandataire), pas contre le Souverain (Peuple)}\\\textit{Refus d'obéir à une loi non générale} };

% Flèches
\draw[arrow] (hobbes_nat) -- (hobbes_cont) node[labelarrow] {Contrat};
\draw[arrow] (hobbes_cont) -- (hobbes_pow) node[labelarrow] {Instaure};
\draw[arrow] (hobbes_pow) -- (hobbes_res) node[labelarrow] {Conséquence};

\draw[arrow] (locke_nat) -- (locke_cont) node[labelarrow] {Contrat};
\draw[arrow] (locke_cont) -- (locke_pow) node[labelarrow] {Instaure};
\draw[arrow] (locke_pow) -- (locke_res) node[labelarrow] {Garantit};

\draw[arrow] (rousseau_nat) -- (rousseau_cont) node[labelarrow] {Contrat};
\draw[arrow] (rousseau_cont) -- (rousseau_pow) node[labelarrow] {Instaure};
\draw[arrow] (rousseau_pow) -- (rousseau_res) node[labelarrow] {Préserve};

% Étiquettes colonnes
\node[above=0.2cm of hobbes_nat, font=\bfseries\large, text=popblue] {HOBBES};
\node[above=0.2cm of locke_nat, font=\bfseries\large, text=popgreen] {LOCKE};
\node[above=0.2cm of rousseau_nat, font=\bfseries\large, text=poppurple] {ROUSSEAU};

\node[left=0.3cm of hobbes_nat, rotate=90, anchor=south, font=\bfseries\large, text=popdark, xshift=-1.5cm] {\textbf{État de nature}};
\node[left=0.3cm of hobbes_cont, rotate=90, anchor=south, font=\bfseries\large, text=popdark, xshift=-1.5cm] {\textbf{Contrat social}};
\node[left=0.3cm of hobbes_pow, rotate=90, anchor=south, font=\bfseries\large, text=popdark, xshift=-1.5cm] {\textbf{Type de pouvoir}};
\node[left=0.3cm of hobbes_res, rotate=90, anchor=south, font=\bfseries\large, text=popdark, xshift=-1.5cm] {\textbf{Droit de résistance}};
\end{tikzpicture}
\legende{Arbre comparatif des trois théories classiques du contrat social. Hobbes fonde l'absolutisme par la peur ; Locke fonde le libéralisme constitutionnel par la protection des biens ; Rousseau fonde la démocratie radicale par la souveraineté populaire.}
\end{popfigure}

\begin{exemplebox}[Application camerounaise : La Constitution de 1996 comme contrat social]
Le Préambule de la Constitution camerounaise proclame : \og Nous, Peuple du Cameroun... proclamons notre attachement aux libertés fondamentales... et adoptons la présente Constitution. \fg C'est un acte \textbf{lockeen} (protection des droits : vie, liberté, propriété, Art. 1er et 65) et \textbf{rousseauiste} (souveraineté populaire, Art. 2 : \og Le peuple exerce sa souveraineté par ses représentants et par voie de référendum \fg). Mais la concentration des pouvoirs exécutifs (Président élu au suffrage universel direct, nommant le PM, les ministres, les hauts fonctionnaires, les juges du Conseil Constitutionnel) rappelle une tentation \textbf{hobbésienne} de l'efficacité au prix de la séparation des pouvoirs.
\end{exemplebox}

\courssub{La typologie weberienne appliquée au Cameroun}

Max Weber (\emph{Économie et Société}, 1922) offre la grille de lecture la plus opératoire pour analyser la légitimité réelle des régimes. Aucun régime pur n'existe ; ils se combinent.

\begin{popfigure}
\centering
\begin{tikzpicture}[
    node distance=1.5cm and 1cm,
    every node/.style={font=\small},
    type/.style={rectangle, draw, thick, rounded corners, minimum width=3.8cm, minimum height=2.2cm, text centered, text width=3.5cm, fill=popcream},
    trad/.style={type, draw=poporange, fill=poporangeL},
    char/.style={type, draw=poppink, fill=poppinkL},
    legal/.style={type, draw=popblue, fill=popblueL},
    ex/.style={rectangle, draw=popgreen, thick, fill=popgreenL, rounded corners, minimum width=3.5cm, text centered, text width=3.2cm, font=\scriptsize},
    arrow/.style={-Stealth, thick, popdark},
]
% Trois types purs
\node[trad, align=center] (T){ \textbf{Légitimité TRADITIONNELLE}\\\og C'est la coutume \fg\\Base : Habitude, sacralité de l'ancienneté\\Obéissance au \emph{chef personnel} (père, roi, chef coutumier) };
\node[char, right=of T, align=center] (C){ \textbf{Légitimité CHARISMATIQUE}\\\og C'est la grâce / le héros \fg\\Base : Dévouement personnel au leader\\Obéissance au \emph{chef missionné} (prophète, révolutionnaire, père de la nation) };
\node[legal, right=of C, align=center] (L){ \textbf{Légitimité LÉGALE-RATIONNELLE}\\\og C'est la loi / la procédure \fg\\Base : Croyance en la légalité des règles\\Obéissance à la \emph{fonction}, pas à la personne\\Règles impersonnelles, prévisibles, bureaucratiques };

% Exemples Cameroun
\node[ex, below=of T, align=center] (Tcam){ \textbf{Exemples CM} :\\• Lamido, Fon, Sultan (Art. 37 Const.)\\• Justice coutumière (Alkali, Cour coutumière)\\• Transmission héréditaire / désignation par notables };
\node[ex, below=of C, align=center] (Ccam){ \textbf{Exemples CM} :\\• Ahidjo (1960-82) : Père de la Nation, Réunification\\• Biya (début 1982) : \og Rénovation \fg, Rigueur morale\\• Figures de l'opposition historique (ex: Fru Ndi) };
\node[ex, below=of L, align=center] (Lcam){ \textbf{Exemples CM} :\\• Élections (Président 7 ans, Députés/Sénateurs/Conseillers 5 ans)\\• ELECAM (organe de gestion), Conseil Constitutionnel\\• Constitution 1996, Codes électoral, pénal, travail\\• Administration territoriale (Gouverneurs, Préfets, Sous-préfets) };

% Flèches d'évolution (routinisation)
\draw[arrow, dashed, bend left=20] (C) to node[above, font=\scriptsize, text=poppurple] {Routinisation charismatique} (L);
\draw[arrow, dashed, bend right=20] (T) to node[below, font=\scriptsize, text=poporange] {Institutionnalisation coutumière} (L);

% Titre
\node[above=0.3cm of C, font=\bfseries\large, text=popdark] {Typologie weberienne de la légitimité (avec ancrage camerounais)};
\end{tikzpicture}
\legende{Les trois types purs de légitimité selon Max Weber et leur incarnation dans le paysage politique camerounais. La flèche pointillée indique le processus historique de \emph{routinisation} : le charisme et la tradition tendent à se bureaucratiser pour durer.}
\end{popfigure}

\courssub{La légitimité démocratique au Cameroun : le pari de la légalité-rationnelle}

Depuis la conférence nationale (1991) et la Constitution de 1996, le Cameroun revendique explicitement le modèle de la \textbf{légitimité légale-rationnelle} : le pouvoir émane du peuple par des élections régulières, libres et transparentes.

\begin{propriete}[Souveraineté populaire et suffrage universel (Art. 2 Const. Cameroun)]
\og La souveraineté nationale appartient au peuple qui l'exerce soit par ses représentants, soit par voie de référendum. \fg Le suffrage est \textbf{universel} (tout citoyen 20 ans), \textbf{égal} (une personne = une voix), \textbf{secret} et \textbf{libre}. Mandats : Président 7 ans (renouvelable indéfiniment depuis la révision de 2008), Députés 5 ans, Sénateurs 5 ans, Conseillers municipaux 5 ans.
\end{propriete}

\begin{popfigure}
\centering
\begin{tikzpicture}
\begin{axis}[
    width=\linewidth,
    height=9cm,
    ybar,
    bar width=12pt,
    ymin=0, ymax=90,
    ylabel={Taux de participation (\%)},
    xlabel={Élections présidentielles},
    symbolic x coords={1992, 1997, 2004, 2011, 2018},
    xtick=data,
    xticklabel style={rotate=0, font=\small},
    ytick={0,10,20,30,40,50,60,70,80},
    ymajorgrids=true,
    grid style={popline},
    legend style={at={(0.5,-0.25)}, anchor=north, legend columns=-1, font=\small, cells={anchor=west}},
    nodes near coords={\pgfmathprintnumber[fixed,precision=2]{\pgfplotspointmeta}\%},
    nodes near coords align={vertical},
    every node near coord/.append style={font=\tiny, text=popdark},
    popblue,
    fill=popblueL,
]
\addplot coordinates {
    (1992, 68.93)
    (1997, 82.45)
    (2004, 78.12)
    (2011, 68.28)
    (2018, 53.85)
};
\addlegendentry{Participation officielle (ELECAM / Conseil Constitutionnel)}
\end{axis}
\end{tikzpicture}
\legende{Évolution du taux de participation aux élections présidentielles au Cameroun (Sources : Rapports ELECAM, Décisions Conseil Constitutionnel). La chute brutale de 2018 (53,85\%) signale une érosion de la légitimité légale-rationnelle. L'opposition (MRC) a revendiqué un taux réel \textless 30\% dans certaines zones.}
\end{popfigure}

\begin{exemplebox}[Chiffres clés : Présidentielle 2018 — Indicateur de crise]
\begin{itemize}
    \item \textbf{Inscrits :} 6 667 754 (fichier biométrique ELECAM).
    \item \textbf{Votants :} 3 591 097.
    \item \textbf{Participation officielle :} 53,85 \%.
    \item \textbf{Résultat :} Paul Biya réélu 71,28 \% ; Maurice Kamto 14,23 \% ; Cabral Libii 6,26 \%.
    \item \textbf{Revendication opposition (MRC) :} Taux réel \textless 30 \% (boycott actif, zones anglophones quasi désertées).
    \item \textbf{Analyse :} L'écart entre chiffre officiel et perception réelle (ou chiffres contestés) = \textbf{déficit de légitimité}. Le Conseil Constitutionnel a validé les résultats, assurant la \emph{légalité}, mais la \emph{légitimité} (acceptation morale) est fracturée.
\end{itemize}
\end{exemplebox}

\courssub{Crise de légitimité : diagnostic et régulateurs}

\courssub{Les symptômes au Cameroun}

\begin{itemize}
    \item \textbf{Abstention structurelle :} Désintérêt, défiance, sentiment d'inutilité du vote (\og Ça ne change rien \fg).
    \item \textbf{Contestation post-électorale récurrente :} 1992 (villes mortes), 2018 (revendication de victoire Kamto, marche pacifique réprimée), perspectives 2025.
    \item \textbf{Déficit de confiance institutionnel :} Corruption (Cameroun souvent mal classé par Transparency International), népotisme (gestion ethnorégionale perçue des postes), justice aux ordres.
    \item \textbf{Crise anglophone (depuis 2016) :} Crise majeure du \textbf{contrat social}. Revendication initiale : common law, système éducatif, décentralisation effective. Dérive : revendications sécessionnistes (Ambazonie), violence armée, répression militaire. Le \og vivre-ensemble \fg (slogan officiel) est rompu pour une partie de la population qui ne se sent plus représentée par l'État de Yaoundé.
\end{itemize}

\begin{methode}[Analyser une crise de légitimité (Grille weberienne appliquée)]
\textbf{Étape 1 : Identifier le type de légitimité \emph{revendiqué} par le pouvoir.} (Ex: Légale-rationnelle : \og Nous avons gagné les élections selon la Constitution \fg).
\textbf{Étape 2 : Identifier le type de légitimité \emph{attendu/perçu} par les gouvernés.} (Ex: Légitimité de performance / justice sociale / proximité / respect identitaire).
\textbf{Étape 3 : Mesurer l'écart (le déficit).} (Ex: Urnes pleines mais rues vides ; lois votées mais non appliquées ; institutions existantes mais perçues comme partisanes).
\textbf{Étape 4 : Proposer des régulateurs institutionnels.}
\begin{enumerate}
    \item \textbf{Dialogue inclusif} (Grand Dialogue National 2019 : nécessaire mais insuffisant si non suivi d'effets).
    \item \textbf{Réforme électorale} (Indépendance réelle d'ELECAM, transparence du vote, biométrie fiable).
    \item \textbf{Décentralisation effective} (Transfert de compétences/ressources aux Régions/CTD pour légitimité de proximité).
    \item \textbf{Redevabilité (Accountability)} : Contrôle parlementaire réel, Cour des Comptes indépendante, protection lanceurs d'alerte.
\end{enumerate}
\end{methode}

\courssub{Exercice résolu : Qualifier la nature d'un pouvoir}

\begin{exoresolu}[Analyse de la légitimité d'un pouvoir local]
\textbf{Énoncé :} Dans une commune rurale de l'Extrême-Nord, le Maire (élu RDPC) gère les affaires courantes. Le Lamido (chef traditionnel héréditaire) règle les conflits fonciers et matrimoniaux selon la coutume musulmane. Un leader jeune, charismatique, mobilise les jeunes contre la corruption du Maire et l'archaïsme du Lamido, proposant une \og nouvelle gouvernance \fg. Analysez la configuration de légitimité.
\\
\tcblower
\textbf{Solution détaillée :}
\begin{enumerate}
    \item \textbf{Maire :} Légitimité \textbf{légale-rationnelle} (élection au suffrage universel, mandat 5 ans, Code des collectivités territoriales). \emph{Faiblesse :} Perception de corruption, clientélisme $\rightarrow$ érosion de la légalité-rationnelle vers une légitimité de fait (contrôle de l'administration/budget).
    \item \textbf{Lamido :} Légitimité \textbf{traditionnelle} pure (Art. 37 Const., coutume, hérédité, sacralité). \emph{Force :} Acceptation profonde par la population musulmane pour le droit personnel/foncier. \emph{Faiblesse :} Inadaptation aux enjeux modernes (développement, jeunesse, genre).
    \item \textbf{Leader jeune :} Légitimité \textbf{charismatique} naissante (dévouement personnel, rupture, modernité). \emph{Risque :} Instable, dépend de la personne. Pour durer, il doit se \emph{routiniser} : créer un parti, se présenter aux élections (légalité-rationnelle) ou être reconnu par le Lamido (traditionnelle).
    \item \textbf{Synthèse :} \textbf{Pluralisme légitime concurrentiel.} L'État (Maire) ne détient pas le monopole de la légitimité. La gouvernance effective nécessite une \textbf{articulation} : Maire = gestion administrative/développement ; Lamido = cohésion sociale/coutume ; Jeunes = innovation/contrôle citoyen. Une réforme de la décentralisation (transferts de compétences/clarté) est le régulateur institutionnel.
\end{enumerate}
\end{exoresolu}

\courssub{Exercices d'entraînement (Baccalauréat technique)}

\begin{exercice}[Sujet de dissertation type]
\textbf{Sujet :} \og La légalité suffit-elle à fonder la légitimité du pouvoir politique ? \fg
\textbf{Pistes :} 1. Distinction conceptuelle (légalité = conformité à la loi ; légitimité = acceptation morale). 2. Légalité nécessaire mais insuffisante (lois injustes, élections truquées validées par cour complaisante, pouvoir légal mais tyrannique). 3. Conditions de la légitimité : consentement (contrat social), justice (droits fondamentaux), performance (bien commun), responsabilité (reddition de comptes). 4. Exemples : Cameroun (légalité formelle vs contestation réelle), Afrique du Sud post-apartheid (légalité de l'apartheid vs illégitimité morale), Résistance française 1940 (illégalité formelle vs légitimité morale).
\end{exercice}

\begin{exercice}[Commentaire de texte guidé]
\textbf{Extrait :} \og Le plus fort n'est jamais assez fort pour être toujours le maître, s'il ne transforme sa force en droit, et l'obéissance en devoir. \fg — \textbf{Rousseau}, \emph{Du Contrat social}, Livre I, Chap. 3.
\textbf{Questions :}
\begin{enumerate}
    \item Expliquez la thèse de l'auteur sur le rapport entre force et droit.
    \item Pourquoi la force brute ne suffit-elle pas à fonder l'autorité politique ?
    \item Montrez comment le contrat social opère la \og transformation \fg dont parle Rousseau. Illustrer par l'exemple camerounais (Constitution, élections).
\end{enumerate}
\end{exercice}

\begin{corrige}[Exercice Commentaire de texte]
\textbf{1. Thèse :} Rousseau réfute Thrasymaque/Machiavel. La force est un fait physique ; le droit est une construction morale. Si le pouvoir repose seulement sur la force, il disparaît dès qu'une force supérieure apparaît. Pour durer, il doit transformer la contrainte en obligation (\emph{devoir}).
\textbf{2. Argument :} La force produit la soumission (peur), non l'obéissance (libre adhésion à une norme). L'esclave obéit au maître par force, pas par devoir. Si le maître dort, l'esclave fuit. Le citoyen obéit à la loi qu'il s'est prescrite (autonomie).
\textbf{3. Application CM :} La Constitution de 1996 est l'acte de \og transformation \fg : le peuple (souverain) se donne une loi fondamentale. Les élections (ELECAM) sont la procédure légale-rationnelle de renouvellement du consentement. Mais si le processus électoral est biaisé (fraude, exclusion), la \og transformation \fg est inachevée : on retombe dans la force (maintien au pouvoir par l'armée/administration), perdant le \og devoir \fg d'obéissance.
\end{corrige}

\begin{aretenir}[Synthèse : Les 3 piliers de la légitimité moderne au Cameroun]
\begin{enumerate}
    \item \textbf{Légalité constitutionnelle :} Respect scrupuleux des procédures (mandats, séparation des pouvoirs, ELECAM indépendant, Conseil Constitutionnel impartial).
    \item \textbf{Justice sociale / Performance :} L'État-providence (éducation, santé, routes, emploi jeunes) tient sa part du contrat social (Locke/Rousseau). L'échec ici détruit la légitimité plus vite qu'une élection imparfaite.
    \item \textbf{Reconnaissance identitaire / Proximité :} Intégration effective des chefferies traditionnelles (Art. 37) et décentralisation réelle (Régions/CTD) pour que le citoyen anglophone, bamiléké, peul ou bassa se sente \emph{représenté} et non seulement \emph{administré}.
\end{enumerate}
\end{aretenir}

\coursec{Les formes de gouvernement et le régime politique camerounais}

Après avoir examiné les fondements qui légitiment le pouvoir politique (section 2), il convient d'analyser son \textbf{exercice concret}. Un pouvoir légitime doit s'organiser institutionnellement : comment le chef de l'État accède-t-il à sa fonction ? Comment le gouvernement est-il formé et contrôlé ? Quel est le rapport entre l'exécutif et le législatif ? C'est l'objet de la distinction des \cle{formes de gouvernement} (ou régimes politiques). Nous appliquerons ces grilles d'analyse classiques au cas camerounais, tel que défini par la Constitution de 1996 (révisée en 2008).

\courssub{Critères de classification des régimes politiques (Duverger, Carcassonne)}

La science politique ne classe pas les régimes selon leur idéologie (démocratie/dictature) mais selon leur \textbf{architecture institutionnelle}. Maurice Duverger et Guy Carcassonne ont dégagé trois critères structurels majeurs pour qualifier un régime :

\begin{enumerate}
    \item \textbf{Le mode de désignation du Chef de l'État :} Élection (suffrage universel direct ou indirect, parlement) vs Hérédité (monarchie). C'est le critère qui distingue la \cle{République} de la \cle{Monarchie}. Le Cameroun est une République (Art. 1er Const.) : le Président est élu au suffrage universel direct.
    \item \textbf{La responsabilité du Gouvernement devant le Parlement :} Le gouvernement doit-il la confiance de l'assemblée élue pour gouverner ? Si oui (motion de censure, vote de confiance, engagement de responsabilité), le régime a une composante \cle{parlementaire}. Si non, il est \cle{présidentiel}.
    \item \textbf{La dissolution du Parlement :} Le Chef de l'État peut-il mettre fin prématurément au mandat de l'assemblée ? Ce pouvoir est l'arme de l'exécutif face au législatif dans les régimes parlementaires et semi-présidentiels.
\end{enumerate}

\begin{methode}[Comment classifier un régime politique ?]
Appliquez systématiquement le questionnaire à 4 questions (Grille Duverger/Carcassonne) :
\begin{enumerate}
    \item \textbf{Origine du Chef de l'État :} Élu (République) ou Héritier (Monarchie) ?
    \item \textbf{Origine du Chef du Gouvernement :} Nommé par le Chef de l'État seul ? Issu de la majorité parlementaire ? Élu directement ?
    \item \textbf{Responsabilité politique :} Le Gouvernement rend-il des comptes au Parlement (motion de censure, vote budget) ? Le Parlement peut-il le renverser ?
    \item \textbf{Dissolution :} Le Chef de l'État peut-il dissoudre l'Assemblée basse ?
\end{enumerate}
\emph{Application au Cameroun :} 1. Président élu au suffrage universel direct (République). 2. PM nommé par le Président (Art. 10), non nécessairement issu de la majorité. 3. Gouvernement responsable devant l'Assemblée Nationale (Art. 13), mais motion de censure très encadrée. 4. Président peut dissoudre l'Assemblée (Art. 14). \textbf{Conclusion : Régime semi-présidentiel (mixte), fortement présidentialisé.}
\end{methode}

\courssub{Les trois archétypes classiques et le cas camerounais}

\begin{definition}[Régime semi-présidentiel (Maurice Duverger)]
Régime où le \cle{Président de la République} est élu au suffrage universel et dispose de pouvoirs réels propres (exécutif, dissolution, référendum), coexistant avec un \cle{Premier Ministre} et un \cle{Gouvernement} qui sont \textbf{responsables devant le Parlement} (Assemblée Nationale). Il y a une \textbf{dualité de l'exécutif} (Diarchie).
\end{definition}

\begin{center}
\begin{popfigure}
\begin{tikzpicture}[
    node distance=2mm and 4mm,
    header/.style={fill=popblue, text=white, font=\bfseries\small, minimum height=8mm, align=center},
    camcol/.style={fill=popgold!30, font=\small, minimum height=7mm, align=center},
    stdcol/.style={fill=popcream, font=\small, minimum height=7mm, align=center},
    rowtitle/.style={fill=popblueL, font=\bfseries\small, minimum height=7mm, align=center, text width=3.2cm},
    cell/.style={font=\small, minimum height=7mm, align=center},
    thickline/.style={line width=1.5pt, draw=popdark}
]
% Largeurs colonnes : Critère (3.2cm), Présidentiel (3.2cm), Parlementaire (3.2cm), Semi-présidentiel (3.2cm), Cameroun (3.2cm) -> Total ~16cm
\matrix (tab) [
    matrix of nodes,
    nodes in empty cells,
    column 1/.style={nodes={rowtitle, minimum width=3.2cm}},
    column 2/.style={nodes={stdcol, minimum width=3.2cm}},
    column 3/.style={nodes={stdcol, minimum width=3.2cm}},
    column 4/.style={nodes={stdcol, minimum width=3.2cm}},
    column 5/.style={nodes={camcol, minimum width=3.2cm}},
    row 1/.style={nodes={header, minimum width=3.2cm}},
    row sep=-\pgflinewidth,
    column sep=-\pgflinewidth,
] {
    Critère & \textbf{Régime Présidentiel} & \textbf{Régime Parlementaire} & \textbf{Régime Semi-présidentiel} & \textbf{CAMEROUN (1996/2008)} \\
    \textbf{Chef de l'État} & Président (Élu suffrage universel) & Monarque ou Président (Élu Parlement/Grands électeurs) & Président (Élu suffrage universel) & Président (Élu suffrage universel direct, 1 tour, 7 ans) \\
    \textbf{Chef du Gouvernement} & Président (fusionné) & Premier Ministre (Issu majorité parlementaire) & Premier Ministre (Nommé par Président) & Premier Ministre (Nommé par Président Art. 10) \\
    \textbf{Responsabilité Govt} & Aucune devant Parlement & Oui : Motion censure, Vote confiance, Eng. responsabilité & Oui : Responsable devant Assemblée Nationale & Responsable devant Assemblée Nationale (Art. 13) \\
    \textbf{Dissolution Assemblée} & Non (mandats fixes) & Oui (par Chef État sur avis PM) & Oui (Par Président) & Oui : Président peut dissoudre Assemblée (Art. 14) \\
    \textbf{Exemple type} & États-Unis (USA) & Royaume-Uni, Allemagne, Espagne & France (Vème République) & Cameroun (Constitution 1996) \\
};
\draw[thickline] (tab.north west) rectangle (tab.south east);
\draw[thickline] (tab.north west) -- (tab.north east);
\draw[thickline] (tab.south west) -- (tab.south east);
% Ligne séparation header
\draw[thickline] ([yshift=-\pgflinewidth]tab-1-1.south west) -- ([yshift=-\pgflinewidth]tab-1-5.south east);
% Lignes horizontales internes
\foreach \i in {2,3,4,5,6} {
    \draw[draw=popline, line width=0.5pt] ([yshift=-\pgflinewidth]tab-\i-1.south west) -- ([yshift=-\pgflinewidth]tab-\i-5.south east);
}
% Lignes verticales
\foreach \j in {2,3,4,5} {
    \draw[draw=popline, line width=0.5pt] ([xshift=-\pgflinewidth]tab-1-\j.north east) -- ([xshift=-\pgflinewidth]tab-6-\j.south east);
}
% Highlight colonne Cameroun (déjà fait via style camcol)
\end{tikzpicture}
\legende{Tableau comparatif des régimes politiques (Grille Duverger/Carcassonne). La colonne Cameroun (or) montre l'hybridation : légitimité présidentielle + responsabilité parlementaire encadrée.}
\end{popfigure}
\end{center}

\courssub{Analyse de la Constitution camerounaise : Un semi-présidentialisme « à la camerounaise »}

La Constitution du 18 janvier 1996 (révisée par la loi n°2008/001 du 14 avril 2008) instaure un régime \cle{semi-présidentiel} mais avec un déséquilibre net en faveur du Président. Analysons les articles clés.

\begin{propriete}[Architecture institutionnelle camerounaise (Art. 5, 6, 10, 12, 13, 14, 15)]
\begin{itemize}
    \item \textbf{Président de la République (Art. 5, 6) :} Chef de l'État, Chef des Armées, \textbf{détermine la politique de la Nation}. Élu au suffrage universel direct, \textbf{majorité relative (1 tour)}, mandat de \textbf{7 ans renouvelable une fois} (Art. 6, révisé 2008). Bénéficie d'\textbf{immunité présidentielle} (Art. 53). Pouvoir de nomination étendu : PM, Ministres, hauts fonctionnaires, magistrats (via CSM), dirigeants entreprises publiques, préfets.
    \item \textbf{Premier Ministre (Art. 10, 12, 13) :} Chef du Gouvernement. \textbf{Nommé et révoqué par le Président} (Art. 10). \textbf{Dirige l'action du Gouvernement} (Art. 12) et \textbf{est responsable devant l'Assemblée Nationale} (Art. 13). Contre-signe les actes présidentiels (sauf nomination PM, dissolution, référendum).
    \item \textbf{Parlement Bicaméral (Art. 14, 15) :} Assemblée Nationale (députés élus 5 ans) + Sénat (créé 2013, conseillers régionaux/municipaux + nommés Président). L'Assemblée a le dernier mot en cas de désaccord (navette, Art. 26).
    \item \textbf{Relations Exécutif/Législatif :} Président \textbf{peut dissoudre l'Assemblée} (Art. 14) — jamais utilisé depuis 1996. Gouvernement \textbf{responsable devant Assemblée} (motion censure Art. 34 Règlement AN, vote budget) — jamais renversé.
\end{itemize}
\end{propriete}

\begin{attention}[Piège majeur : Confondre « Premier Ministre » et « Régime Parlementaire »]
\textbf{Erreur fréquente :} « Il y a un Premier Ministre au Cameroun, donc c'est un régime parlementaire. »
\textbf{FAUX.} L'existence d'un PM ne suffit pas. En régime parlementaire (UK, Allemagne), le PM \emph{émane} de la majorité parlementaire (il en est l'expression) et le Chef de l'État est neutre/arbitre. Au Cameroun, le PM est \emph{choisi} par le Président (souvent hors majorité parlementaire stricte), le Président \textbf{détermine la politique} (Art. 5) et dispose de l'administration/armée. C'est un \textbf{régime présidentialisé} (ou semi-présidentiel « présidentialiste ») : le Président est le \cle{vrai chef de l'exécutif}, le PM est un \textbf{collaborateur direct} (subordonné), fusible politique.
\end{attention}

\begin{exemplebox}[Le « Présidentialisme » camerounais en pratique : La nomination de 2019]
En janvier 2019, le Président Paul Biya reconduit Joseph Dion Ngute comme Premier Ministre après l'élection présidentielle d'octobre 2018. Le RDPC (parti présidentiel) dispose de la majorité absolue à l'Assemblée. Pourtant, le PM n'est pas « issu » de cette majorité au sens parlementaire (chef de file de la majorité), il est \textbf{choisi personnellement} par le Président. Le gouvernement formé comprend des ministres de divers horizons, mais la cohésion politique dépend de la volonté présidentielle, non d'un contrat de coalition parlementaire. Cela illustre la \textbf{subordination du PM au Président}.
\end{exemplebox}

\begin{center}
\begin{popfigure}
\begin{tikzpicture}[
    node distance=8mm and 12mm,
    box/.style={rectangle, draw=popblue, thick, fill=popblueL, rounded corners, minimum height=10mm, minimum width=3.5cm, align=center, font=\small\bfseries},
    boxpm/.style={rectangle, draw=poporange, thick, fill=poporangeL, rounded corners, minimum height=10mm, minimum width=3.5cm, align=center, font=\small\bfseries},
    boxgvt/.style={rectangle, draw=popgreen, thick, fill=popgreenL, rounded corners, minimum height=10mm, minimum width=4cm, align=center, font=\small\bfseries},
    boxan/.style={rectangle, draw=poppurple, thick, fill=poppurpleL, rounded corners, minimum height=10mm, minimum width=3.5cm, align=center, font=\small\bfseries},
    arrow/.style={-Stealth, thick, draw=popdark},
    labelarrow/.style={midway, fill=white, font=\tiny, inner sep=1pt, text=popdark}
]
% Nœuds
\node[box, align=center] (peuple){PEUPLE\\Souvereigneté (Art. 2)};
\node[box, right=of peuple, align=center] (pres){PRÉSIDENT DE LA RÉPUBLIQUE\\Chef État, Armées, Politique Nation (Art 5,6)};
\node[boxan, below=of pres, align=center] (an){ASSEMBLÉE NATIONALE\\Députés élus 5 ans (Art 14)};
\node[boxpm, below=of peuple, align=center] (pm){PREMIER MINISTRE\\Chef Gouvernement (Art 10, 12)};
\node[boxgvt, right=of pm, align=center] (gvt){GOUVERNEMENT\\Ministres (Art 10, 12)};
\node[boxan, right=of gvt, align=center] (senat){SÉNAT\\Conseillers régionaux + Nomination PR (Art 14)};
% Flèches
\draw[arrow] (peuple) -- node[labelarrow] {Élection directe (1 tour)} (pres);
\draw[arrow] (peuple) -- node[labelarrow, left] {Élection députés} (an);
\draw[arrow] (pres) -- node[labelarrow] {Nomination / Révocation (Art 10)} (pm);
\draw[arrow] (pres) -- node[labelarrow, above] {Nomination Ministres (sur prop. PM)} (gvt);
\draw[arrow] (gvt) -- node[labelarrow] {Action administrative} (pm);
\draw[arrow, bend left=20] (gvt) to node[labelarrow, align=center]{Responsabilité politique\\Vote budget, Questions} (an);
\draw[arrow, bend right=20, draw=poporange] (an) to node[labelarrow, below, align=center]{Motion de censure (Art 34 Règl.)\\Jamais aboutie} (gvt);
\draw[arrow, draw=poppurple, dashed] (pres) -- node[labelarrow, above, align=center]{Dissolution possible (Art 14)\\Jamais utilisée} (an);
\draw[arrow] (pres) -- node[labelarrow] {Promulgation lois / Veto} (an);
\draw[arrow] (an) -- node[labelarrow] {Vote lois (Navette)} (senat);
\draw[arrow] (senat) -- node[labelarrow] {Vote lois} (an);
% Légende boîte
\node[draw=popdark, fill=popcream, rounded corners, font=\tiny, align=left, anchor=north east] at (current bounding box.north east) {
    \textbf{Légende :}\\
    \textcolor{popblue}{---} Flux légitimité / Nomination\\
    \textcolor{poporange}{---} Responsabilité politique (Gvt $\to$ AN)\\
    \textcolor{poppurple}{--- - -} Pouvoir de dissolution (PR $\to$ AN)
};
\end{tikzpicture}
\legende{Organigramme des flux de responsabilité politique au Cameroun. Le Président (bleu) est le pivot central : source de la nomination du PM/Gouvernement (orange/vert) et détenteur du pouvoir de dissolution sur l'Assemblée (violet). Le Gouvernement rend des comptes à l'Assemblée (orange), mais l'Assemblée n'a jamais renversé le Gouvernement.}
\end{popfigure}
\end{center}

\courssub{Les pouvoirs du Président : Une concentration exceptionnelle}

Au-delà de l'architecture formelle, la pratique et le texte constitutionnel confèrent au Président camerounais un éventail de pouvoirs qui dépasse largement le modèle français (Vème République), justifiant la qualification de \cle{régime présidentialisé}.

\begin{center}
\begin{popfigure}
\begin{tikzpicture}[
    mindmap,
    concept color=popblue,
    level 1 concept/.append style={level distance=4.5cm, sibling angle=90, font=\bfseries\small, minimum size=2.5cm},
    level 2 concept/.append style={level distance=3cm, sibling angle=45, font=\small, minimum size=2cm},
    level 3 concept/.append style={level distance=2.5cm, font=\tiny, minimum size=1.5cm},
    every node/.style={concept, execute at begin node=\hskip0pt},
    edge from parent/.style={draw=popblue!50, thick}
]
\node[concept, align=center]{Pouvoirs du\\ Président Camerounais}
    child[concept color=popgreen] { node[concept] {Exécutif Propre}
        child { node[concept] {Nomination PM / Ministres (Art 10)} }
        child { node[concept] {Nomination hauts fonct. (Préfets, DG, Ambassadeurs)} }
        child { node[concept] {Chef Armées / Sécurité (Art 6)} }
        child { node[concept] {Politique Nation (Art 5)} }
    }
    child[concept color=poporange] { node[concept] {Législatif Partiel}
        child { node[concept] {Initiative lois (concurrence)} }
        child { node[concept] {Promulgation / Veto suspensif} }
        child { node[concept] {Dissolution Assemblée (Art 14)} }
        child { node[concept] {Référendum (Art 28)} }
        child { node[concept] {Pouvoir réglementaire (Art 9)} }
    }
    child[concept color=poppurple] { node[concept] {Judiciaire / Institutionnel}
        child { node[concept] {Président CSM (Art 37) $\to$ Carrière magistrats} }
        child { node[concept] {Droit de grâce (Art 30)} }
        child { node[concept] {Nomination 30\% Sénateurs (Art 14)} }
        child { node[concept] {Nomination Conseil Constitutionnel} }
    }
    child[concept color=poppink] { node[concept] {Exceptionnels / Crise}
        child { node[concept] {Délégation de pouvoirs (Art 9 al. 3)} }
        child { node[concept] {État d'urgence / Siège (Art 20)} }
        child { node[concept] {Immunité présidentielle (Art 53)} }
    };
\end{tikzpicture}
\legende{Carte mentale des pouvoirs du Président camerounais. Quatre pôles : Exécutif (nomination massive), Législatif (initiative, dissolution, référendum, règlement), Judiciaire (CSM, grâce, nominations), Exceptionnels. Cette concentration explique le « présidentialisme » réel.}
\end{popfigure}
\end{center}

\courssub{Le Parlement camerounais : Bicaméralisme asymétrique}

\begin{aretenir}[Parlement camerounais (Art. 14, 15, 26)]
\begin{itemize}
    \item \textbf{Assemblée Nationale (Chambre basse) :} 180 députés, élus au suffrage universel direct (\textbf{scrutin majoritaire à un tour avec prime majoritaire} selon les départements), mandat 5 ans. \textbf{Seule chambre habilitée à renverser le Gouvernement} (motion de censure, Art. 13/34 Règl.) et à voter le budget en dernier ressort.
    \item \textbf{Sénat (Chambre haute, créé 2013) :} 100 sénateurs (10 par région). \textbf{70 élus} au suffrage indirect par les conseillers municipaux et régionaux (collège électoral), \textbf{30 nommés par le Président}. Mandat 5 ans.
    \item \textbf{Navette parlementaire :} En cas de désaccord, l'Assemblée a \textbf{le dernier mot} (Art. 26). Le Sénat a un rôle consultatif et de représentation des collectivités territoriales, mais pas de pouvoir de blocage définitif ni de contrôle du gouvernement.
\end{itemize}
\end{aretenir}

\begin{savaistu}[Le Sénat : Une chambre de « sages » sous influence présidentielle]
Institué par la loi n°2013/004 (application Art. 14 Const. 1996), le Sénat a tenu sa première session en juin 2013. La règle des \textbf{30 sénateurs nommés par le Président} (sur 100) garantit une majorité structurelle au pouvoir en place, même si l'opposition gagne les élections municipales. En 2023, le RDPC contrôle la quasi-totalité des sièges (élus + nommés). Le Sénat vote les lois, mais ne peut censurer le gouvernement. C'est un \textbf{bicaméralisme faible, asymétrique}, typique des régimes présidentialisés où la chambre haute sert de chambre de relecture et d'intégration des élites locales, non de contre-pouvoir.
\end{savaistu}

\courssub{Évaluation quantitative : La responsabilité parlementaire en pratique}

\begin{exoresolu}[Analyse de la responsabilité gouvernementale devant l'Assemblée (1992-2024)]
\textbf{Énoncé :} À partir des données ci-dessous, caractérisez la réalité du contrôle parlementaire au Cameroun.
\begin{center}
\begin{tabularx}{\linewidth}{|l|X|X|X|}\hline
\textbf{Législature} & \textbf{Gouvernements formés} & \textbf{Motions de censure déposées} & \textbf{Motions adoptées} \\ \hline
1992--1997 & 2 (Sadou Hayatou, Simon Achidi Achu) & 1 (Opposition, 1992) & 0 \\ \hline
1997--2002 & 2 (Achidi Achu, Musonge) & 0 & 0 \\ \hline
2002--2007 & 1 (Ephraïm Inoni) & 0 & 0 \\ \hline
2007--2013 & 1 (Philemon Yang) & 0 & 0 \\ \hline
2013--2018 & 1 (Philemon Yang) & 0 & 0 \\ \hline
2018--2020 & 1 (Joseph Dion Ngute) & 0 & 0 \\ \hline
2020--2025 & 1 (Joseph Dion Ngute) & 1 (SDF, 2022 - rejetée) & 0 \\ \hline
\end{tabularx}
\end{center}
\tcblower
\textbf{Solution :}
\begin{enumerate}
    \item \textbf{Stabilité gouvernementale extrême :} 7 gouvernements en 32 ans (durée moyenne $\approx$ 4,5 ans), alignée sur le mandat présidentiel (7 ans). Le remaniement suit l'élection présidentielle, pas le calendrier parlementaire.
    \item \textbf{Responsabilité théorique, inexistante en pratique :} 2 motions déposées en 32 ans, \textbf{0 adoptée}. L'Assemblée n'a jamais renversé le gouvernement.
    \item \textbf{Cause structurelle :} Majorité mécanique du RDPC à l'Assemblée (système électoral + discipline de parti) + Subordination du PM au Président (le Gouvernement ne « tombe » que si le Président le décide).
    \item \textbf{Conclusion :} Le critère « Responsabilité du gouvernement devant le Parlement » (critère parlementaire) est \textbf{formel mais inopérant}. Le régime fonctionne comme un \textbf{régime présidentiel} sur le plan de la stabilité de l'exécutif, malgré l'existence textuelle d'un PM responsable.
\end{enumerate}
\end{exoresolu}

\begin{exemplebox}[Chiffres clés : L'exécutif camerounais en 2024]
\begin{itemize}
    \item \textbf{Président :} Paul Biya (au pouvoir depuis 1982, 42 ans). Réélu 2018 (71,28\%).
    \item \textbf{Premier Ministre :} Joseph Dion Ngute (depuis 2019).
    \item \textbf{Gouvernement :} ~60 membres (Ministres, Ministres délégués, Secrétaires d'État). Turn-over faible aux postes régaliens (Défense, Intérieur, Finances, Justice souvent stables > 10 ans).
    \item \textbf{Assemblée Nationale :} 180 sièges. RDPC : $\approx$ 140-150 sièges (majorité absolue large). Opposition : $\approx$ 30-40 sièges (fragmentée).
    \item \textbf{Sénat :} 100 sièges. RDPC : $\approx$ 90+ sièges (grâce aux 30 nommés + victoire municipales).
\end{itemize}
\end{exemplebox}

\courssub{Synthèse : Qualification juridique et politique du régime camerounais}

Comment qualifier *in fine* le régime camerounais pour une dissertation ou un commentaire ?

\begin{definition}[Régime semi-présidentiel présidentialisé (ou « présidentialisme tempéré »)]
Système où les \textbf{institutions formelles} sont celles d'un régime semi-présidentiel (Président élu + PM responsable devant Parlement + Dissolution possible), mais où la \textbf{pratique politique} (majorité présidentielle stable, discipline de parti, contrôle de l'administration/justice par le Président, absence de cohabitation réelle) concentre l'effectivité du pouvoir exécutif entre les mains du Président de la République. Le Premier Ministre est un \textbf{chef de gouvernement d'exécution}, non de direction politique.
\end{definition}

\begin{attention}[Nuance pour le Bac : Éviter les étiquettes caricaturales]
\begin{itemize}
    \item Ne pas écrire : « Le Cameroun est une dictature » (jugement politique, hors grille d'analyse institutionnelle).
    \item Ne pas écrire : « Le Cameroun est une démocratie parlementaire » (faux institutionnellement).
    \item \textbf{Écrire :} « Le Cameroun est une République à régime \textbf{semi-présidentiel} (Constitution 1996), caractérisé par un \textbf{présidentialisme fort} (élection directe, détermination politique, nomination massive, immunité) et un \textbf{parlementarisme atténué} (PM responsable mais subordonné, Assemblée sans pouvoir de renversement effectif, Sénat sans pouvoir de censure). »
\end{itemize}
\end{attention}

\begin{exercice}[Entraînement Bac - Application de la grille Duverger]
\textbf{Sujet :} « Classer le régime politique camerounais à l'aide des critères de Maurice Duverger. »
\textbf{Consignes :} Utilisez la méthode en 4 questions (Section 3.1). Citez les articles constitutionnels pertinents. Concluez sur la nature « hybride » du régime.
\textit{(Corrigé type en fin de chapitre / section 6).}
\end{exercice}

\coursec{Rôle, finalité et fonctions de l'État : L'État-providence en question}

Nous avons vu dans la section précédente que le pouvoir politique peut s'exercer selon des formes de gouvernement variées (monarchie, république, régimes parlementaire ou présidentiel) et que le Cameroun est une République décentralisée. Mais une fois le pouvoir installé et légitime, une question demeure : \emph{que doit faire l'État ?} Doit-il se limiter à garantir l'ordre et la sécurité, ou doit-il aussi assurer le bien-être de ses citoyens ? C'est le débat entre l'\cle{État gendarme} et l'\cle{État-providence}. Cette section examine les fonctions de l'État, ses instruments (le budget, la décentralisation) et les défis que rencontre l'État camerounais.

\courssub{Les fonctions régaliennes : le cœur incompressible de l'État}

\begin{definition}[Fonctions régaliennes]
Les \cle{fonctions régaliennes} (du latin \emph{regalia}, « droits du roi ») sont les missions fondamentales que l'État ne peut déléguer à aucun autre acteur : la \cle{justice}, la \cle{police} et la \cle{défense} (sécurité intérieure et extérieure), la \cle{diplomatie}, la \cle{monnaie} et la \cle{législation} (fabriquer la loi).
\end{definition}

Commençons par une intuition. Imagine un quartier de Douala où deux voisins se disputent un terrain. Si l'État n'existait pas, chacun ferait justice lui-même : le plus fort gagnerait. L'État intervient précisément pour \emph{remplacer la vengeance privée par la justice publique}. C'est le sens du célèbre monopole défini par le sociologue allemand Max Weber.

\begin{propriete}[Le monopole de la violence physique légitime (Weber)]
Selon Max Weber (\emph{Le Savant et le Politique}, 1919), l'État est « une communauté humaine qui revendique, avec succès, le \cle{monopole de la violence physique légitime} ». Seuls la police, la gendarmerie, l'armée et la justice peuvent user légalement de la contrainte physique (arrestation, garde à vue, peine de prison). Toute violence privée est illégitime et punie par la loi.
\end{propriete}

Ce monopole n'est pas une tyrannie : il est encadré par le droit. Le policier ne peut user de la force que dans des conditions strictement définies ; le juge ne peut condamner qu'au terme d'un procès équitable. La violence de l'État est \emph{légitime} parce qu'elle est \emph{légale} et \emph{contrôlée}.

Au Cameroun, ces fonctions régaliennes s'incarnent concrètement :
\begin{itemize}
\item \textbf{Justice} : les tribunaux, la Cour suprême, le Conseil constitutionnel ;
\item \textbf{Sécurité intérieure} : la Police nationale (Délégation générale à la Sûreté nationale) et la Gendarmerie nationale ;
\item \textbf{Défense extérieure} : les Forces armées camerounaises (FAC), engagées notamment contre Boko Haram / ISWAP dans l'Extrême-Nord ;
\item \textbf{Diplomatie} : le MINREX (Ministère des Relations extérieures), les ambassades, la participation à l'ONU et à l'Union africaine ;
\item \textbf{Monnaie} : le Cameroun partage le franc CFA avec la CEMAC ; l'émission monétaire est assurée par la \cle{BEAC} (Banque des États de l'Afrique centrale), dont le siège est à Yaoundé ;
\item \textbf{Législation} : l'Assemblée nationale et le Sénat votent les lois.
\end{itemize}

\begin{attention}[Piège fréquent]
Ne confonds pas \emph{monopole de la violence légitime} et \emph{violence arbitraire}. Dire que l'État a le monopole de la violence ne signifie pas qu'il peut tout se permettre : la contrainte étatique est légitime seulement si elle respecte la loi et les droits de l'homme. Un État qui use de la force hors du droit devient un État de fait, non un État de droit.
\end{attention}

\courssub{Les fonctions sociales : la naissance de l'État-providence}

Après la crise économique de 1929 et la Seconde Guerre mondiale, une idée nouvelle s'impose : l'État ne doit pas seulement maintenir l'ordre, il doit aussi \emph{protéger les citoyens contre les risques de la vie} (maladie, chômage, vieillesse, pauvreté). L'économiste John Maynard \cle{Keynes} montre que l'État doit intervenir dans l'économie pour soutenir l'activité ; le Britannique William \cle{Beveridge} (rapport de 1942) propose un système de sécurité sociale « du berceau au tombeau ».

\begin{definition}[L'État-providence]
L'\cle{État-providence} (ou État social, \emph{Welfare State}) est une forme d'État qui assume la responsabilité du bien-être économique et social de ses citoyens par des politiques de \cle{redistribution}, de \cle{protection sociale} et de \cle{services publics} accessibles à tous : éducation, santé, logement, emploi, retraite.
\end{definition}

Au Cameroun, les fonctions sociales se traduisent par :
\begin{itemize}
\item \textbf{Éducation} : gratuité de l'enseignement primaire public (depuis 2000), construction de salles de classes, recrutement massif d'enseignants ;
\item \textbf{Santé} : hôpitaux publics, campagnes de vaccination, gratuité du traitement du paludisme simple pour les enfants de moins de 5 ans ;
\item \textbf{Protection sociale} : la \cle{CNPS} (Caisse nationale de prévoyance sociale) pour les retraites des salariés du secteur privé, et la \cle{CNAMGS} (Caisse nationale de l'assurance maladie universelle) pour la couverture santé des pauvres et vulnérables ;
\item \textbf{Logement et urbanisme} : la SIC (Société immobilière du Cameroun) et le MAETUR (Mission d'aménagement et d'équipement des terrains urbains et ruraux) construisent des logements sociaux ;
\item \textbf{Emploi} : le FNE (Fonds national de l'emploi) finance la formation et l'insertion des jeunes ;
\item \textbf{Infrastructures} : routes, énergie (barrages de Memve'ele, Nachtigal, Lom Pangar), eau potable.
\end{itemize}

\begin{exemplebox}[La CNAMGS en chiffres (données 2023-2024)]
La CNAMGS vise environ 4 millions de pauvres et de personnes vulnérables pris en charge par l'État (volet « Assurance Maladie Universelle » non contributif). La couverture effective atteint environ 60 \% de la cible, avec un budget alloué de l'ordre de 50 milliards de FCFA par an. Les défis majeurs restent l'identification fiable des bénéficiaires (registre social unique en cours de déploiement) et le paiement rapide des prestataires (hôpitaux, pharmacies), dont les retards découragent la prise en charge.
\end{exemplebox}

\begin{attention}[L'État-providence ne signifie pas « tout gratuit »]
Erreur fréquente : croire que l'État-providence distribue des biens gratuits. En réalité, il s'agit d'une \cle{gestion mutualisée des risques} (solidarité nationale) financée par l'impôt et les cotisations sociales : chacun contribue selon ses moyens et reçoit selon ses besoins. Sa \cle{soutenabilité financière} est donc cruciale : si les recettes baissent ou si les dépenses explosent, le système s'effondre.
\end{attention}

\courssub{Les fonctions économiques : réguler, planifier, investir}

L'État n'est pas seulement un protecteur social : il est aussi un \emph{acteur économique}. Trois grandes missions se distinguent.

\textbf{1. La régulation.} L'État fixe les règles du jeu économique : il sanctionne les ententes et les abus de position dominante (concurrence), encadre certains prix (carburant, transport urbain), protège les consommateurs et les travailleurs (Code du travail).

\textbf{2. La planification.} L'État trace les grandes orientations du développement.

\begin{savaistu}[De la DSCE à la SND30]
Le Cameroun a adopté la \cle{SND30} (Stratégie Nationale de Développement 2020--2030) comme cadre unique de planification, en remplacement de la DSCE (Document de Stratégie pour la Croissance et l'Emploi). Objectif affiché : faire du Cameroun un \emph{pays émergent à revenu intermédiaire à l'horizon 2035}, avec un PIB par habitant supérieur à 4 000 dollars. La SND30 repose sur la transformation structurelle de l'économie (industrialisation, agriculture modernisée) et l'amélioration de la gouvernance.
\end{savaistu}

\textbf{3. L'investissement public et les entreprises publiques.} L'État réalise directement des investissements : le \cle{PID} (Programme d'Infrastructures Durables), les routes (Yaoundé--Nsimalen, pénétrantes nord), les barrages hydroélectriques (Memve'ele, Nachtigal, Lom Pangar). Il possède ou contrôle aussi des entreprises publiques et parapubliques : la \cle{SONARA} (raffinerie de Limbé), \cle{CAMTEL} (télécommunications), \cle{CAMWATER} (eau potable), tandis que la distribution d'électricité est concédée à \cle{ENEO} (concession : l'État reste propriétaire du réseau mais en confie l'exploitation à un opérateur privé).

\courssub{Le budget de l'État : la politique chiffrée}

\begin{definition}[Budget de l'État]
Le \cle{budget} est l'acte qui prévoit et autorise, pour une année, l'ensemble des recettes et des dépenses de l'État. Il est voté par le Parlement (loi de finances) : c'est la traduction chiffrée des choix politiques de la nation.
\end{definition}

Les \textbf{recettes} proviennent principalement des impôts et taxes (TVA, impôt sur les sociétés), des droits de douane, des revenus pétroliers (via la SNH) et des dons et emprunts extérieurs. Les \textbf{dépenses} se divisent en dépenses de \emph{fonctionnement} (salaires du personnel, achats courants, subventions) et dépenses d'\emph{investissement} (équipements, infrastructures), auxquelles s'ajoutent le service de la dette et les transferts aux collectivités territoriales décentralisées (CTD).

\begin{popfigure}
\begin{center}
\begin{tikzpicture}
\begin{axis}[
  ybar, bar width=0.5cm,
  width=0.48\linewidth, height=6cm,
  symbolic x coords={Pétrole,Impôts,Douanes,Dons/Emprunts,Autres},
  xtick=data,
  ymin=0, ymax=2800,
  ylabel={Milliards de FCFA},
  title={Recettes principales (LFI 2024)},
  nodes near coords, every node near coord/.append style={font=\tiny, rotate=90, anchor=west, yshift=2pt},
  enlarge x limits=0.15,
  x tick label style={rotate=30, anchor=east},
  tick label style={font=\scriptsize}]
\addplot[fill=popblue,draw=popdark] coordinates {(Pétrole,850) (Impôts,2500) (Douanes,1050) (Dons/Emprunts,1800) (Autres,540)};
\end{axis}
\end{tikzpicture}
\hfill
\begin{tikzpicture}
\begin{axis}[
  ybar, bar width=0.5cm,
  width=0.48\linewidth, height=6cm,
  symbolic x coords={Personnel,Investissement,Dette,Transferts CTD,Fonct. hors pers.},
  xtick=data,
  ymin=0, ymax=2800,
  ylabel={Milliards de FCFA},
  title={Dépenses principales (LFI 2024)},
  nodes near coords, every node near coord/.append style={font=\tiny, rotate=90, anchor=west, yshift=2pt},
  enlarge x limits=0.15,
  x tick label style={rotate=30, anchor=east},
  tick label style={font=\scriptsize}]
\addplot[fill=poporange,draw=popdark] coordinates {(Personnel,1550) (Investissement,1800) (Dette,1050) (Transferts CTD,500) (Fonct. hors pers.,1100)};
\end{axis}
\end{tikzpicture}
\end{center}
\legende{Structure indicative du Budget général du Cameroun 2024 (Loi de finances n°2023/017, ~6 740 milliards FCFA). À gauche : recettes principales. À droite : dépenses principales. Les montants sont arrondis. Source : Minfi.}
\end{popfigure}
Le budget obéit à des principes classiques de droit financier public :
\begin{itemize}
\item \textbf{Annualité} : le budget est voté et exécuté pour un an (1$^{\text{er}}$ janvier -- 31 décembre) ;
\item \textbf{Unité} : toutes les recettes et dépenses figurent dans un document unique (le budget général) ;
\item \textbf{Universalité} : toutes les recettes et toutes les dépenses doivent y figurer, sans compensation (pas de netting) ;
\item \textbf{Non-affectation} : en principe, aucune recette n'est réservée à une dépense précise (toutes les recettes financent toutes les dépenses) ; des dérogations existent (comptes d'affectation spéciale) ;
\item \textbf{Sincérité} : les prévisions doivent être honnêtes et réalistes, sans sous-évaluation volontaire des dépenses ni surévaluation des recettes.
\end{itemize}
\courssub{La décentralisation : rapprocher l'action publique du citoyen}
La Constitution de 1996 (révisée en 2008) fait du Cameroun une « République décentralisée ». La \cle{décentralisation} consiste à transférer des compétences et des ressources de l'État central vers les \cle{collectivités territoriales décentralisées} (CTD) : les \textbf{Régions} (Conseil régional, président du conseil régional) et les \textbf{Communes} (Conseil municipal, maire élu). Les lois de 2019/2020 précisent l'organisation et les compétences.
Les compétences transférées aux communes portent notamment sur les écoles maternelles et primaires (construction, entretien, équipements), les centres de santé intégrés, les marchés, la voirie urbaine et l'assainissement. Les ressources suivent : \cle{centimes additionnels communaux} (part des impôts d'État reversée aux communes) et la \cle{DGD} (Dotation générale de décentralisation), versée par le FEICOM (Fonds d'équipement et d'intervention intercommunal).
\begin{popfigure}
\begin{center}
\begin{tikzpicture}[
  box/.style={draw=popdark, thick, rounded corners, align=center, minimum width=3.8cm, minimum height=1.3cm, font=\small},
  fleche/.style={-{Stealth[length=3mm]}, very thick, popblue}]
\node[box, fill=popblueL, align=center] (etat) at (0,3.5){\textbf{ÉTAT CENTRAL}\\ (Niveau national)\\ Ministères, Parlement, Présidence};
\node[box, fill=popgreenL, align=center] (region) at (-4,0){\textbf{RÉGION}\\ Conseil régional\\ (Développement économique,\\ aménagement du territoire)};
\node[box, fill=poporangeL, align=center] (commune) at (4,0){\textbf{COMMUNE}\\ Conseil municipal, Maire\\ (Services de proximité)};
\draw[fleche] (etat.south west) -- (region.north) node[midway, above left, font=\scriptsize, align=center, text=popdark]{Transfert de\\ compétences + Ressources};
\draw[fleche] (etat.south east) -- (commune.north) node[midway, above right, font=\scriptsize, align=center, text=popdark]{Transfert de\\ compétences + Ressources\\(DGD, Centimes add.)};
\node[draw=popmuted, dashed, rounded corners, align=left, font=\scriptsize, text width=5cm, fill=popcream] at (0,-2.6) {\textbf{Compétences communales transférées :}\\ écoles maternelles/primaires, centres de santé, marchés, voirie/assainissement, état civil, action sociale, environnement local.};
\draw[popmuted, dashed] (region.south) -- (-1.5,-1.8);
\draw[popmuted, dashed] (commune.south) -- (1.5,-1.8);
\end{tikzpicture}
\end{center}
\legende{Schéma de la décentralisation camerounaise : l'État central transfère compétences et ressources (DGD, centimes additionnels via FEICOM) aux Régions et aux Communes, responsables des services publics de proximité.}
\end{popfigure}

\courssub{Les défis de l'État camerounais}

L'ambition de l'État-providence se heurte au Cameroun à des obstacles majeurs :
\begin{itemize}
\item \textbf{La corruption} : l'indice de perception de la corruption (CPI) de Transparency International classe régulièrement le Cameroun parmi les pays les plus mal notés (rang 140/180 en 2023). La \cle{CONAC} (Commission nationale anti-corruption) et l'opération « Épervier » (contrôle des comptes publics) luttent contre le détournement des deniers publics, mais la corruption détourne les ressources qui devraient financer écoles et hôpitaux ;
\item \textbf{La dette publique} : elle atteint environ 45 \% du PIB (seuil de vigilance CEMAC : 70 \%) ; son service (intérêts et remboursements) réduit la marge de manœuvre budgétaire pour l'investissement ;
\item \textbf{Le secteur informel} : il fournit environ 90 \% des emplois ; or ces travailleurs échappent largement à l'impôt et à la protection sociale, ce qui affaiblit à la fois les recettes de l'État et la couverture des citoyens ;
\item \textbf{L'insécurité} : les attaques de Boko Haram / ISWAP dans l'Extrême-Nord et la crise dans les régions du Nord-Ouest et du Sud-Ouest (NOSO) imposent des dépenses de défense élevées et freinent le développement local ;
\item \textbf{L'objectif d'émergence 2035} : devenir un pays à revenu intermédiaire suppose une croissance forte, durable et partagée, ce qui exige une amélioration drastique du climat des affaires et de la gouvernance.
\end{itemize}

\courssub{Justice sociale et égalité des chances : Rawls et Sen}

Pourquoi l'État doit-il redistribuer ? Deux philosophes éclairent la finalité de l'État-providence.

John \cle{Rawls} (\emph{Théorie de la justice}, 1971) propose une expérience de pensée : imaginons des individus choisissant les règles de la société derrière un « \cle{voile d'ignorance} », sans savoir s'ils seront riches ou pauvres, sains ou malades. Ils choisiraient rationnellement deux principes : l'égalité des libertés fondamentales pour tous, et le \cle{principe de différence} : les inégalités ne sont justes que si elles profitent \emph{d'abord aux plus défavorisés}. L'État-providence se justifie ainsi : il corrige les inégalités de naissance pour garantir une véritable \cle{égalité des chances}.

Amartya \cle{Sen} (prix Nobel d'économie 1998) va plus loin avec la théorie des \cle{capabilités} (ou \emph{capacités}) : le développement n'est pas seulement la croissance du revenu, c'est l'élargissement des \emph{libertés réelles} des personnes — être en bonne santé, être éduqué, participer à la vie de la cité. Un revenu ne suffit pas si l'école ou l'hôpital est inaccessible (absence de \emph{capabilité}).

Ces théories s'appliquent directement aux politiques camerounaises : la \textbf{gratuité de l'enseignement primaire} (2000) vise l'égalité des chances scolaires (Rawls) et la capabilité « être éduqué » (Sen) ; la \textbf{CNAMGS} vise la couverture sanitaire universelle (capabilité « être en bonne santé ») ; les \textbf{bourses d'excellence} permettent aux meilleurs élèves pauvres de poursuivre leurs études (principe de différence).

\begin{methode}[Analyser une politique publique]
Pour étudier une politique publique (exigible en dissertation et en étude de cas) :
\begin{enumerate}
\item \textbf{Identifier le problème public} : quel besoin social non satisfait (ex. : faible accès aux soins des pauvres) ?
\item \textbf{Analyser les acteurs} : État (niveau central/déconcentré), collectivités territoriales, ONG, secteur privé, usagers/bénéficiaires ;
\item \textbf{Étudier les instruments} : loi, règlement, budget, norme, incitation fiscale, subvention, conventionnement ;
\item \textbf{Évaluer les résultats} : comparer les indicateurs obtenus aux objectifs fixés (ex. : taux de couverture CNAMGS vs cible ; taux de scolarisation vs objectif ODD4).
\end{enumerate}
\end{methode}

\begin{popfigure}
\begin{center}
\begin{tikzpicture}[
  etape/.style={draw=popdark, thick, rounded corners, fill=popblueL, align=center, minimum width=2.6cm, minimum height=1.1cm, font=\scriptsize},
  fl/.style={-{Stealth[length=2.5mm]}, thick, poporange}]
\node[etape, align=center] (a) at (90:2.8){\textbf{1. Mise à l'agenda}\\ Problème : coût\\ des soins inaccessible\\ pour les pauvres};
\node[etape, align=center] (b) at (18:2.8){\textbf{2. Formulation}\\ Loi 2017/014 créant\\ l'Assurance Maladie\\ Universelle (AMU)};
\node[etape, align=center] (c) at (-54:2.8){\textbf{3. Décision}\\ Création de la\\ CNAMGS (2017)\\ Décrets d'application};
\node[etape, align=center] (d) at (-126:2.8){\textbf{4. Mise en œuvre}\\ Immatriculation,\\ carte d'assuré,\\ prise en charge\\ tiers payant};
\node[etape, align=center] (e) at (162:2.8){\textbf{5. Évaluation}\\ Couverture $\approx$ 60\% cible,\\ délais paiement,\\ satisfaction usagers};
\draw[fl] (a) -- (b);
\draw[fl] (b) -- (c);
\draw[fl] (c) -- (d);
\draw[fl] (d) -- (e);
\draw[fl] (e) -- (a);
\node[font=\scriptsize, align=center, text=popmuted] at (0,0) {Cycle continu\\ (réajustements)};
\end{tikzpicture}
\end{center}
\legende{Le cycle de la politique publique appliqué à la CNAMGS/AMU : mise à l'agenda, formulation, décision, mise en œuvre, évaluation (qui relance un nouveau cycle).}
\end{popfigure}

\begin{exoresolu}[Étude de cas : la gratuité de l'enseignement primaire]
\textbf{Énoncé.} En appliquant la méthode d'analyse d'une politique publique, montre que la gratuité de l'enseignement primaire au Cameroun (2000) illustre à la fois la fonction sociale de l'État et le principe de différence de Rawls.
\tcblower
\textbf{Solution détaillée.}
\begin{enumerate}
\item \emph{Problème public} : avant 2000, les frais de scolarité et les charges indirectes (tenues, fournitures, cotisations APEE) excluaient de l'école les enfants des familles pauvres, surtout en milieu rural. Taux de scolarisation bas, inégalités filles/garçons marquées.
\item \emph{Acteurs} : l'État (MINEDUB, MINESEC, MINFI), les communes (maîtres d'ouvrage délégués pour l'entretien/construction), les parents d'élèves, les partenaires techniques et financiers (UNICEF, Banque mondiale, AFD).
\item \emph{Instruments} : décision présidentielle de suppression des frais de scolarité (2000), subventions de fonctionnement aux écoles publiques (subvention par élève), recrutement d'enseignants contractuels et fonctionnaires, programmes de construction de salles de classes (PNDP, BIP).
\item \emph{Résultats} : forte hausse du taux brut de scolarisation (TBS) primaire (dépassant 100 \% du fait des surâges/redoublants), parité filles/garçons améliorée. \textbf{Limites} : persistance de coûts indirects (fournitures, cantines, APEE), surcharge des classes, disparités régionales (Nord vs Sud), qualité des apprentissages (PASEC) encore perfectible.
\end{enumerate}
\emph{Lien philosophique} : cette mesure relève de la fonction sociale de l'État (service public d'éducation accessible à tous, financé par l'impôt). Elle illustre le \textbf{principe de différence de Rawls} : la gratuité profite \emph{d'abord aux plus défavorisés}, qui seuls étaient exclus par le paiement ; elle réalise ainsi une égalité des chances \emph{réelle} (accès effectif), et non seulement \emph{formelle} (droit abstrait). Au sens de \textbf{Sen}, elle accroît la \textbf{capabilité « être éduqué »}, condition de l'exercice de toutes les autres libertés (choisir son métier, participer à la vie citoyenne, s'extraire de la pauvreté).
\end{exoresolu}

\begin{aretenir}[L'essentiel de la section]
\begin{itemize}
\item Les fonctions régaliennes (justice, police, défense, diplomatie, monnaie, législation) sont le cœur incompressible de l'État, fondées sur le monopole de la violence légitime (Weber).
\item L'État-providence (Keynes, Beveridge) assume le bien-être social par la redistribution et les services publics, financés par l'impôt et les cotisations — non par la « gratuité » magique.
\item L'État agit sur l'économie par la régulation, la planification (SND30) et l'investissement public ; son budget (LFI) est la traduction chiffrée de ses choix politiques, soumise à des principes stricts (annualité, unité, universalité, sincérité).
\item La décentralisation (lois 2019/2020) transfère compétences et ressources (DGD, centimes additionnels) aux Communes et Régions pour rapprocher le service public du citoyen.
\item La justice sociale (Rawls : principe de différence / voile d'ignorance ; Sen : capabilités / libertés réelles) donne sa finalité morale à l'État-providence : garantir l'égalité réelle des chances et le développement humain.
\end{itemize}
\end{aretenir}

\coursec{TP Intégré : Cartographie institutionnelle et Analyse de l'action publique locale}

Après avoir étudié le rôle et la finalité de l'État, et interrogé le modèle de l'État-providence, il devient nécessaire de \emph{vérifier} ces notions sur le terrain. La philosophie politique ne se contente pas de spéculer : elle observe, mesure et juge l'action publique concrète. Ce TP intégré vous propose de passer de la théorie à l'enquête : cartographier les institutions, dessiner la décentralisation, lire un budget, analyser un service public local et débattre d'une réforme constitutionnelle. Autant de gestes qui transforment le citoyen spectateur en citoyen éclairé.

\courssub{TP 1 : Construire l'organigramme des institutions de la République}

\textbf{Problématique.} Qui détient le pouvoir au Cameroun, et comment ces pouvoirs s'articulent-ils ?

\textbf{Intuition.} Un État est comme un organisme : il possède des organes spécialisés (décider, exécuter, juger, contrôler) reliés entre eux. L'organigramme rend visible cette architecture invisible.

\begin{methode}[Démarche du TP 1]
\begin{enumerate}
\item Recenser les institutions à partir de la Constitution du 18 janvier 1996 (révisée en 2008) : Présidence de la République, Gouvernement, Parlement (Assemblée nationale et Sénat), Conseil constitutionnel, Cour suprême.
\item Ajouter les organes de contrôle : CONSUPE (Contrôle supérieur de l'État), CONAC (Commission nationale anti-corruption), Cour des Comptes.
\item Classer chaque institution selon sa fonction : exécutive, législative, juridictionnelle, de contrôle, de régulation constitutionnelle.
\item Relier par des flèches les relations de tutelle, de nomination, de contrôle ou de saisine.
\item Vérifier la cohérence : chaque pouvoir est-il contrôlé par un autre ? La séparation est-elle effective ?
\end{enumerate}
\end{methode}

\begin{popfigure}
\begin{tikzpicture}[
box/.style={draw=popblue, fill=popblueL, rounded corners, text width=3.0cm, align=center, font=\small, minimum height=0.9cm},
ctrl/.style={draw=poporange, fill=poporangeL, rounded corners, text width=2.5cm, align=center, font=\small, minimum height=0.9cm},
const/.style={draw=poppurple, fill=poppurpleL, rounded corners, text width=2.8cm, align=center, font=\small, minimum height=0.9cm},
rel/.style={-{Stealth}, thick, popdark},
relc/.style={-{Stealth}, thick, poppurple, dashed}]
\node[box, align=center] (pr) at (0,5){\textbf{Président de la République}\\ Chef de l'État};
\node[box, align=center] (gouv) at (-4.5,3){\textbf{Gouvernement}\\ Premier Ministre};
\node[box, align=center] (parl) at (0,3){\textbf{Parlement}\\ Assemblée nationale + Sénat};
\node[box, align=center] (just) at (4.5,3){\textbf{Cour suprême}\\ Ordres judiciaire, adm., comptable};
\node[const, align=center] (cc) at (0,1){\textbf{Conseil constitutionnel}\\ Régularité lois/élections};
\node[ctrl] (consupe) at (-4.5,1) {CONSUPE};
\node[ctrl] (conac) at (-1.5,-1) {CONAC};
\node[ctrl] (cdc) at (4.5,1) {Cour des Comptes};
% Relations
\draw[rel] (pr) -- (gouv) node[midway, left, font=\scriptsize]{nomme / révoque};
\draw[rel] (pr) -- (just) node[midway, right, font=\scriptsize]{nomme magistrats (sur avis CSM)};
\draw[rel] (pr) -- (cc) node[midway, right, font=\scriptsize]{nomme 3 membres};
\draw[rel] (parl) -- (gouv) node[midway, right, font=\scriptsize]{contrôle (qst, interpellation)};
\draw[rel] (parl) -- (cc) node[midway, left, font=\scriptsize]{élit 3 membres / saisine};
\draw[relc] (cc) -- (parl) node[midway, left, font=\scriptsize]{contrôle constitutionnalité lois};
\draw[relc] (cc) -- (pr) node[midway, right, font=\scriptsize]{proclame élection};
\draw[rel, dashed, poporange] (consupe) -- (gouv) node[midway, left, font=\scriptsize]{contrôle admin.};
\draw[rel, dashed, poporange] (cdc) -- (gouv) node[midway, right, font=\scriptsize]{audit finances};
\draw[rel, dashed, poporange] (conac) -- (parl) node[midway, below, font=\scriptsize]{rapporte};
\draw[rel, dashed, poporange] (conac) -- (gouv) node[midway, below, font=\scriptsize]{enquête};
\end{tikzpicture}
\legende{Organigramme simplifié des institutions de la République du Cameroun (Constitution 1996, révisée 2008). Bleu : pouvoirs classiques ; Violet : régulation constitutionnelle (Conseil constitutionnel) ; Orange : organes de contrôle (flèches pointillées = contrôle/saisine).}
\end{popfigure}

\begin{savaistu}[Automatiser l'organigramme]
On peut générer automatiquement cet organigramme à partir d'un fichier de données CSV (colonnes : \emph{Institution, Type, Tutelle, Effectifs}). Un petit script Python lit le fichier ligne par ligne et écrit pour chaque institution un nœud TikZ (\texttt{\textbackslash node}) placé selon son type, puis trace les flèches de tutelle. C'est le principe des outils professionnels de visualisation administrative : la donnée d'abord, le dessin ensuite.
\end{savaistu}

\begin{definition}[Reddition des comptes]
La \cle{reddition des comptes} est l'obligation pour tout gestionnaire de deniers publics de justifier leur emploi devant l'organe compétent (Cour des Comptes, Conseil régional ou communal). Elle est un pilier de la \cle{bonne gouvernance} : nul ne peut administrer l'argent des citoyens sans en rendre compte.
\end{definition}

\courssub{TP 2 : Croquis de la décentralisation au Cameroun}

\textbf{Problématique.} Comment l'État rapproche-t-il le pouvoir des citoyens ?

\textbf{Intuition.} Gouverner plus de 25 millions d'habitants depuis Yaoundé seul serait inefficace : l'État délègue donc des compétences à des collectivités proches du terrain. La carte rend visible cette répartition.

\begin{methode}[Réussir un croquis géographique]
\begin{enumerate}
\item Tracer un \textbf{fond de carte précis} : frontières, chefs-lieux de région.
\item Choisir les \textbf{symboles} adaptés : ponctuels (point = commune), linéaires (flèche = transfert), surfaciques (hachure = région).
\item Organiser la \textbf{légende} par thèmes (découpage, flux, réalisations).
\item Rédiger un \textbf{titre informatif} : Quoi ? Où ? Quand ? Source.
\item Soigner la \textbf{propreté} : traits nets, couleurs limitées, orientation (nord) et échelle indicative.
\end{enumerate}
\end{methode}

\begin{attention}[Erreur fréquente en TP cartographie]
Ne \textbf{surchargez pas} la carte ! Ne représentez que l'essentiel pour répondre à la problématique. Hiérarchisez l'information : l'élément principal (ici, les transferts de compétences) doit dominer visuellement ; l'élément secondaire (réalisations locales) reste discret. Une carte illisible est une carte fausse, car elle cache l'essentiel.
\end{attention}

\begin{popfigure}
\begin{tikzpicture}[scale=0.85]
% Contour très simplifié du Cameroun (polygone approximatif)
\draw[thick, popdark, fill=popcream] (1.2,0) -- (3.2,0.2) -- (4.6,1.4) -- (4.4,3.2) -- (5.2,4.6) -- (4.2,6.4) -- (2.8,7.6) -- (1.6,6.8) -- (0.6,5.2) -- (0.2,3.4) -- (0.6,1.6) -- cycle;
% Régions (hachurées, schématiques - 10 régions)
\fill[pattern=north east lines, pattern color=popblue!60] (0.4,4.8) -- (1.6,6.8) -- (2.8,7.6) -- (4.2,6.4) -- (3.0,5.0) -- cycle; % Nord/Adamaoua/Est approx
\fill[pattern=dots, pattern color=popgreen!70] (0.6,1.6) -- (0.2,3.4) -- (2.2,3.0) -- (2.4,1.2) -- cycle; % Sud/Ouest/Littoral approx
% Communes (points - 374 communes)
\foreach \p in {(1.0,2.0),(1.6,2.6),(2.2,2.0),(1.2,3.2),(2.8,3.4),(3.4,2.6),(3.8,3.6),(2.0,4.2),(3.0,4.6),(1.4,5.4),(2.4,6.0),(3.4,6.6),(4.2,5.2),(4.4,2.0),(0.9,4.4)}
  \fill[popink] \p circle (0.05);
% Chefs-lieux de région (carrés orange)
\fill[poporange] (1.8,1.6) rectangle (1.98,1.78); \node[font=\scriptsize, right] at (2.0,1.7) {Yaoundé (Centre)};
\fill[poporange] (0.8,2.4) rectangle (0.98,2.58); \node[font=\scriptsize, left] at (0.7,2.5) {Douala (Littoral)};
\fill[poporange] (3.6,3.0) rectangle (3.78,3.18); \node[font=\scriptsize, right] at (3.8,3.1) {Bafoussam (Ouest)};
\fill[poporange] (2.6,6.0) rectangle (2.78,6.18); \node[font=\scriptsize, right] at (2.8,6.1) {Ngaoundéré (Adamaoua)};
\fill[poporange] (4.0,5.0) rectangle (4.18,5.18); \node[font=\scriptsize, right] at (4.2,5.1) {Bertoua (Est)};
\fill[poporange] (1.2,5.8) rectangle (1.38,5.98); \node[font=\scriptsize, left] at (1.1,5.9) {Bamenda (N-O)};
\fill[poporange] (0.5,4.0) rectangle (0.68,4.18); \node[font=\scriptsize, left] at (0.4,4.1) {Buea (S-O)};
\fill[poporange] (3.0,2.0) rectangle (3.18,2.18); \node[font=\scriptsize, right] at (3.2,2.1) {Ebolowa (Sud)};
\fill[poporange] (4.8,3.8) rectangle (4.98,3.98); \node[font=\scriptsize, right] at (5.0,3.9) {Maroua (Extrême-Nord)};
\fill[poporange] (2.0,7.0) rectangle (2.18,7.18); \node[font=\scriptsize, right] at (2.2,7.1) {Garoua (Nord)};
% Flèches de transferts (État -> CTD)
\draw[-{Stealth}, very thick, poppurple] (1.8,1.8) -- (2.4,5.6) node[midway, right, font=\scriptsize]{compétences};
\draw[-{Stealth}, very thick, poppurple] (1.8,1.8) -- (4.0,3.4) node[midway, above, font=\scriptsize]{ressources};
% Réalisations (pictogrammes simples)
\draw[popgreen, thick] (3.6,4.2) rectangle (3.9,4.5); \node[font=\scriptsize, below] at (3.75,4.15) {marché};
\draw[popgreen, thick] (1.0,0.9) -- (1.35,0.9) -- (1.35,1.2) -- (1.0,1.2) -- cycle; \node[font=\scriptsize, below] at (1.17,0.85) {école};
% Nord
\draw[-{Stealth}, thick] (5.6,6.6) -- (5.6,7.4) node[above, font=\scriptsize]{N};
% Échelle indicative
\draw[thick] (0.2,0.2) -- (1.2,0.2) node[midway, below, font=\scriptsize]{~200 km};
\end{tikzpicture}
\legende{Croquis schématique de la décentralisation au Cameroun (Loi n°2019/024). Point noir = commune (374 communes) ; hachures/points = 10 régions ; carré orange = chef-lieu de région ; flèche violette épaisse = transferts de compétences et de ressources (dotations) de l'État vers les CTD ; pictogrammes verts = réalisations locales concrètes. Tracé approximatif, non contractuel.}
\end{popfigure}

\begin{exemplebox}[Lecture du croquis]
Sur le croquis, la flèche Yaoundé $\rightarrow$ région du Nord (Garoua) symbolise le transfert de compétences (éducation de base, santé, voirie rurale) et de ressources (dotations budgétaires, fiscalité locale) vers les Conseils régionaux. Le pictogramme « marché » près de Bafoussam illustre une réalisation concrète financée par la commune : la décentralisation n'est pas une abstraction, c'est un marché construit, une école rénovée, une piste réparée.
\end{exemplebox}

\courssub{TP 3 : Lire le budget de l'État — investissement contre fonctionnement}

\textbf{Problématique.} L'État camerounais fait-il un effort suffisant pour le développement ?

\textbf{Intuition.} Un budget révèle les priorités réelles d'un État, mieux que les discours. Dépenser pour fonctionner (salaires, carburant, biens et services) maintient la machine ; dépenser pour investir (routes, hôpitaux, écoles, barrages) prépare l'avenir.

\begin{experience}[Exploitation de données chiffrées (MINFI / INS / Open Data)]
\textbf{Matériel} : données des Lois de finances 2019--2023, tableur ou pgfplots.\\
\textbf{Protocole} :
\begin{enumerate}
\item Relever pour chaque année le budget total et la part des dépenses d'investissement (BCI - Budget des Investissements Publics).
\item Calculer le ratio : $\frac{\text{Investissement}}{\text{Budget total}} \times 100$.
\item Tracer l'évolution sur un graphique.
\item Comparer à l'objectif de la SND30 (Stratégie Nationale de Développement 2020--2030) : porter la part de l'investissement public à plus de 30\,\% du budget total.
\end{enumerate}
\textbf{Observations} (données arrondies issues des Lois de finances) :
\begin{center}
\begin{tabularx}{\linewidth}{|l|X|X|X|X|X|}\hline
\rowcolor{popblueL} Année & 2019 & 2020 & 2021 & 2022 & 2023 \\ \hline
Budget total (milliards FCFA) & 4\,850 & 4\,950 & 4\,700 & 5\,100 & 5\,800 \\ \hline
Investissement (milliards FCFA) & 1\,067 & 891 & 940 & 1\,071 & 1\,334 \\ \hline
Ratio Invest./Budget & 22\,\% & 18\,\% & 20\,\% & 21\,\% & 23\,\% \\ \hline
\end{tabularx}
\end{center}
\textbf{Interprétation} : le ratio reste durablement inférieur à l'objectif de 30\,\% de la SND30. La chute de 2020 (18\,\%) s'explique par la crise de Covid-19, qui a gonflé les dépenses de fonctionnement (santé, soutien d'urgence, baisse des recettes). La remontée lente après 2021 (+5 points en 3 ans) reste insuffisante au regard de l'ambition d'émergence à l'horizon 2035.
\end{experience}

\begin{popfigure}
\begin{tikzpicture}
\begin{axis}[
width=0.9\linewidth, height=6.5cm,
xlabel={Année}, ylabel={Part de l'investissement (\%)},
xtick={2019,2020,2021,2022,2023},
ymin=0, ymax=35,
grid=major, grid style={popline!40},
legend style={at={(0.5,-0.3)}, anchor=north, font=\small, legend columns=2},
mark size=2.5pt]
\addplot[popcurve, mark=*, mark options={fill=popblue, draw=popblue}] coordinates {(2019,22) (2020,18) (2021,20) (2022,21) (2023,23)};
\addlegendentry{Ratio observé (MINFI)}
\addplot[poporange, dashed, thick, domain=2019:2023, samples=2] {30};
\addlegendentry{Objectif SND30 : 30\,\%}
\end{axis}
\end{tikzpicture}
\legende{Évolution de la part des dépenses d'investissement dans le budget de l'État (2019--2023, en \,\%). Source : Lois de finances, MINFI. La courbe reste sous la ligne d'objectif SND30 : l'effort de développement est réel mais insuffisant pour atteindre l'émergence.}
\end{popfigure}

\begin{exoresolu}[Calculer et interpréter un taux de variation]
Énoncé : la part de l'investissement passe de 18\,\% (2020) à 23\,\% (2023). Calculer le taux de variation et conclure philosophiquement.
\tcblower
Solution : $t = \dfrac{23-18}{18} \times 100 \approx 27,8\,\%$. La part de l'investissement a augmenté d'environ 27,8\,\% en trois ans. Interprétation : la reprise est nette, mais le niveau absolu (23\,\%) demeure en dessous du seuil SND30 (30\,\%). Philosophiquement, cela illustre la distinction entre \emph{finalité proclamée} (l'émergence) et \emph{moyens effectifs} : un État se juge à l'adéquation entre ses fins et ses actes budgétaires. Le budget est la loi la plus politique qui soit : il dit la vérité du pouvoir.
\end{exoresolu}

\begin{savaistu}[Le Budget Citoyen]
Chaque année, le MINFI publie un \cle{Budget Citoyen} : une version simplifiée et illustrée de la Loi de finances, accessible à tous. C'est un outil de transparence démocratique — et le document idéal pour ce TP : nul besoin d'être expert pour vérifier ce que l'État fait de l'argent public.
\end{savaistu}

\courssub{TP 4 : Analyse de cas — la gestion des ordures ménagères à Douala et Yaoundé}

\textbf{Problématique.} Qui est responsable quand les poubelles débordent ?

\begin{experience}[Grille d'analyse d'un service public local]
\textbf{Observation} : accumulation d'ordures dans certains quartiers de Douala et de Yaoundé, plaintes des riverains, risques sanitaires.\\
\textbf{Analyse en quatre questions} :
\begin{enumerate}
\item \textbf{Compétence} : la collecte, le transport et le traitement des déchets ménagers relèvent des \cle{communes} (compétence propre transférée par la Loi n°2019/024, Code général des CTD, art. 231).
\item \textbf{Acteurs} : les communes (maîtres d'ouvrage, signataires des contrats), la société HYSACAM (Hygiène et Salubrité du Cameroun, opérateur privé délégataire), les riverains (usagers, payeurs de redevances), l'État (régulateur, financeur d'appui via la FEICOM et dotations).
\item \textbf{Financement} : redevances d'enlèvement des ordures ménagères (REOM) perçues par les communes, subventions d'équilibre de l'État, dotations de la FEICOM (Fonds d'équipement et d'intervention intercommunal).
\item \textbf{Dysfonctionnements identifiés} : retards de paiement des communes à l'opérateur (trésorerie), insuffisance du parc de bennes et camions, incivisme (dépôts sauvages, non-tri), faible taux de recouvrement de la REOM, absence de centres de traitement/valorisation (décharges à ciel ouvert).
\end{enumerate}
\textbf{Interprétation} : l'État n'exécute plus directement (désengagement opérationnel), mais il conserve un rôle de \cle{régulation} (normes d'hygiène, cahiers des charges, contrôle des contrats de délégation de service public) et d'\cle{appui} (subventions, équipements via FEICOM, formation). C'est l'État-providence transformé : moins opérateur direct, plus garant de l'intérêt général et régulateur du service public.
\end{experience}

\begin{attention}[Piège d'analyse]
Ne concluez pas trop vite « c'est la faute de l'État » ou « c'est la faute des citoyens ». Un service public défaillant résulte presque toujours d'une \emph{chaîne} de responsabilités : financement structurellement insuffisant, gouvernance locale fragile (gestion, recouvrement), comportements individuels (incivisme), capacité technique de l'opérateur. La rigueur philosophique exige de distinguer les niveaux de causalité (structurels, conjoncturels, individuels) avant de juger.
\end{attention}

\courssub{TP 5 : Débat contradictoire — le mandat présidentiel unique}

\textbf{Motion.} \emph{« Faut-il revenir à un mandat présidentiel unique de 5 ou 7 ans non renouvelable ? »}

\begin{methode}[Conduire un débat contradictoire]
\begin{enumerate}
\item Former deux camps (pour / contre) et un jury ; chaque camp prépare trois arguments hiérarchisés (le plus fort en premier).
\item Argumenter en mobilisant trois registres : \textbf{juridique} (la Constitution fixe la durée et la renouvelabilité du mandat — art. 6 ; sa révision suit une procédure stricte — art. 62), \textbf{politique} (alternance démocratique contre stabilité de l'action publique), \textbf{sociologique} (confiance des citoyens, fatigue du pouvoir long, légitimité).
\item Réfuter l'adversaire en attaquant ses prémisses ou son raisonnement, jamais sa personne (argument \emph{ad hominem} interdit).
\item Le jury évalue la solidité logique, la pertinence des exemples et la maîtrise du vocabulaire constitutionnel ; conclure par une synthèse nuancée.
\end{enumerate}
\end{methode}

\begin{exemplebox}[Arguments attendus]
\textbf{Pour le mandat unique} : il garantit l'\emph{alternance}, empêche la confiscation du pouvoir et la personnalisation de l'État, rapproche l'élu du peuple (la réélection n'étant plus un enjeu, le président gouverne pour tous et non pour sa réélection). \textbf{Contre} : un mandat unique trop court (5 ans) fragilise la \emph{continuité} des politiques publiques de long terme (grands travaux, SND30, éducation) et peut rendre le président « \emph{lame duck} » (canard boiteux) dès son investiture, sans levier sur sa majorité ; la stabilité institutionnelle est aussi une valeur démocratique (ex. mandat unique de 7 ans au Mexique : bilan mitigé). \textbf{Synthèse possible} : la vraie question n'est pas seulement la durée mais l'effectivité des \emph{contre-pouvoirs} (justice indépendante, presse libre, reddition des comptes, Cour des Comptes forte) qui contraignent le pouvoir, quel que soit le nombre de mandats.
\end{exemplebox}

\courssub{Dossier documentaire guidé : la loi n° 2019/024 du 24 décembre 2019}

Ce texte porte \cle{Code général des Collectivités Territoriales Décentralisées} (CTD). Il organise ce que l'État confie aux régions et aux communes.

\begin{methode}[Fiche de lecture structurée d'un texte juridique]
\begin{enumerate}
\item \textbf{Identification} : titre exact, date de promulgation, auteur (législateur : Parlement), objet du texte (organisation, fonctionnement, tutelle des CTD).
\item \textbf{Thèse centrale} : les CTD disposent de compétences \emph{propres} (exercées sous leur seule responsabilité), \emph{partagées} (exercées conjointement avec l'État) et \emph{transférées} (déléguées par l'État avec les ressources correspondantes — principe de concomitance).
\item \textbf{Concepts clés} : décentralisation (autonomie administrative et financière), déconcentration (représentants de l'État en région : gouverneurs, préfets), tutelle administrative (contrôle de légalité a posteriori), dotation (Fonds d'équipement, dotation générale de décentralisation).
\item \textbf{Question critique} : les ressources transférées (dotations, fiscalité locale) sont-elles à la hauteur des compétences transférées ? (Lien direct avec le TP 3 : ratio investissement/budget, capacité d'autofinancement des communes).
\item \textbf{Lien philosophique} : ce partage illustre la tension permanente entre unité de l'État (souveraineté une, indivisible) et autonomie locale (proximité, efficacité, démocratie participative). La souveraineté reste une, mais son exercice se pluralise.
\end{enumerate}
\end{methode}

\begin{aretenir}[L'essentiel du TP intégré]
\begin{itemize}
\item L'organigramme institutionnel révèle l'architecture des pouvoirs et des contre-pouvoirs : Exécutif (Président, PM, Gouvernement), Législatif (Parlement bicaméral), Juridictionnel (Cour suprême), Constitutionnel (Conseil constitutionnel), Contrôle (CONSUPE, CONAC, Cour des Comptes).
\item Le croquis de décentralisation montre concrètement comment l'État se rapproche du citoyen : 10 régions, 374 communes, des transferts de compétences (éducation, santé, voirie, environnement) et de ressources (dotations, fiscalité).
\item La lecture du budget (ratio investissement/fonctionnement) permet de juger l'État sur ses actes : environ 23\,\% en 2023, sous l'objectif SND30 de 30\,\% — écart entre discours d'émergence et réalité budgétaire.
\item L'analyse d'un service public (ordures ménagères) distingue compétences communales, acteurs privés (HYSACAM) et rôle régulateur/appui de l'État (FEICOM, normes).
\item Le débat sur le mandat présidentiel apprend à argumenter juridiquement (Constitution), politiquement (alternance/stabilité) et sociologiquement (confiance) — compétence reine de la dissertation et du commentaire.
\end{itemize}
\end{aretenir}

\begin{exercice}[Pour s'entraîner]
\begin{enumerate}
\item Complétez l'organigramme du TP 1 en ajoutant le Conseil économique et social (CES) et la Commission des droits de l'homme (CDHC), en précisant leurs rattachements.
\item Réalisez le croquis de votre région d'origine : chef-lieu, trois communes, deux réalisations locales visibles (marché, route, école, centre de santé), une flèche de transfert de l'État. N'oubliez ni titre (Quoi ? Où ? Quand ? Source), ni légende hiérarchisée, ni nord, ni échelle indicative.
\item À partir du Budget Citoyen de l'année en cours (disponible sur minfi.gov.cm), calculez le ratio investissement/budget total et comparez-le à la série 2019--2023. La tendance est-elle à la consolidation ou à l'infléchissement ?
\end{enumerate}
\end{exercice}


\coursec{Dissertation type Bac : Méthodologie et Sujets probables}

Après le TP intégré consacré à la cartographie institutionnelle et à l'analyse de l'action publique locale, nous disposons désormais d'une connaissance concrète du fonctionnement de l'État camerounais. Il reste à transformer ce savoir en \cle{compétence d'examen} : au Probatoire comme au Baccalauréat technique, l'épreuve de philosophie exige de rédiger, en quatre heures, une dissertation structurée, problématisée et argumentée. Cette dernière section fournit la méthode complète, quatre sujets types corrigés dans leur démarche, et les outils d'auto-évaluation.

\courssub{La structure type d'une dissertation philosophique}

\begin{definition}[La dissertation philosophique]
La \cle{dissertation} est un exercice de réflexion écrite qui, à partir d'une question, construit une réponse personnelle et argumentée en mobilisant des concepts, des thèses d'auteurs et des exemples. Elle n'est ni un résumé de cours, ni une opinion : c'est une \emph{démonstration} conduite pas à pas.
\end{definition}

\begin{propriete}[Architecture et proportions de la copie]
Une copie équilibrée respecte les proportions suivantes :
\begin{itemize}
\item \textbf{Introduction} : 15 à 20 lignes (accroche, définitions des termes, problématique, annonce du plan).
\item \textbf{Développement} : environ 80\,\% de la copie, en 2 ou 3 parties équilibrées. Chaque paragraphe = 1 idée directrice + arguments + exemples précis (Cameroun, monde, auteurs).
\item \textbf{Conclusion} : 10 à 15 lignes (bilan répondant à la problématique + ouverture).
\end{itemize}
\end{propriete}

\begin{popfigure}
\begin{center}
\begin{tikzpicture}[font=\small]
% Entonnoir introduction
\draw[fill=popblueL,draw=popblue] (0,6.4) -- (6,6.4) -- (4.6,5.2) -- (1.4,5.2) -- cycle;
\node at (3,5.85) {\textbf{Introduction : du général au sujet}};
\node[align=center] at (3,4.85) {\footnotesize Accroche $\rightarrow$ Définitions $\rightarrow$ Problématique $\rightarrow$ Plan};
% Colonnes développement
\draw[fill=poporangeL,draw=poporange] (0.2,4.3) rectangle (2.9,1.6);
\node[align=center] at (1.55,2.95) {\textbf{Partie 1}\\ \footnotesize Thèse (l'évidence)\\ \footnotesize Argument + auteur\\ \footnotesize Exemple};
\draw[fill=popgreenL,draw=popgreen] (3.1,4.3) rectangle (5.8,1.6);
\node[align=center] at (4.45,2.95) {\textbf{Partie 2}\\ \footnotesize Antithèse (la limite)\\ \footnotesize Argument + auteur\\ \footnotesize Exemple};
\draw[fill=poppurpleL,draw=poppurple] (6.0,4.3) rectangle (8.7,1.6);
\node[align=center] at (7.35,2.95) {\textbf{Partie 3}\\ \footnotesize Synthèse (dépassement)\\ \footnotesize Argument + auteur\\ \footnotesize Exemple};
% Entonnoir inversé conclusion
\draw[fill=popgoldL,draw=popgold] (1.4,1.1) -- (4.6,1.1) -- (6,-0.1) -- (0,-0.1) -- cycle;
\node at (3,0.5) {\textbf{Conclusion : bilan + ouverture}};
\end{tikzpicture}
\end{center}
\legende{Architecture type d'une dissertation philosophique : l'introduction resserre le propos vers le sujet (entonnoir), le développement confronte thèse, antithèse et synthèse en colonnes équilibrées, la conclusion rouvre la réflexion (entonnoir inversé).}
\end{popfigure}

\begin{methode}[Problématiser : transformer le sujet en tension]
\begin{enumerate}
\item Repérer la \textbf{thèse évidente} contenue dans le sujet. Exemple : « L'État assure la sécurité » (c'est le sens commun, Hobbes le confirme).
\item Formuler l'\textbf{antithèse} qui la conteste ou la dépasse : « Mais l'État doit-il aussi assurer le bonheur, la justice sociale, l'éducation ? »
\item Construire la \textbf{problématique} comme question ouverte sur cette tension : « Jusqu'où l'État doit-il intervenir sans aliéner la liberté des citoyens ? »
\item Vérifier que la problématique reprend \emph{tous} les termes du sujet et annonce le plan en deux ou trois temps.
\end{enumerate}
\end{methode}

\begin{attention}[Le piège du hors-sujet disciplinaire]
Erreur fréquente : faire un cours de \textbf{droit constitutionnel} ou d'\textbf{histoire} au lieu de philosophie. Réciter les articles de la Constitution de 1996 ou la chronologie des indépendances ne rapporte aucun point de réflexion. Le droit et l'histoire fournissent des \cle{exemples} ; la philosophie fournit des \cle{concepts} (légitimité, contrat, justice) et des \cle{arguments} logiques. Chaque fait cité doit servir une idée, jamais s'y substituer.
\end{attention}

\courssub{Techniques de rédaction : définir, argumenter, lier}

\textbf{Définir les termes du sujet.} Tout mot du sujet doit être défini dans l'introduction selon trois registres : l'\emph{étymologie} (\emph{État} vient du latin \emph{status}, « manière d'être, situation stable »), le \emph{sens philosophique} (l'État comme forme d'organisation du pouvoir politique fondée sur la souveraineté), le \emph{sens juridique} (personne morale de droit public regroupant un peuple, un territoire, un gouvernement). Cette triple définition montre la maîtrise du concept.

\textbf{Argumenter.} Chaque paragraphe suit la chaîne : \textbf{Thèse $\rightarrow$ Argument $\rightarrow$ Exemple ou référence d'auteur $\rightarrow$ Lien} (phrase de transition vers le paragraphe suivant). Sans lien, la copie devient une juxtaposition ; avec le lien, elle devient un raisonnement.

\begin{propriete}[Boîte à outils linguistique]
\textbf{Mots-clés de liaison} : \emph{en outre, par conséquent, cependant, certes... mais, néanmoins, en définitive, dès lors, or}. \textbf{Verbes d'analyse} : \emph{illustre, démontre, contredit, nuance, fonde, légitime, remet en cause, éclaire}. Éviter « je pense que » ; préférer « il apparaît que », « on peut soutenir que ».
\end{propriete}

\courssub{Quatre sujets types analysés}

\begin{exoresolu}[Sujet type 1 : « L'État a-t-il pour seule fonction d'assurer la sécurité ? »]
Énoncé : construire la problématique et le plan détaillé de ce sujet.
\tcblower
\textbf{Définitions} : État (souveraineté, territoire, population) ; fonction (tâche assignée à une institution) ; sécurité (protection des personnes et des biens).
\textbf{Problématique} : si la sécurité est la fonction première de l'État (Hobbes), est-elle sa \emph{seule} fonction ? L'État peut-il se réduire à un gendarme, ou doit-il aussi promouvoir la justice et le bien-être ?
\textbf{Plan en trois parties} :
\begin{itemize}
\item \emph{Partie 1 — Oui, la sécurité est la fonction essentielle} : Hobbes, \emph{Léviathan} : hors de l'État, « la vie de l'homme est solitaire, misérable, dangereuse, animale et brève » ; l'État naît de la peur de la mort violente. Exemple : les fonctions régaliennes (police, armée, justice) au Cameroun.
\item \emph{Partie 2 — Non, l'État a aussi des fonctions sociales} : Locke (protection de la propriété et des libertés), Rousseau (l'État vise le bien commun), Rawls (justice comme équité). Exemple : la SND30 (Stratégie Nationale de Développement 2020-2030) assigne à l'État camerounais des objectifs d'éducation, de santé, d'infrastructures.
\item \emph{Partie 3 — Synthèse} : la sécurité est la \emph{condition} des autres fonctions, non leur substitut ; un État qui ne fait que sécuriser devient policier, un État qui néglige la sécurité devient impuissant.
\end{itemize}
\end{exoresolu}

\begin{exemplebox}[Sujet type 2 : « La légitimité du pouvoir repose-t-elle uniquement sur le consentement des gouvernés ? »]
\textbf{Problématique} : le consentement (élections, contrat social) suffit-il à fonder l'obéissance, ou d'autres sources de légitimité existent-elles ?
\textbf{Axes} : (1) Weber et les trois légitimités (traditionnelle, charismatique, rationnelle-légale) : le consentement n'est qu'une source parmi d'autres ; (2) la théorie du contrat social (Rousseau : la volonté générale) fait du consentement le fondement moderne par excellence — exemple : les élections comme procédure de légitimation ; (3) limites : crises de légitimité (abstention, contestation), droit de résistance à l'oppression (Locke, Déclaration de 1789, art. 2). Conclusion : le consentement est nécessaire dans une démocratie moderne, mais il doit être éclairé, renouvelé et accompagné de l'État de droit.
\end{exemplebox}

\begin{exemplebox}[Sujet type 3 : « Le régime semi-présidentiel est-il le meilleur compromis pour la stabilité démocratique ? »]
\textbf{Problématique} : combiner un président élu et un gouvernement responsable devant le parlement garantit-il l'équilibre des pouvoirs ?
\textbf{Axes} : (1) la typologie de Duverger : le semi-présidentialisme combine les avantages du présidentialisme (stabilité de l'exécutif) et du parlementarisme (responsabilité du gouvernement) — cas de la France de la V\textsuperscript{e} République ; (2) le cas camerounais : Constitution de 1996, président élu au suffrage universel, Premier ministre responsable devant l'Assemblée nationale ; (3) les risques : dérive présidentialiste (hyperconcentration du pouvoir) ou instabilité parlementaire (cohabitation, renversements de majorité). Conclusion nuancée : le meilleur régime dépend moins du texte que de la culture démocratique et du respect effectif de la séparation des pouvoirs (Montesquieu).
\end{exemplebox}

\begin{exemplebox}[Sujet type 4 : « Justice sociale et inégalités : l'État peut-il tout corriger ? »]
\textbf{Problématique} : l'État a-t-il le devoir — et les moyens — de supprimer les inégalités ?
\textbf{Axes} : (1) Rawls : les inégalités ne sont justes que si elles bénéficient aux plus défavorisés (principe de différence) ; Sen : la justice exige le développement des \emph{capacités} (santé, éducation) ; (2) les politiques publiques comme instruments : impôt progressif, éducation gratuite, subventions — exemples camerounais (bourses, hôpitaux de district) ; (3) les limites : contraintes budgétaires, corruption qui détourne la redistribution, rôle irremplaçable de la société civile et des ONG. Conclusion : l'État ne peut tout corriger seul, mais il demeure le principal garant de la justice sociale.
\end{exemplebox}

\begin{popfigure}
\begin{center}
\begin{tikzpicture}[mindmap, grow cyclic, every node/.style={concept, font=\scriptsize, align=center},
root concept/.append style={concept color=popblue, font=\small\bfseries},
level 1/.append style={level distance=3.4cm, sibling angle=60},
level 2/.append style={level distance=2.1cm, sibling angle=45}]
\node[root concept, text=white, align=center]{Bac Philo Tech\\ État et pouvoir}
  child[concept color=poporange] { node {Contractualistes}
    child { node[align=center]{Hobbes\\ sécurité,\\ Léviathan} }
    child { node[align=center]{Locke\\ propriété,\\ résistance} }
    child { node[align=center]{Rousseau\\ volonté\\ générale} } }
  child[concept color=popgreen] { node[align=center]{Sociologie\\ du pouvoir}
    child { node[align=center]{Weber\\ 3 légitimités} }
    child { node[align=center]{Marx\\ État et\\ domination} } }
  child[concept color=poppurple] { node {Justice}
    child { node[align=center]{Rawls\\ voile\\ d'ignorance} }
    child { node[align=center]{Sen\\ capacités} } }
  child[concept color=popgold] { node {Droit}
    child { node[align=center]{Kelsen\\ hiérarchie\\ des normes} }
    child { node[align=center]{Montesquieu\\ séparation\\ des pouvoirs} } }
  child[concept color=poppink] { node[align=center]{Pensée\\ africaine}
    child { node[align=center]{Mbembe\\ postcolonie} }
    child { node[align=center]{Wiredu\\ consensus} }
    child { node[align=center]{Mudimbe\\ invention\\ de l'Afrique} } }
  child[concept color=popblueL] { node[align=center]{Liberté\\ politique}
    child { node[align=center]{Arendt\\ pouvoir et\\ violence} } };
\end{tikzpicture}
\end{center}
\legende{Carte mentale du répertoire d'auteurs à mobiliser au Bac : chaque branche associe un thème du programme à ses références majeures, y compris les auteurs africains valorisés dans les copies.}
\end{popfigure}

\courssub{Entraînement chronométré et évaluation}

\begin{methode}[Gérer les 4 heures de l'épreuve]
\begin{enumerate}
\item \textbf{1\textsuperscript{re} heure} : analyse du sujet (définitions, présupposés), recherche d'idées au brouillon, construction de la problématique et du plan détaillé.
\item \textbf{2 heures} : rédaction de l'introduction, du développement et de la conclusion au propre.
\item \textbf{Dernière heure} : relecture (orthographe, cohérence, transitions), correction croisée avec un camarade à l'aide de la grille d'auto-évaluation.
\end{enumerate}
\end{methode}

\begin{exemplebox}[Barème indicatif du Bac Philo Tech (coefficient 2 à 3 selon les séries)]
\begin{center}
\begin{tabularx}{\linewidth}{|c|X|}\hline
\rowcolor{popblueL} \textbf{Note /20} & \textbf{Appréciation} \\ \hline
0--5 & Hors sujet ou copie insuffisante \\ \hline
10--12 & Passable : connaissances présentes mais pas de problématique \\ \hline
14--16 & Bien : problématique, structure claire, exemples pertinents \\ \hline
18--20 & Excellent : originalité, style, pertinence des références \\ \hline
\end{tabularx}
\end{center}
Grille détaillée sur 26 points convertie sur 20 : compréhension du sujet (4), problématique (4), plan (4), connaissances (4), argumentation (4), expression (4), conclusion et ouverture (2), présentation (4).
\end{exemplebox}

\begin{popfigure}
\begin{center}
\begin{tikzpicture}[font=\scriptsize]
\draw[fill=popblueL] (0,0) rectangle (12.6,0.7);
\node[anchor=west] at (0.1,0.35) {\textbf{Critère}};
\node[anchor=west] at (3.6,0.35) {\textbf{Indicateur de réussite}};
\node[anchor=west] at (9.4,0.35) {\textbf{Auto-note}};
\node[anchor=west] at (11.2,0.35) {\textbf{Note prof}};
\foreach \y/\a/\b in {
 -0.7/{Compréhension}/{Tous les termes du sujet sont définis},
 -1.4/{Problématique}/{Une vraie question de tension est posée},
 -2.1/{Plan}/{2 ou 3 parties équilibrées et annoncées},
 -2.8/{Connaissances}/{Auteurs cités avec exactitude},
 -3.5/{Argumentation}/{Thèse + argument + exemple + lien},
 -4.2/{Expression}/{Connecteurs logiques, orthographe soignée},
 -4.9/{Conclusion}/{Réponse claire + ouverture pertinente}}{
 \draw[fill=white] (0,\y) rectangle (12.6,\y+0.7);
 \node[anchor=west] at (0.1,\y+0.35) {\a};
 \node[anchor=west] at (3.6,\y+0.35) {\b};
 \draw (9.4,\y) rectangle (11.0,\y+0.7);
 \draw (11.2,\y) rectangle (12.6,\y+0.7);}
\end{tikzpicture}
\end{center}
\legende{Grille d'auto-évaluation de la dissertation : à compléter lors de la correction croisée pour repérer les points faibles avant l'examen.}
\end{popfigure}

\begin{savaistu}[Dissertation ou explication de texte ?]
Au Baccalauréat technique, le sujet de dissertation porte le plus souvent sur les notions du programme de l'année : l'État, la justice, le travail, la technique, la religion. L'explication de texte, elle, porte sur un extrait d'un auteur au programme — par exemple un passage du \emph{Contrat social} de Rousseau, du \emph{Léviathan} de Hobbes ou d'\emph{Économie et Société} de Weber. Réviser les thèses centrales de ces auteurs, c'est se préparer aux deux épreuves à la fois.
\end{savaistu}

\begin{aretenir}[L'essentiel de la méthode]
Une dissertation réussie repose sur quatre piliers : des \textbf{définitions} précises des termes du sujet, une \textbf{problématique} construite comme tension entre deux thèses, un \textbf{plan} équilibré où chaque paragraphe enchaîne thèse, argument, exemple et lien, et une \textbf{conclusion} qui répond à la question posée avant d'ouvrir sur un élargissement. Les exemples camerounais (SND30, Constitution de 1996, action publique locale étudiée au TP) donnent à la copie son ancrage et son originalité.
\end{aretenir}

\begin{exercice}[Entraînement]
En une heure, rédiger l'introduction complète (accroche, définitions, problématique, annonce du plan) du sujet : « L'État peut-il tout corriger ? », puis échanger sa copie avec un camarade et la noter à l'aide de la grille d'auto-évaluation.
\end{exercice}

\newpage

\coursec{Activité d'intégration}

\coursec{Leçon 19 : L’État et le pouvoir}

\courssub{TP Intégré : Cartographie institutionnelle et Analyse de l'action publique locale}

\courssub{Objectifs de l'activité}

Cette activité d'intégration prend la forme d'un \cle{débat citoyen simulé} intitulé : « Réforme de la Constitution : touche-t-on à l'intouchable ? ». Elle vise à mobiliser l'ensemble des savoirs de la leçon sur l'État et le pouvoir dans une situation concrète de la vie politique camerounaise. À l'issue de l'activité, l'élève doit être capable de :

\begin{itemize}
\item définir l'État et identifier ses éléments constitutifs (territoire, population, souveraineté) dans un cas réel ;
\item distinguer les \cle{fondements} du pouvoir politique (force, droit divin, contrat social) et évaluer la \cle{légitimité} d'une décision politique ;
\item appliquer la théorie de la \cle{séparation des pouvoirs} (Montesquieu) au fonctionnement des institutions camerounaises ;
\item conduire une \cle{argumentation} philosophique structurée (thèse, arguments, objections, réponses) ;
\item rédiger et justifier une conclusion personnelle sous forme d'éditorial argumenté.
\end{itemize}

\courssub{Situation}

\textbf{Contexte.} À Yaoundé, en mars 2026, un groupe de députés dépose à l'Assemblée nationale une \emph{proposition de loi de révision constitutionnelle} portant sur l'article 6 de la Constitution du 18 janvier 1996 (révisée par la loi n°2008/001 du 14 avril 2008), qui fixe les modalités de l'élection du Président de la République : durée du mandat, limitation ou non du nombre de mandats, âge d'éligibilité, mode de scrutin. La rumeur de cette réforme enflamme le pays : meetings politiques à Douala et à Bafoussam, pétitions en ligne ayant recueilli \num{85000} signatures en trois semaines, communiqués contradictoires des partis, éditoriaux dans la presse. Les uns dénoncent une « confiscation du pouvoir », les autres invoquent la « stabilité indispensable au développement ».

\textbf{Données du dossier.} La classe de Terminale technique est transformée en \emph{Commission nationale de réflexion}. Elle est divisée en quatre groupes de 8 à 10 élèves :

\begin{enumerate}
\item \textbf{Commission 1 — Majorité présidentielle} : elle défend la réforme au nom de la stabilité des institutions, de l'expérience du chef de l'État et de la continuité des grands chantiers (routes, barrages, industrialisation).
\item \textbf{Commission 2 — Opposition radicale} : elle combat la réforme au nom de l'alternance démocratique, de la souveraineté du peuple et de l'exigence de renouvellement du personnel politique.
\item \textbf{Commission 3 — Société civile et ONG} : elle défend la paix sociale, le consensus national et la préservation de la cohésion dans un pays de plus de 250 groupes ethniques.
\item \textbf{Commission 4 — Cour constitutionnelle} : elle veille à la légalité de la procédure de révision et rappelle que la forme républicaine de l'État ne peut faire l'objet d'aucune révision (article 63 de la Constitution).
\end{enumerate}

\textbf{Problématique du débat.} Une révision de l'article 6 est-elle un exercice légitime de la souveraineté nationale, ou une atteinte au contrat social qui fonde l'État de droit ? Autrement dit : \emph{le pouvoir peut-il modifier les règles mêmes qui le fondent sans se délégitimer ?}

\courssub{Consignes}

\begin{enumerate}
\item \textbf{Préparation en commission (30 min).} À partir du dossier documentaire remis (extraits de la Constitution de 1996, extraits du \emph{Léviathan} de Hobbes et du \emph{Contrat social} de Rousseau, communiqués de partis, avis de juristes), rédigez une fiche d'argumentation comportant : votre thèse, trois arguments étayés par un auteur ou un texte, deux objections adverses anticipées et vos réponses.
\item \textbf{Question de cadrage 1.} Rappelez la définition de l'État et ses trois éléments constitutifs. La réforme envisagée menace-t-elle l'un de ces éléments ? Justifiez.
\item \textbf{Question de cadrage 2.} Selon la théorie du contrat social (Hobbes, Locke, Rousseau), d'où le pouvoir tire-t-il sa légitimité ? Une révision constitutionnelle votée par le Parlement respecte-t-elle ce fondement ? Discutez selon votre commission.
\item \textbf{Question de cadrage 3.} En vous appuyant sur Montesquieu (\emph{De l'esprit des lois}), montrez comment la séparation des pouvoirs permet de contrôler une telle réforme. Quel rôle précis joue la Cour constitutionnelle dans ce débat ?
\item \textbf{Débat public (40 min).} Chaque commission dispose de 5 minutes de présentation, puis le débat est ouvert sous la modération du professeur. Chaque prise de parole doit citer au moins un auteur ou un article du dossier.
\item \textbf{Vote final.} À main levée, la classe se prononce : la réforme renforce-t-elle ou affaiblit-elle l'État de droit au Cameroun ?
\item \textbf{Devoir maison.} Rédigez individuellement un éditorial d'environ \num{200} mots intitulé : « Pourquoi cette réforme renforce (ou affaiblit) l'État de droit au Cameroun ». Structure exigée : accroche, thèse, deux arguments, une objection traitée, conclusion.
\end{enumerate}

\courssub{Ressources à mobiliser}

\begin{aretenir}[Notions utiles pour le débat]
\begin{itemize}
\item \textbf{L'État} : organisation politique souveraine réunissant une \cle{population} sur un \cle{territoire} délimité, dotée d'un \cle{pouvoir} institutionnalisé.
\item \textbf{Fondements du pouvoir} : la force (Thrasymachos), le droit divin (Bossuet), le \cle{contrat social} (Hobbes : sécurité ; Locke : protection des droits naturels ; Rousseau : volonté générale).
\item \textbf{Légitimité} : un pouvoir est légitime s'il repose sur le consentement des gouvernés et s'exerce dans le cadre du droit.
\item \textbf{Séparation des pouvoirs} (Montesquieu) : « Pour qu'on ne puisse abuser du pouvoir, il faut que, par la disposition des choses, le pouvoir arrête le pouvoir. »
\item \textbf{État de droit} : État soumis à sa propre loi ; la Constitution est la norme suprême.
\item \textbf{Formes de gouvernement} : monarchie, aristocratie, démocratie (Aristote) ; régimes présidentiel, parlementaire, mixte.
\item \textbf{Finalité de l'État} : ordre et sécurité (Hobbes), justice et liberté (Locke, Rousseau), bien-être social (État-providence).
\end{itemize}
\end{aretenir}

\begin{attention}[Pièges à éviter]
\begin{itemize}
\item Confondre \cle{légalité} (conformité à la procédure écrite) et \cle{légitimité} (acceptation morale et fondement du pouvoir).
\item Réduire l'État à son gouvernement : l'État est permanent, le gouvernement est temporaire.
\item Oublier que la souveraineté nationale appartient au peuple (art. 2 Constitution), qui l'exerce par référendum ou par l'intermédiaire de ses représentants.
\item Négliger le rôle du juge constitutionnel (Cour constitutionnelle) comme gardien de la hiérarchie des normes et de l'intangibilité de la forme républicaine (art. 63).
\end{itemize}
\end{attention}

\courssub{Démarche et corrigé commenté}

\begin{exoresolu}[Corrigé commenté du débat et de l'éditorial]
\textbf{Énoncé :} À partir de la situation de la révision de l'article 6, analyser la distinction légalité/légitimité, appliquer les théories du contrat social et de la séparation des pouvoirs, et rédiger un éditorial argumenté.
\tcblower
\textbf{Étape 1 — Cadrage conceptuel (questions 1 et 2).} L'État camerounais repose sur un territoire (\SI{475442}{km^2}), une population (plus de \num{28000000} d'habitants) et une souveraineté exercée par des institutions. La réforme de l'article 6 ne supprime aucun de ces éléments : elle touche au \emph{mode d'exercice} du pouvoir, non à l'existence de l'État. En revanche, elle engage le \emph{fondement} du pouvoir : si, selon Rousseau, « la souveraineté ne peut être représentée », le peuple reste le détenteur originaire du pouvoir. Une révision votée par les seuls députés est juridiquement valable si la procédure de l'article 63 est respectée (vote des 3/5 du Parlement réuni en Congrès ou référendum), mais sa légitimité démocratique peut être contestée si elle n'est pas soumise au référendum. C'est la distinction capitale entre \cle{légalité} et \cle{légitimité}.

\textbf{Étape 2 — Arguments attendus de chaque commission.}
\begin{itemize}
\item \emph{Majorité} : Hobbes montre que la première finalité de l'État est la sécurité et l'ordre (\emph{Léviathan}) ; la stabilité institutionnelle garantit la continuité des politiques publiques (grands travaux, lutte contre l'insécurité). La légalité de la procédure parlementaire suffit : les députés représentent la nation (mandat représentatif).
\item \emph{Opposition} : Locke soutient que le pouvoir est un dépôt (\emph{trust}) révocable (\emph{Second Traité du gouvernement civil}) ; sans alternance, le consentement devient fictif. Rousseau ajoute que la volonté générale exige la participation du peuple : la réforme doit passer par le référendum. La limitation des mandats est un rempart contre la confiscation du pouvoir (dérive oligarchique).
\item \emph{Société civile} : la finalité de l'État est la paix sociale et le bien commun ; une réforme imposée sans consensus risque la radicalisation et la violence. Elle propose une médiation et un large dialogue national (conférence nationale) avant tout vote.
\item \emph{Cour constitutionnelle} : elle rappelle la hiérarchie des normes (Kelsen) et l'article 63 (intangibilité de la forme républicaine, de l'unité de l'État, de la laïcité). Elle illustre la thèse de Montesquieu : le pouvoir judiciaire « arrête le pouvoir » en censurant toute révision contraire à la Constitution (contrôle de constitutionnalité a priori).
\end{itemize}

\textbf{Étape 3 — Traitement de l'objection centrale.} Objection de la majorité : « La révision est légale, donc légitime. » Réponse : la légalité est une condition nécessaire mais non suffisante de la légitimité ; un pouvoir peut respecter les formes tout en vidant le contrat social de sa substance, ce qu'Aristote appellerait une dérive de la démocratie en démagogie (gouvernement de la multitude pauvre pour son intérêt) ou en oligarchie (gouvernement des riches pour leur intérêt).

\textbf{Étape 4 — Modèle d'éditorial (extrait, environ \num{200} mots).} « Touche-t-on à l'intouchable ? En réformant l'article 6, nos élus affirment agir au nom de la stabilité. Mais la stabilité n'est pas une fin en soi : Hobbes lui-même subordonnait l'ordre à la protection des citoyens. Or un pouvoir qui modifie seul les règles de sa propre succession rompt le pacte qui le fonde. Selon Rousseau, la souveraineté appartient au peuple : toute révision de cette ampleur exige donc le référendum. On objectera que les députés représentent la nation ; mais la représentation n'est pas une délégation absolue, elle suppose la confiance, que seul le vote populaire peut renouveler. En exigeant le contrôle de la Cour constitutionnelle et la consultation du peuple, nous ne fragilisons pas l'État : nous le ramenons à sa finalité véritable, la justice. Une réforme acceptée renforce l'État de droit ; une réforme imposée le dissout. »

\textbf{Étape 5 — Grille d'évaluation (20 pts).} Exactitude des notions définies (5 pts) ; Qualité de l'argumentation et usage pertinent des auteurs/textes (5 pts) ; Prise de parole, écoute et respect du contradicteur (5 pts) ; Éditorial : structure, cohérence, style (5 pts).
\end{exoresolu}

\courssub{Bilan de l'activité}

Ce débat simulé a permis de mettre en évidence trois acquis essentiels. D'abord, la distinction décisive entre \cle{légalité} et \cle{légitimité} : une décision peut être conforme à la procédure tout en étant contestée dans son fondement, ce qui renvoie à la théorie du contrat social. Ensuite, l'utilité pratique de la \cle{séparation des pouvoirs} : la Cour constitutionnelle incarne concrètement le principe de Montesquieu selon lequel « le pouvoir arrête le pouvoir ». Enfin, l'activité a montré que les notions philosophiques (souveraineté, contrat social, État de droit) ne sont pas des abstractions : elles éclairent directement les crises politiques réelles du Cameroun et outillent l'élève pour la dissertation du Baccalauréat technique, où il devra, comme ici, problématiser, argumenter et conclure de manière personnelle et justifiée.

\courssub{Dissertation type Bac : Méthodologie et Sujets probables}

\begin{methode}[Rédiger une dissertation philosophique au Bac Technique]
\begin{enumerate}
\item \textbf{Analyser le sujet} : Définir les termes, identifier la tension (opposition, paradoxe), formuler la \cle{problématique} (question philosophique précise).
\item \textbf{Construire le plan} (dialectique recommandé) :
      \begin{itemize}
      \item \emph{Thèse} : Réponse affirmative immédiate (arguments + exemples/auteurs).
      \item \emph{Antithèse} : Réponse négative / limites de la thèse (arguments + exemples/auteurs).
      \item \emph{Synthèse} : Dépassement, nuance, condition de validité, ouverture.
      \end{itemize}
\item \textbf{Rédiger} :
      \item \emph{Introduction} : Accroche $\rightarrow$ Définitions $\rightarrow$ Problématique $\rightarrow$ Annonce du plan.
      \item \emph{Développement} : Un argument par paragraphe (Idée directrice $\rightarrow$ Explication $\rightarrow$ Exemple/Autorité $\rightarrow$ Lien).
      \item \emph{Conclusion} : Bilan $\rightarrow$ Réponse à la problématique $\rightarrow$ Ouverture.
\item \textbf{Relire} : Orthographe, syntaxe, cohérence logique, respect du programme.
\end{enumerate}
\end{methode}

\begin{exercice}[Sujets probables pour l'entraînement]
\begin{enumerate}
\item L'État a-t-il pour seule finalité d'assurer la sécurité ?
\item La séparation des pouvoirs est-elle une condition suffisante de la liberté politique ?
\item Obéir aux lois, est-ce renoncer à sa liberté ?
\item La démocratie se réduit-elle au vote majoritaire ?
\item Une réforme constitutionnelle peut-elle être illégitime tout en étant légale ?
\item L'État-providence est-il une nécessité ou une utopie ?
\end{enumerate}
\end{exercice}

\begin{corrige}[Exercice n°5 : Une réforme constitutionnelle peut-elle être illégitime tout en étant légale ?]
\textbf{Problématique :} La conformité à la procédure juridique (légalité) garantit-elle l'accord moral et le fondement du pouvoir (légitimité) ?
\textbf{Plan :}
\begin{itemize}
\item \textbf{I. La légalité : condition nécessaire de la légitimité.} Dans l'État de droit, le pouvoir ne s'exerce que par la loi (Kelsen, Constitution). La révision de l'art. 6 suit l'art. 63 (procédure parlementaire ou référendum). Respect des formes = légalité.
\item \textbf{II. La légalité ne suffit pas à fonder la légitimité.} Fondement du pouvoir = consentement (Contrat social : Hobbes, Locke, Rousseau). Légitimité = adhésion du peuple. Risque de « légalité sans légitimité » : révision adoptée par une majorité parlementaire issue d'un scrutin contesté, sans consultation populaire (référendum), modifiant les règles du jeu démocratique (mandats, éligibilité). Dérive oligarchique (Aristote).
\item \textbf{III. L'articulation : le contrôle de constitutionnalité et le référendum comme garants.} La Cour constitutionnelle (Montesquieu) veille à la hiérarchie des normes (art. 63). Le référendum (souveraineté primaire, art. 2) restaure la légitimité directe. La réforme renforce l'État de droit si elle unit légalité formelle et légitimité populaire.
\end{itemize}
\end{corrige}

\newpage

\coursec{Méthodes et exercices résolus}

\courssub{Fiches méthodes et exercices résolus}

\begin{methode}[Analyser une situation concrète à la lumière des éléments constitutifs de l'État]
Pour appliquer la notion d'État à une situation de la vie courante (un conflit de voisinage, une élection, une crise sécuritaire), suivre une démarche rigoureuse :
\begin{enumerate}
\item \textbf{Identifier les éléments constitutifs} présents ou absents dans la situation : la \cle{population}, le \cle{territoire}, la \cle{souveraineté} (pouvoir de contrainte légitime).
\item \textbf{Repérer l'organe étatique concerné} : exécutif (préfet, maire, gouverneur), législatif (Assemblée nationale), judiciaire (tribunal).
\item \textbf{Qualifier le pouvoir en jeu} : est-il légitime (fondé sur le droit, la tradition ou le charisme selon Weber) ou seulement de fait (force brute) ?
\item \textbf{Formuler le problème philosophique} : la situation révèle-t-elle une défaillance de l'État, un abus de pouvoir, un conflit de légitimités ?
\item \textbf{Conclure} en mobilisant une thèse d'auteur (Hobbes, Rousseau, Weber) pour justifier la réponse.
\end{enumerate}
\textbf{Réflexe} : toujours nommer précisément l'institution camerounaise concernée (sous-préfecture, tribunal de première instance, conseil municipal) plutôt que de parler de « l'État » en général.
\end{methode}

\begin{exoresolu}[Le préfet et la manifestation interdite]
\textbf{Énoncé} : Dans une ville du Cameroun, une association annonce une marche pacifique pour protester contre la dégradation d'une route. L'autorité administrative interdit la manifestation au motif de « trouble à l'ordre public ». Les organisateurs estiment cette interdiction abusive. À partir des éléments constitutifs de l'État et de la notion de pouvoir légitime, analysez cette situation.
\tcblower
\textbf{Solution détaillée.}
\begin{enumerate}
\item \textbf{Éléments constitutifs mobilisés} : la situation se déroule sur un \cle{territoire} délimité (la ville), concerne une fraction de la \cle{population} (les habitants usagers de la route) et met en jeu la \cle{souveraineté} de l'État, c'est-à-dire son pouvoir de contrainte.
\item \textbf{Organe concerné} : l'autorité administrative (préfet ou sous-préfet) représente le pouvoir \cle{exécutif} déconcentré de l'État sur le terrain.
\item \textbf{Qualification du pouvoir} : l'interdiction est un acte de puissance publique. Elle est \emph{légale} si elle s'appuie sur la réglementation des manifestations (déclaration préalable, menace réelle pour l'ordre public). Mais la légalité ne suffit pas : selon Max Weber, le pouvoir légitime de type \cle{rationnel-légal} suppose que la règle soit appliquée de manière impartiale et non arbitraire. Si l'interdiction vise seulement à empêcher toute critique, elle devient un abus de pouvoir.
\item \textbf{Problème philosophique} : où s'arrête le droit de l'État de maintenir l'ordre et où commence la liberté d'expression des citoyens ? C'est le problème classique de la limite du pouvoir politique.
\item \textbf{Conclusion} : l'État, selon Hobbes, a pour fonction d'assurer la sécurité et l'ordre ; mais selon Rousseau, sa légitimité repose sur la volonté générale, donc sur le respect des libertés des citoyens. L'interdiction n'est donc légitime que si elle protège réellement l'ordre public et non si elle étouffe une revendication légitime. La voie de droit (recours devant le tribunal administratif) est le moyen prévu par l'État de droit pour trancher ce litige.
\end{enumerate}
\textbf{Phrase de conclusion rédigée} : Cette situation montre que le pouvoir de l'État n'est légitime que s'il concilie l'ordre public, sa fonction essentielle, avec le respect des libertés qu'il a précisément pour mission de garantir.
\end{exoresolu}

\begin{methode}[Distinguer et comparer les formes de gouvernement dans un commentaire ou une dissertation]
Face à un sujet portant sur les régimes politiques :
\begin{enumerate}
\item \textbf{Définir chaque forme} avec un critère précis : le nombre de gouvernants (monarchie, aristocratie, démocratie selon Aristote) ou le mode de séparation des pouvoirs (régime présidentiel, parlementaire, semi-présidentiel).
\item \textbf{Identifier le critère de distinction} demandé par le sujet : ne pas mélanger la classification d'Aristote (nombre) et la classification moderne (organisation des pouvoirs).
\item \textbf{Illustrer par des exemples} : monarchie constitutionnelle (Royaume-Uni), république présidentielle (États-Unis), république camerounaise (régime à dominante présidentielle).
\item \textbf{Comparer avec nuance} : montrer qu'aucun régime n'est parfait et que la démocratie elle-même connaît des dérives (tyrannie de la majorité, démagogie dénoncées par Platon).
\item \textbf{Conclure} en dégageant l'enjeu : la forme de gouvernement importe moins que le respect de l'État de droit et de la séparation des pouvoirs (Montesquieu).
\end{enumerate}
\textbf{Formule à retenir} : « Qui gouverne ? » (nombre) et « Comment gouverne-t-on ? » (organisation et limites du pouvoir) sont deux questions distinctes.
\end{methode}

\begin{exoresolu}[Démocratie et régime camerounais]
\textbf{Énoncé} : Un élève affirme : « Le Cameroun est une république, donc le pouvoir du président est limité. » Corrigez et précisez cette affirmation en distinguant les formes de gouvernement et en vous appuyant sur la Constitution camerounaise de 1996.
\tcblower
\textbf{Solution détaillée.}
\begin{enumerate}
\item \textbf{Définitions} : une \cle{république} est un régime où le chef de l'État est élu (directement ou indirectement) pour une durée déterminée, par opposition à la monarchie où le pouvoir est héréditaire. La \cle{démocratie} est un régime où le pouvoir appartient au peuple (\emph{demos} : peuple, \emph{kratos} : pouvoir).
\item \textbf{Correction de l'erreur} : l'affirmation confond « république » et « limitation du pouvoir ». Une république n'implique pas automatiquement un pouvoir limité : il existe des républiques autoritaires. Ce qui limite le pouvoir, c'est la \cle{séparation des pouvoirs} (exécutif, législatif, judiciaire) théorisée par Montesquieu dans \emph{De l'esprit des lois}, et l'État de droit.
\item \textbf{Application au Cameroun} : la Constitution du 18 janvier 1996 (loi n°96-06) fait du Cameroun une « République décentralisée, unitaire », un État « démocratique » où « la souveraineté nationale appartient au peuple camerounais ». Le président de la République est élu au suffrage universel direct pour sept ans. Mais le régime est à \emph{dominante présidentielle} : le président nomme le Premier ministre et les membres du Gouvernement, qui sont responsables devant lui ; le Parlement (Assemblée nationale et Sénat) contrôle l'action du Gouvernement mais ne peut le renverser par une motion de censure.
\item \textbf{Nuance} : le pouvoir présidentiel est donc juridiquement limité (durée du mandat, existence d'un Parlement bicaméral, d'un Conseil constitutionnel), mais la pratique montre une forte concentration des pouvoirs au profit de l'exécutif.
\item \textbf{Conclusion} : dire qu'un État est une république ne suffit pas à garantir la limitation du pouvoir ; seuls la séparation effective des pouvoirs, l'indépendance de la justice et le contrôle des citoyens assurent cette limitation.
\end{enumerate}
\textbf{Phrase de conclusion rédigée} : La qualité d'un régime ne se mesure pas à son nom (république ou monarchie) mais à l'effectivité des contre-pouvoirs qui empêchent, selon Montesquieu, que « le pouvoir n'arrête le pouvoir ».
\end{exoresolu}

\begin{methode}[Problématiser un sujet de dissertation sur le rôle et la finalité de l'État]
\begin{enumerate}
\item \textbf{Dégager les termes clés} du sujet : « État », « liberté », « sécurité », « bonheur », « intervention ». Définir chacun brièvement.
\item \textbf{Repérer la tension} : le sujet oppose généralement deux exigences (ordre contre liberté ; État gendarme contre État-providence).
\item \textbf{Formuler la problématique} sous forme de question : « L'État doit-il seulement garantir l'ordre, ou doit-il aussi assurer le bien-être des citoyens ? »
\item \textbf{Construire le plan dialectique} en trois parties :
\begin{itemize}
\item Thèse : l'État a pour fonction essentielle la sécurité (Hobbes : sortir de l'état de nature où « l'homme est un loup pour l'homme »).
\item Antithèse : mais cette fonction est insuffisante ; l'État doit aussi promouvoir la justice sociale et le bien-être (État-providence, Rousseau, Rawls).
\item Synthèse : la finalité de l'État est la réalisation de la liberté de tous, ce qui exige à la fois l'ordre et la justice (Hegel : l'État comme réalisation de la liberté concrète).
\end{itemize}
\item \textbf{Rédiger la conclusion} en répondant clairement à la problématique et en ouvrant éventuellement sur les limites de l'action étatique (dette publique, bureaucratie).
\end{enumerate}
\textbf{Réflexe} : chaque partie doit contenir au moins un auteur cité, un argument et un exemple concret (camerounais de préférence).
\end{methode}

\begin{exoresolu}[L'État doit-il seulement maintenir l'ordre ?]
\textbf{Énoncé} (type épreuve) : « L'État a-t-il pour seule fonction de maintenir l'ordre ? » Rédigez l'introduction problématisée et le plan détaillé de la dissertation.
\tcblower
\textbf{Solution détaillée.}

\textbf{Introduction.} Amorce : Toute vie en société suppose un pouvoir capable de faire respecter des règles communes. Définition : l'\cle{État} est l'organisation politique souveraine qui exerce l'autorité sur une population et un territoire donnés. Sa fonction la plus visible est le maintien de l'ordre : police, justice, armée. Tension : mais réduire l'État à cette fonction répressive, c'est oublier qu'il intervient aussi dans l'éducation, la santé, l'économie. Problématique : l'État a-t-il pour seule fonction de maintenir l'ordre, ou doit-il aussi œuvrer au bien-être et à la justice sociale ? Annonce du plan : nous verrons d'abord que le maintien de l'ordre est la fonction première de l'État, puis qu'elle est insuffisante, enfin que la finalité véritable de l'État est la liberté de tous.

\textbf{Plan détaillé.}
\begin{itemize}
\item \textbf{I. Le maintien de l'ordre, fonction première de l'État.} Sans ordre, la vie sociale est impossible : Hobbes montre dans le \emph{Léviathan} que l'état de nature est une « guerre de tous contre tous » et que seul un pouvoir souverain fort garantit la sécurité. Exemple : en période de crise sécuritaire, c'est à l'État qu'incombe la protection des populations (rôle des forces de défense et de sécurité au Cameroun).
\item \textbf{II. Mais cette fonction est insuffisante.} Un État qui ne fait que réprimer devient policier et injuste. Rousseau montre que le pouvoir légitime repose sur le contrat social et la volonté générale : l'État doit servir l'intérêt commun, pas seulement imposer l'obéissance. L'\cle{État-providence} intervient dans la santé (hôpitaux publics), l'éducation (écoles publiques), la protection sociale. Exemple : les programmes de lutte contre la pauvreté ou les subventions aux intrants agricoles.
\item \textbf{III. La finalité de l'État : la liberté et la justice pour tous.} Hegel voit dans l'État la réalisation de la liberté concrète : l'ordre n'a de sens que s'il permet à chacun de développer ses capacités. L'État de droit concilie autorité et liberté en soumettant le pouvoir lui-même à la loi.
\end{itemize}

\textbf{Conclusion attendue} : le maintien de l'ordre est la condition, non la finalité de l'État ; sa fin véritable est de rendre possible une vie libre et digne pour tous les citoyens, ce qui justifie son action sociale autant que sécuritaire.
\end{exoresolu}

\begin{methode}[Justifier une réponse ou une conclusion dans une copie de philosophie]
Le savoir-faire « rédiger et justifier » exige une écriture argumentative rigoureuse :
\begin{enumerate}
\item \textbf{Affirmer clairement} la réponse (une phrase de thèse sans ambiguïté : « oui », « non », ou « oui, mais à condition que... »).
\item \textbf{Justifier par un argument} : donner la raison qui fonde la réponse, introduite par « car », « en effet », « parce que ».
\item \textbf{Étayer par une référence} : citer brièvement un auteur avec sa thèse (pas de citation apprise par cœur sans lien avec l'argument).
\item \textbf{Illustrer par un exemple} concret et daté si possible (institution camerounaise, fait d'actualité).
\item \textbf{Conclure en répondant à la problématique} : reprendre les termes exacts du sujet, sans introduire d'idée nouvelle.
\end{enumerate}
\textbf{Formule à retenir} : une idée = un paragraphe = thèse + argument + référence + exemple + mini-conclusion.
\end{methode}

\begin{exoresolu}[Rédiger une conclusion de dissertation]
\textbf{Énoncé} : Rédigez la conclusion (10 à 15 lignes) d'une dissertation sur le sujet : « Le pouvoir politique a-t-il besoin d'être légitime ? »
\tcblower
\textbf{Solution détaillée.}

\textit{Étape 1 : rappel de la problématique.} Nous nous demandions si le pouvoir politique a besoin d'être légitime, c'est-à-dire si la force suffit à gouverner ou si le pouvoir doit être reconnu comme juste par ceux qui obéissent.

\textit{Étape 2 : bilan du parcours.} Nous avons d'abord montré, avec Hobbes, que tout pouvoir repose sur une force capable d'imposer l'ordre ; puis, avec Max Weber, que la domination n'est durable que lorsqu'elle est crue légitime, que ce soit par la tradition, par le charisme ou par la légalité rationnelle ; enfin, avec Rousseau, que la force ne crée aucun droit : « le plus fort n'est jamais assez fort pour être toujours le maître, s'il ne transforme sa force en droit et l'obéissance en devoir ».

\textit{Étape 3 : réponse claire.} Le pouvoir politique a donc bien besoin d'être légitime : un pouvoir fondé sur la seule contrainte engendre la peur, la révolte et l'instabilité, comme le montrent les régimes qui s'effondrent dès que la force leur fait défaut.

\textit{Étape 4 : exemple.} Au Cameroun, le respect des procédures électorales et constitutionnelles vise précisément à fonder l'autorité des gouvernants sur le consentement des citoyens et non sur la seule force.

\textit{Étape 5 : ouverture.} Reste à savoir si la légitimité légale suffit, ou si le pouvoir doit aussi être juste dans ses actes pour mériter durablement l'adhésion des citoyens.

\textbf{Remarque méthodologique} : cette conclusion répond aux termes du sujet (« besoin », « légitime »), synthétise les trois parties sans les répéter, et ouvre sur une question voisine sans y répondre.
\end{exoresolu}

\begin{methode}[Conduire une analyse de l'action publique locale (TP intégré)]
Pour cartographier les institutions et analyser une action de l'État dans sa commune :
\begin{enumerate}
\item \textbf{Recenser les institutions} présentes localement : sous-préfecture, mairie, tribunal, gendarmerie, lycée public, hôpital public, et les classer selon le pouvoir dont elles relèvent (exécutif, judiciaire, services publics).
\item \textbf{Identifier une action publique} concrète : construction d'une route, d'un forage, d'une salle de classe, campagne de vaccination.
\item \textbf{Déterminer la fonction étatique} qu'elle illustre : régalienne (sécurité, justice), sociale (santé, éducation), économique (infrastructures).
\item \textbf{Évaluer l'action} : répond-elle à un besoin de la population ? Rencontre-t-elle des limites (retards, financement, corruption) ?
\item \textbf{Dégager l'enjeu philosophique} : cette action manifeste-t-elle un État gendarme ou un État-providence ? Est-elle conforme à la finalité de l'État (bien commun) ?
\end{enumerate}
\end{methode}

\begin{exoresolu}[Analyse d'une action publique : la construction d'un centre de santé]
\textbf{Énoncé} : Dans un arrondissement de la région du Centre, l'État construit un centre de santé intégré et y affecte du personnel. Appliquez la démarche d'analyse de l'action publique et dégagez l'enjeu philosophique.
\tcblower
\textbf{Solution détaillée.}
\begin{enumerate}
\item \textbf{Institutions impliquées} : le ministère de la Santé publique (pouvoir exécutif central), le gouverneur et le préfet (administration déconcentrée), la commune (collectivité territoriale décentralisée), le personnel de santé (fonction publique).
\item \textbf{Action publique} : construction et équipement d'un centre de santé, affectation de médecins et d'infirmiers.
\item \textbf{Fonction étatique illustrée} : il s'agit d'une fonction \cle{sociale} (protection de la santé), et non d'une fonction régalienne. L'État n'agit pas ici comme gardien de l'ordre mais comme pourvoyeur de bien-être : c'est la logique de l'\cle{État-providence}.
\item \textbf{Évaluation} : l'action répond à un besoin réel (accès aux soins des populations rurales, réduction des évacuations sanitaires coûteuses vers Yaoundé). Ses limites possibles : sous-équipement, pénurie de médicaments, absentéisme du personnel, qui montrent l'écart entre le droit et sa réalisation effective.
\item \textbf{Enjeu philosophique} : cette action confirme que la finalité de l'État dépasse le maintien de l'ordre. Comme le montre Rousseau, le pouvoir légitime vise la volonté générale, c'est-à-dire l'intérêt commun ; or la santé est une condition de la liberté réelle des citoyens : un malade sans soins n'est pas véritablement libre. L'État-providence apparaît ainsi comme l'accomplissement, et non comme la trahison, de la mission de l'État.
\end{enumerate}
\textbf{Phrase de conclusion rédigée} : La construction d'un centre de santé manifeste que l'État, loin de se réduire à un gendarme, assume une finalité de justice sociale qui conditionne l'exercice effectif des libertés de tous.
\end{exoresolu}

\begin{methode}[Traiter un sujet de type épreuve sur l'État et le pouvoir : gestion du temps et rédaction]
\begin{enumerate}
\item \textbf{Lire tous les sujets} et choisir celui dont on maîtrise les notions (5 minutes).
\item \textbf{Analyser le sujet au brouillon} : définir les termes, dégager le présupposé, formuler la problématique (15 minutes).
\item \textbf{Mobiliser le cours} : lister auteurs, thèses et exemples pertinents ; écarter ce qui est hors sujet (10 minutes).
\item \textbf{Rédiger} : introduction (amorce, définitions, problématique, annonce du plan), développement en trois parties avec transitions, conclusion (environ 2 heures).
\item \textbf{Relire} : orthographe, accords, cohérence, présence des termes du sujet dans la conclusion (10 minutes).
\end{enumerate}
\textbf{Sujets probables} : « L'État est-il un mal nécessaire ? » ; « La force suffit-elle à fonder le pouvoir ? » ; « L'État doit-il tout faire pour les citoyens ? » ; « Obéir à l'État, est-ce renoncer à sa liberté ? »
\end{methode}

\begin{exoresolu}[La force suffit-elle à fonder le pouvoir ?]
\textbf{Énoncé} (type épreuve) : « La force suffit-elle à fonder le pouvoir politique ? » Rédigez la problématique, le plan et le développement de la première partie.
\tcblower
\textbf{Solution détaillée.}

\textbf{Problématique} : le pouvoir politique dispose de la force (police, armée, prisons) ; mais la force seule peut-elle fonder un pouvoir durable et juste, ou le pouvoir a-t-il besoin d'être reconnu comme légitime ?

\textbf{Plan} : I. La force est la condition du pouvoir (Hobbes). II. Mais la force seule ne fonde aucun droit (Rousseau). III. Le pouvoir durable repose sur la légitimité (Weber).

\textbf{Développement de la première partie.}

\textit{Thèse} : la force est la condition première de tout pouvoir politique.

\textit{Argument} : sans capacité de contrainte, les lois resteraient des vœux pieux ; l'État se distingue précisément de toute autre association par le monopole de la violence légitime. Un État qui ne peut plus faire respecter ses décisions cesse d'exister comme État : on parle alors d'État failli.

\textit{Référence} : Hobbes, dans le \emph{Léviathan} (1651), montre qu'à l'état de nature, chacun ayant droit sur tout, la vie est « solitaire, misérable, dangereuse, animale et brève ». Seule l'institution d'un souverain doté d'une puissance redoutable met fin à cette guerre de tous contre tous : c'est la crainte du châtiment qui rend les conventions respectées.

\textit{Exemple} : lorsqu'une crise sécuritaire éclate, c'est la capacité de l'État à déployer ses forces de défense et de sécurité qui restaure l'autorité de la loi sur le territoire ; inversement, dans les zones où l'État est absent, ce sont des groupes armés qui imposent leur loi par la seule force.

\textit{Transition} : mais si la force est nécessaire, est-elle suffisante ? Un pouvoir qui ne repose que sur la peur est-il véritablement fondé ? C'est ce qu'il faut examiner dans une seconde partie, avec Rousseau : « la force est une puissance physique ; je ne vois point quelle moralité peut résulter de ses effets ».
\end{exoresolu}

\begin{attention}[Les 5 erreurs qui coûtent des points au Baccalauréat technique]
\begin{enumerate}
\item \textbf{Confondre État, gouvernement et nation.} L'État est l'organisation politique permanente (population, territoire, souveraineté) ; le gouvernement n'en est que l'organe exécutif, changeant ; la nation est une communauté culturelle et historique. Écrire « le gouvernement est souverain » est une erreur de notion.
\item \textbf{Confondre les classifications des régimes.} Ne pas mélanger le critère du nombre de gouvernants (Aristote : monarchie, aristocratie, démocratie) avec le critère de l'organisation des pouvoirs (régime présidentiel, parlementaire). Dire « le Cameroun est une monarchie parce que le président est puissant » est un contresens.
\item \textbf{Citer un auteur sans l'utiliser.} Juxtaposer « Hobbes, Rousseau, Weber » sans expliquer leur thèse ni la relier à l'argument ne rapporte aucun point : chaque référence doit servir le raisonnement.
\item \textbf{Oublier de répondre à la problématique dans la conclusion.} Une conclusion qui résume sans trancher (« nous avons vu que... ») est pénalisée : il faut répondre explicitement, avec les termes du sujet.
\item \textbf{Rester dans le vague et l'abstrait.} Une copie sans aucun exemple concret (institution camerounaise, fait précis) paraît hors sol : chaque partie doit contenir au moins un exemple précis et correctement analysé, sans se transformer en récit journalistique.
\end{enumerate}
\end{attention}

\newpage

\coursec{Exercices}

\courssub{Je vérifie mes connaissances}

\begin{exercice}[QCM : les fondamentaux de l'État]
Pour chaque question, une seule réponse est exacte. Recopie la lettre correspondante.
\begin{enumerate}
\item Le monopole de la violence légitime, selon Max Weber, désigne :
\begin{enumerate}
\item[a)] le droit de tout citoyen de se défendre ;
\item[b)] le fait que l'État seul peut user légalement de la contrainte physique ;
\item[c)] l'interdiction faite à l'armée d'intervenir sur le territoire.
\end{enumerate}
\item Les quatre éléments constitutifs de l'État (Convention de Montevideo, 1933) sont :
\begin{enumerate}
\item[a)] un territoire, une population, un gouvernement, la souveraineté ;
\item[b)] un drapeau, un hymne, une monnaie, une armée ;
\item[c)] un roi, un peuple, une religion, une capitale.
\end{enumerate}
\item Selon Weber, la légitimité fondée sur la croyance en des qualités exceptionnelles d'un chef est la légitimité :
\begin{enumerate}
\item[a)] traditionnelle ; \item[b)] légale-rationnelle ; \item[c)] charismatique.
\end{enumerate}
\item Au Cameroun, la souveraineté nationale appartient :
\begin{enumerate}
\item[a)] au Président de la République ; \item[b)] au peuple, qui l'exerce par voie de référendum ou par ses représentants ; \item[c)] à l'Assemblée nationale seule.
\end{enumerate}
\end{enumerate}
\end{exercice}

\begin{exercice}[Vrai ou faux — justifier]
Indique si chaque affirmation est vraie ou fausse, puis justifie ta réponse en deux ou trois lignes.
\begin{enumerate}
\item La Nation et l'État sont deux notions strictement identiques.
\item Une loi légale est toujours une loi légitime.
\item Le Premier Ministre, au Cameroun, est le chef du Gouvernement, chargé de la mise en œuvre de la politique de la Nation définie par le Président de la République.
\item Le principe de non-affectation budgétaire signifie que chaque impôt doit financer obligatoirement une dépense précise.
\end{enumerate}
\end{exercice}

\begin{exercice}[Définitions]
Définis en trois lignes maximum chacune des notions suivantes, en donnant un exemple :
\begin{enumerate}
\item l'État ;
\item la légitimité ;
\item la séparation des pouvoirs ;
\item l'État-providence.
\end{enumerate}
\end{exercice}

\begin{exercice}[Questions de cours]
Réponds brièvement aux questions suivantes.
\begin{enumerate}
\item Cite et caractérise en une phrase les trois types de légitimité du pouvoir selon Max Weber.
\item Distingue la \cle{légalité} de la \cle{légitimité}.
\item Quelle est la différence entre un régime présidentiel et un régime parlementaire ?
\item Rappelle la définition du contrat social chez Rousseau.
\end{enumerate}
\end{exercice}

\courssub{Je m'entraîne}

\begin{exercice}[Identifier les formes de gouvernement]
Pour chacune des situations suivantes, indique la forme de gouvernement ou le type de régime concerné, en justifiant ta réponse.
\begin{enumerate}
\item Un pays où le pouvoir est détenu par un monarque héréditaire dont les décisions ne sont limitées par aucune constitution.
\item Un pays où le chef de l'État est élu au suffrage universel direct et où le gouvernement n'est pas responsable devant le Parlement.
\item Un pays où le gouvernement peut être renversé par un vote de censure du Parlement.
\item Un pays où un petit groupe de militaires exerce le pouvoir sans élections.
\end{enumerate}
\end{exercice}

\begin{exercice}[Appliquer la typologie de Weber]
Voici trois situations camerounaises ou africaines. Pour chacune, indique le type de légitimité (traditionnelle, charismatique, légale-rationnelle) qui domine, et explique pourquoi.
\begin{enumerate}
\item Un lamido du Nord-Cameroun exerce une autorité reconnue sur sa communauté, héritée de ses ancêtres.
\item Un chef d'État est élu conformément à la Constitution et exerce son pouvoir dans le cadre des lois en vigueur.
\item Un leader politique entraîne les foules par son éloquence exceptionnelle et la confiance quasi mystique qu'il inspire.
\end{enumerate}
\end{exercice}

\begin{exercice}[Cas pratique : le conflit entre le Maire et le Préfet]
Le Maire de la Commune d'Arrondissement de X refuse d'appliquer une décision du Préfet (représentant de l'État) qu'il juge illégale, car contraire aux compétences reconnues aux collectivités territoriales décentralisées.
\begin{enumerate}
\item Rappelle ce qu'est la décentralisation et ce qu'est la tutelle administrative.
\item Le Préfet peut-il annuler toutes les décisions du Maire ? Justifie ta réponse à l'aide de la hiérarchie des normes.
\item Qui tranche en dernier ressort ce type de conflit de compétence ?
\item Dégage la leçon philosophique : ce conflit montre-t-il une faiblesse ou une force de l'État de droit ?
\end{enumerate}
\end{exercice}

\begin{exercice}[Schéma à compléter]
Reproduis et complète le tableau comparatif suivant :
\begin{center}
\begin{tabularx}{\linewidth}{|l|X|X|}
\hline
\rowcolor{popblueL} \textbf{Critère} & \textbf{Régime présidentiel} & \textbf{Régime parlementaire} \\
\hline Séparation des pouvoirs & & \\
\hline Mode de désignation du chef de l'État & & \\
\hline Responsabilité du gouvernement devant le Parlement & & \\
\hline Exemple de pays & & \\
\hline
\end{tabularx}
\end{center}
Puis réponds : le régime camerounais est-il un régime présidentiel pur ? Argumente en deux arguments.
\end{exercice}

\courssub{Je me prépare au Bac}

\begin{exercice}[Dissertation : sujet type Baccalauréat technique]
\textbf{Sujet :} \emph{Obéir à l'État, est-ce renoncer à sa liberté ?}
\begin{enumerate}
\item Analyse le sujet : dégage les notions principales, formule le problème et construis la problématique.
\item Élabore un plan détaillé en trois parties (thèse, antithèse, synthèse) en mobilisant :
\begin{itemize}
\item Hobbes : l'État comme condition de sortie de l'état de nature ;
\item Locke : l'État garant des droits naturels et le droit de résistance ;
\item Rousseau : la volonté générale et la liberté civile ;
\item la distinction entre légalité et légitimité.
\end{itemize}
\item Rédige l'introduction complète de la dissertation (accroche, définition des termes, problématique, annonce du plan).
\item Illustre ta réflexion par un exemple tiré de la vie politique camerounaise (obéissance aux lois, devoirs du citoyen, contestation d'une décision jugée illégitime).
\end{enumerate}
\end{exercice}

\begin{exercice}[Explication de texte : Rousseau, \emph{Du Contrat social}]
\textbf{Texte :} « Trouver une forme d'association qui défende et protège de toute la force commune la personne et les biens de chaque associé, et par laquelle chacun, s'unissant à tous, n'obéisse pourtant qu'à lui-même, et reste aussi libre qu'auparavant. » (Rousseau, \emph{Du Contrat social}, Livre I, chapitre 6.)
\begin{enumerate}
\item Dégage la thèse de l'auteur et la structure de l'argumentation.
\item Explique l'expression « aliéner sa liberté » : que cède l'individu en entrant dans la société, et que reçoit-il en échange ?
\item En quoi la \cle{volonté générale} diffère-t-elle de la simple volonté de tous (somme des volontés particulières) ?
\item Comment Rousseau résout-il le paradoxe : « n'obéir qu'à soi-même » tout en obéissant aux lois ?
\item Discussion : cette théorie du contrat social explique-t-elle encore la légitimité de l'État moderne, par exemple au Cameroun ? Justifie ta réponse.
\end{enumerate}
\end{exercice}

\begin{exercice}[Analyse de document : dette publique et souveraineté]
\textbf{Document 1 :} Évolution du taux d'endettement public du Cameroun (dette publique en \% du PIB) :
\begin{center}
\begin{tabularx}{\linewidth}{|l|X|X|X|X|X|}
\hline
\rowcolor{popblueL} \textbf{Année} & 2010 & 2015 & 2020 & 2022 & 2024 \\
\hline Dette publique (\% du PIB) & 15 & 31 & 45 & 46 & 42 \\
\hline
\end{tabularx}
\end{center}
\textbf{Document 2 :} Extrait inspiré de la Loi de Finances 2024 : « L'État mobilise des ressources extérieures et intérieures pour financer les grands projets structurants : routes, barrages hydroélectriques, ports et infrastructures sportives. »
\begin{enumerate}
\item Décris l'évolution de la courbe entre 2010 et 2024 (tendance générale, périodes de hausse et de stabilisation).
\item Calcule le taux de variation de l'endettement entre 2010 et 2020, puis entre 2020 et 2024.
\item Explique deux causes possibles de cette évolution (chocs extérieurs, investissements publics, sécurité).
\item Quels risques un endettement élevé fait-il peser sur la souveraineté économique de l'État ?
\item Propose deux mesures de redressement et montre, pour chacune, la conception du rôle de l'État qu'elle suppose (État libéral ou État-providence).
\end{enumerate}
\end{exercice}

\newpage

\coursec{Corrigés des exercices}

\begin{corrige}[Exercice 1 — QCM : les fondamentaux de l'État]
\begin{enumerate}
\item \textbf{b)} Selon Max Weber (\emph{Le Savant et le Politique}), l'État détient le « monopole de la violence physique légitime » : lui seul a le droit d'user de la contrainte physique de manière légale sur un territoire donné.
\item \textbf{a)} La Convention de Montevideo (1933), article 1, pose les quatre critères classiques de l'État en droit international : population permanente, territoire défini, gouvernement effectif, capacité d'entrer en relation avec les autres États (souveraineté).
\item \textbf{c)} Weber distingue trois types de légitimité : traditionnelle (coutume), légale-rationnelle (loi/procédure) et \textbf{charismatique} (dévouement à la personne exceptionnelle du chef).
\item \textbf{b)} Constitution du Cameroun (Loi n°96-06 du 18 janvier 1996), Préambule et Article 2 : « La souveraineté nationale appartient au peuple qui l'exerce par voie de référendum ou par l'intermédiaire de ses représentants. »
\end{enumerate}
\end{corrige}

\begin{corrige}[Exercice 2 — Vrai ou faux — justifier]
\begin{enumerate}
\item \textbf{Faux.} La \cle{Nation} est une communauté humaine consciente de son unité (langue, culture, histoire, volonté de vivre ensemble) : notion sociologique/subjective. L'\cle{État} est une organisation politique juridique dotée de la souveraineté sur un territoire : notion juridique/objective. Il existe des Nations sans État (ex: Kurdes) et des États multi-nationaux (ex: Cameroun, Belgique).
\item \textbf{Faux.} La \cle{légalité} est la conformité formelle à la loi positive en vigueur (procédure respectée, texte publié). La \cle{légitimité} est l'acceptation morale, le bien-fondé aux yeux des gouvernés (justice, consentement). Une loi peut être légale (votée régulièrement) mais illégitime (injuste, discriminatoire) — ex: lois d'apartheid, lois de Vichy.
\item \textbf{Vrai.} Constitution du Cameroun, Art. 10 : « Le Premier Ministre est le Chef du Gouvernement. Il met en œuvre la politique de la Nation définie par le Président de la République. » Il assure l'exécution des lois et dirige l'administration.
\item \textbf{Faux.} Le principe de \textbf{non-affectation} (ou universalité budgétaire) signifie que \emph{toutes} les recettes financent \emph{toutes} les dépenses sans lien direct (budget global). L'affectation spéciale (ex: taxe routière pour l'entretien des routes) en est l'exception stricte, encadrée par la loi.
\end{enumerate}
\end{corrige}

\begin{corrige}[Exercice 3 — Définitions]
\begin{enumerate}
\item \textbf{État :} Personne morale de droit public, dotée de la souveraineté, s'exerçant sur un territoire délimité, une population permanente, par l'intermédiaire d'un gouvernement (Montevideo 1933). \emph{Ex: La République du Cameroun.}
\item \textbf{Légitimité :} Croyance en le bien-fondé de l'autorité qui commande l'obéissance volontaire, sans contrainte pure (Weber). Elle repose sur la tradition, le charisme ou la légalité rationnelle. \emph{Ex: Un président élu au suffrage universel transparent.}
\item \textbf{Séparation des pouvoirs :} Principe constitutionnel (Montesquieu, \emph{De l'esprit des lois}) distinguant trois fonctions — législative (vote la loi), exécutive (applique la loi), judiciaire (sanctionne les violations) — confiées à des organes distincts pour prévenir la tyrannie. \emph{Ex: Au Cameroun, Parlement (législatif), Président/PM (exécutif), Tribunaux (judiciaire).}
\item \textbf{État-providence :} Modèle où l'État intervient massivement dans la sphère sociale (santé, éducation, retraite, chômage) pour corriger les inégalités de marché et garantir un minimum vital (rapport Beveridge 1942). \emph{Ex: Couverture Maladie Universelle (CMU) en cours de déploiement au Cameroun ; modèle scandinave.}
\end{enumerate}
\end{corrige}

\begin{corrige}[Exercice 4 — Questions de cours]
\begin{enumerate}
\item \textbf{Weber — Trois types de légitimité :}
\begin{itemize}
\item \textbf{Traditionnelle :} Repose sur la coutume, l'habitude immémoriale, la sacralité de l'hérédité (ex: monarchie féodale, chefferies traditionnelles).
\item \textbf{Charismatique :} Repose sur la dévotion à la personne du chef perçue comme porteuse de qualités exceptionnelles, surnaturelles ou héroïques (ex: leader révolutionnaire, prophète).
\item \textbf{Légale-rationnelle :} Repose sur la croyance en la légalité des règles édictées rationnellement et en le droit de ceux qui sont désignés par ces règles à donner des ordres (ex: démocratie moderne, bureaucratie).
\end{itemize}
\item \textbf{Légalité vs Légitimité :} La légalité est formelle : conformité à la norme juridique en vigueur (procédure, compétence). La légitimité est morale : adhésion des gouvernés, conformité à la justice ou au bien commun. \emph{Maxime : « Tout ce qui est légal n'est pas légitime ; tout ce qui est légitime n'est pas toujours légal » (résistance à l'oppression).}
\item \textbf{Régime présidentiel vs parlementaire :}
\begin{itemize}
\item \textbf{Présidentiel :} Séparation stricte des pouvoirs. Chef de l'État = Chef de l'exécutif, élu au suffrage universel. Gouvernement non responsable devant le Parlement (pas de motion de censure). Parlement ne peut renverser le gouvernement. Ex: USA.
\item \textbf{Parlementaire :} Séparation souple (collaboration). Dualité exécutif : Chef de l'État (symbolique) \neq Chef de gouvernement (PM, réel). Gouvernement \textbf{responsable} devant le Parlement (vote de censure, question de confiance). Parlement peut dissoudre le gouvernement ; exécutif peut dissoudre le Parlement. Ex: Royaume-Uni, Allemagne.
\end{itemize}
\item \textbf{Contrat social (Rousseau, \emph{Du Contrat social}, L. I, ch. 6) :} « Chacun de nous met en commun sa personne et toute sa puissance sous la suprême direction de la volonté générale ; et nous recevons chaque membre comme partie indivisible du tout. » L'individu s'aliène totalement à la communauté pour retrouver sa liberté sous forme de \textbf{liberté civile} (obéissance à la loi qu'il s'est prescrite) et de \textbf{propriété} (droit sur ses biens garanti par la communauté).
\end{enumerate}
\end{corrige}

\begin{corrige}[Exercice 5 — Identifier les formes de gouvernement]
\begin{enumerate}
\item \textbf{Monarchie absolue.} Justification : Pouvoir détenu par un seul (monarque), transmission héréditaire, absence de constitution limitant le pouvoir (pas de séparation des pouvoirs, pas de contrôle parlementaire).
\item \textbf{Régime présidentiel.} Justification : Chef de l'État élu au suffrage universel direct (légitimité populaire propre) ; il nomme le gouvernement qui lui est responsable (pas de responsabilité politique devant le Parlement) ; séparation stricte des pouvoirs (exécutif ne peut dissoudre le législatif, législatif ne peut renverser l'exécutif).
\item \textbf{Régime parlementaire.} Justification : Caractère distinctif = \textbf{responsabilité politique du gouvernement devant le Parlement}. Le vote de censure (motion de défiance) contraint le gouvernement à démissionner. Le PM émane de la majorité parlementaire.
\item \textbf{Dictature militaire / Junta / Régime autoritaire.} Justification : Accès au pouvoir par la force (coup d'État), non par les élections. Pouvoir concentré entre les mains d'un petit groupe (comité militaire). Suspension de la Constitution, absence de séparation des pouvoirs, répression de l'opposition.
\end{enumerate}
\end{corrige}

\begin{corrige}[Exercice 6 — Appliquer la typologie de Weber]
\begin{enumerate}
\item \textbf{Légitimité traditionnelle.} Le lamido tire son autorité de la coutume, de l'héritage dynastique et de la sacralité de la fonction transmise par les ancêtres. L'obéissance est due à la personne \emph{en raison de sa lignée} (« c'est ainsi de tout temps »).
\item \textbf{Légitimité légale-rationnelle.} Le chef d'État accède au pouvoir par une procédure juridique codifiée (Constitution, Code électoral), dans le respect des règles impersonnelles. L'obéissance est due à la \emph{fonction} et à la \emph{légalité} de la procédure, non à la personne.
\item \textbf{Légitimité charismatique.} L'autorité repose sur la reconnaissance personnelle des qualités exceptionnelles du leader (éloquence, courage, vision mystique). L'obéissance est due à la \emph{personne} perçue comme hors du commun ; elle est instable, liée à la preuve continue du charisme (victoires, miracles).
\end{enumerate}
\end{corrige}

\begin{corrige}[Exercice 7 — Cas pratique : le conflit entre le Maire et le Préfet]
\begin{enumerate}
\item \textbf{Décentralisation :} Transfert par l'État de compétences, de ressources (fiscales, financières) et de moyens à des \textbf{collectivités territoriales décentralisées} (CTD : communes, régions) dotées de la personnalité morale et de l'autonomie de gestion (libre administration, Art. 55 Const.). \textbf{Tutelle administrative :} Contrôle \emph{a posteriori} de la \textbf{légalité} des actes des CTD par l'État (représenté par le Préfet/Sous-préfet) pour garantir le respect de la loi, de l'intérêt national et des libertés publiques. Ce n'est pas un pouvoir de direction (pas d'injonction de faire), mais un pouvoir d'annulation/suspension des actes illégaux.
\item \textbf{Non, le Préfet ne peut pas tout annuler.} Il exerce un \textbf{contrôle de légalité}, pas un \textbf{contrôle d'opportunité} (mérité). Il ne peut annuler que les actes illégaux (incompétence, violation de la loi, procédure irrégulière, détournement de pouvoir). Hiérarchie des normes : Constitution > Lois > Décrets > Arrêtés préfectoraux > Arrêtés municipaux. Le Maire a une \textbf{compétence propre} (libre administration) dans les domaines dévolus à la commune (état civil, marchés, voirie communale, école primaire) : le Préfet ne peut s'y substituer.
\item \textbf{Le Juge administratif (Tribunal Administratif, Cour d'Appel Administrative, Cour Suprême).} En cas de refus du Maire d'exécuter une décision de suspension/annulation du Préfet, ou de contestation de cette décision par le Maire, le contentieux relève de la juridiction administrative (recours pour excès de pouvoir). C'est elle qui tranche en dernier ressort le conflit de compétence.
\item \textbf{Leçon philosophique : Force de l'État de droit.} Ce conflit montre que \textbf{personne n'est au-dessus de la loi}, pas même le représentant de l'État (Préfet) ni l'élu local (Maire). L'existence d'un juge indépendant pour arbitrer le conflit entre deux autorités publiques est la marque de l'État de droit : la norme prime sur la force, la procédure protège l'autonomie locale contre l'arbitraire centralisateur. C'est une \textbf{garantie de la liberté locale}.
\end{enumerate}
\end{corrige}

\begin{corrige}[Exercice 8 — Schéma à compléter]
\begin{center}
\begin{tabularx}{\linewidth}{|l|X|X|}
\hline
\rowcolor{popblueL} \textbf{Critère} & \textbf{Régime présidentiel} & \textbf{Régime parlementaire} \\
\hline Séparation des pouvoirs & Stricte (indépendance mutuelle, pas de responsabilité politique) & Souple / Flexible (collaboration, responsabilité politique, dissolution possible) \\
\hline Mode de désignation du chef de l'État & Suffrage universel direct (ou grand électeurs) & Souvent indirect (élection par le Parlement) ou héréditaire (monarchie) ; si direct, rôle cérémoniel \\
\hline Responsabilité du gouvernement devant le Parlement & \textbf{Aucune} (gouvernement responsable seulement devant le Président) & \textbf{Totale} : motion de censure, vote de confiance, engagement de responsabilité \\
\hline Exemple de pays & États-Unis, Brésil, \textbf{Cameroun (formel)} & Royaume-Uni, Allemagne, Espagne, Italie, Japon \\
\hline
\end{tabularx}
\end{center}
\textbf{Le régime camerounais est-il un régime présidentiel pur ? \textbf{Non}. Deux arguments :}
\begin{enumerate}
\item \textbf{Existence d'un Premier Ministre Chef du Gouvernement responsable devant l'Assemblée nationale.} L'Art. 34 de la Constitution prévoit la motion de censure : l'Assemblée peut mettre en cause la responsabilité du Gouvernement (vote à la majorité des 2/3). Cela introduit une responsabilité politique parlementaire, étrangère au présidentiel pur (USA).
\item \textbf{Pouvoir de dissolution de l'Assemblée nationale par le Président (Art. 18).} Dans un présidentiel pur, le Président ne peut dissoudre le Parlement (mandats fixes, non révocables). Au Cameroun, le Président peut dissoudre l'Assemblée, ce qui est une arme parlementaire classique (responsabilité du législatif devant l'exécutif).
\end{enumerate}
\emph{Conclusion :} Le Cameroun est un \textbf{régime semi-présidentiel} (ou présidentialiste à responsabilité gouvernementale), proche du modèle français de la V\textsuperscript{e} République.
\end{corrige}

\begin{corrige}[Exercice 9 — Dissertation : Obéir à l'État, est-ce renoncer à sa liberté ?]
\textbf{1. Analyse du sujet}
\begin{itemize}
\item \textbf{Notions :} \cle{Obéissance} (soumission à un ordre/loi), \cle{État} (detenteur du monopole de la contrainte légitime), \cle{Liberté} (pouvoir de choisir, autonomie : naturelle vs civile), \cle{Renoncement} (abandon volontaire ou contraint).
\item \textbf{Problème :} L'obéissance implique contrainte (loi, police) ; la liberté implique absence de contrainte. Y a-t-il contradiction nécessaire ?
\item \textbf{Problématique :} \emph{L'obéissance à l'État est-elle une aliénation (perte de soi) ou la condition même de l'exercice d'une liberté réelle et partagée ?}
\end{itemize}

\textbf{2. Plan détaillé (Thèse / Antithèse / Synthèse)}

\textbf{I. Thèse : Obéir à l'État, c'est renoncer à sa liberté naturelle (aliénation nécessaire).}
\begin{itemize}
\item \textbf{Hobbes (\emph{Léviathan}) :} État de nature = guerre de tous contre tous, insécurité permanente. Liberté naturelle = droit à tout (y compris au corps d'autrui) $\rightarrow$ vie « misérable, brutale, courte ».
\item \textbf{Contrat d'aliénation :} Les hommes cèdent \textbf{toute} leur liberté naturelle au Souverain (Léviathan) en échange de la \textbf{sécurité}. Obéir = respecter le pacte pour ne pas retomber dans le chaos. La loi est la volonté du Souverain ; désobéir = retour à l'état de nature.
\item \textbf{Portée :} La liberté ne subsiste que là où la loi est silencieuse (« silence des lois »). Obéissance = renoncement à l'arbitraire individuel pour la paix.
\end{itemize}

\textbf{II. Antithèse : Obéir à un État légitime, c'est exercer sa liberté civile (pas de renoncement, mais transformation).}
\begin{itemize}
\item \textbf{Locke (\emph{Gouvernement civil}) :} État de nature supportable (loi naturelle, droits naturels : vie, liberté, propriété). Contrat = \textbf{fiducie} : on délègue le pouvoir de faire appliquer la loi naturelle, on ne s'aliène pas.
\item \textbf{Droit de résistance :} Si l'État viole la confiance (tyrannie, atteinte aux biens/vie), le peuple reprend sa liberté naturelle : \textbf{droit de résistance} (insurrection). Obéissance due \emph{seulement} aux lois justes (conformes à la fin de l'État : protection des droits).
\item \textbf{Rousseau (\emph{Du Contrat social}) :} « Aliénation totale » de la liberté naturelle $\rightarrow$ \textbf{Liberté civile} (obéissance à la loi qu'on s'est prescrite). \emph{« Obéir à la loi qu'on s'est prescrite, c'est la liberté. »} La \cle{volonté générale} (intérêt commun) $\neq$ volonté de tous (somme égoïste). Le citoyen est \emph{auteur} et \emph{sujet} de la loi.
\end{itemize}

\textbf{III. Synthèse : La distinction Légalité / Légitimité — Obéissance critique et désobéissance civile.}
\begin{itemize}
\item \textbf{Légalité $\neq$ Légitimité (Kant, Thoreau, Gandhi, Rawls) :} Une loi peut être légale (votée) mais illégitime (injuste, raciste, attentatoire aux droits fondamentaux).
\item \textbf{Devoir d'obéissance aux lois justes :} Condition de la vie commune, de l'égalité devant la loi (État de droit).
\item \textbf{Droit / Devoir de désobéissance civile :} Refus public, non-violent, assumé d'obéir à une loi jugée illégitime, en acceptant la sanction légale, pour appeler à la réforme (Thoreau, \emph{Désobéissance civile} ; King, Mandela). C'est l'acte suprême de liberté morale (conscience > loi positive).
\item \textbf{Conclusion :} Obéir à l'État de droit légitime, ce n'est pas renoncer à sa liberté, c'est \textbf{la fonder}. Obéir aveuglément à un État illégitime, c'est y renoncer. La liberté citoyenne est \textbf{vigilance}.
\end{itemize}

\textbf{3. Introduction complète (modèle Bac)}
\begin{quote}
\textbf{[Accroche]} « L'homme est né libre, et partout il est dans les fers. » Par cette formule choc, Rousseau ouvre \emph{Du Contrat social} et pose d'emblée le paradoxe de la condition politique : l'entrée dans la société semble une chaîne, pourtant elle est présentée comme le passage à l'humanité.
\textbf{[Définitions]} Obéir à l'État, c'est se conformer à ses lois et décisions, contrainte par la force publique (Weber). La liberté, c'est l'autonomie : capacité de se donner sa propre loi (Kant) ou absence d'empêchement extérieur (liberté négative). Renoncer, c'est abandonner un droit, une prérogative.
\textbf{[Problématique]} Dès lors, l'obéissance à la loi étatique est-elle nécessairement une mutilation de la liberté, ou peut-elle en être l'accomplissement ? Autrement dit : \emph{la contrainte légale est-elle l'ennemie ou la condition de la liberté ?}
\textbf{[Annonce du plan]} Nous montrerons d'abord que, pour Hobbes, l'obéissance est bien un renoncement à la liberté naturelle indispensable à la sécurité (I). Nous verrons ensuite que, pour Locke et Rousseau, l'obéissance à un État légitime transforme la liberté naturelle en liberté civile supérieure, sans renoncement à l'autonomie (II). Enfin, nous argumenterons que la distinction cruciale entre légalité et légitimité fonde un devoir d'obéissance critique, où la désobéissance civile devient l'acte de liberté par excellence (III).
\end{quote}

\textbf{4. Exemple camerounais}
La \textbf{loi n°2023/004 du 25 juillet 2023 portant Code de la route} impose le port du casque pour les mototaximen et leurs passagers.
\begin{itemize}
\item \textbf{Obéissance (légalité) :} Le mototaximan porte le casque par peur de l'amende (contrainte étatique). Il renonce à sa « liberté » de rouler tête nue.
\item \textbf{Liberté civile (sécurité) :} Cette contrainte protège sa vie et celle du passager (droit à la vie, Art. 1er Const.). Il obéit à une loi qu'il a acceptée via ses députés (représentation) : il s'obéit à lui-même comme citoyen.
\item \textbf{Contestation (légitimité) :} Si le Préfet interdit \emph{toute} manifestation syndicale sans motif valable (atteinte disproportionnée à la liberté de réunion, Art. 29 Const.), le syndicat saisit le Tribunal Administratif (recours pour excès de pouvoir). Le juge annule l'arrêté : l'obéissance n'était pas due car l'acte était \textbf{illégal} (donc illégitime). Le citoyen a usé de sa liberté pour faire respecter la loi supérieure.
\end{itemize}

\textbf{Barème indicatif (sur 20) :} Analyse/Problématique (4) + Plan cohérent (4) + Introduction rédigée (4) + Exemple camerounais pertinent (3) + Mobilisation auteurs/notions (3) + Qualité langue/structure (2).
\textbf{Points de vigilance :} Ne pas confondre liberté naturelle / civile. Ne pas dire « Rousseau = totalitarisme » (erreur classique : la volonté générale protège l'individu). Citer explicitement la distinction légalité/légitimité en synthèse. L'exemple doit être \emph{concret} et \emph{analysé}, pas juste cité.
\end{corrige}

\begin{corrige}[Exercice 10 — Explication de texte : Rousseau, \emph{Du Contrat social}]
\textbf{1. Thèse et structure}
\begin{itemize}
\item \textbf{Thèse :} Le contrat social permet de résoudre le paradoxe de l'association politique : comment s'unir pour se protéger sans s'asservir ? La solution est l'aliénation totale de chaque associé à la communauté entière, qui génère la \textbf{volonté générale}, seule source de lois légitimes où obéir c'est s'obéir à soi-même.
\item \textbf{Structure :} (a) Énoncé du problème : trouver une association protectrice qui préserve la liberté. (b) Solution : aliénation totale de soi et de ses biens à la communauté. (c) Conséquence : chacun, se donnant à tous, ne se donne à personne (échange égalitaire) ; il n'obéit qu'à lui-même (autonomie).
\end{itemize}

\textbf{2. « Aliéner sa liberté »}
\begin{itemize}
\item \textbf{Cède (liberté naturelle) :} Le droit illimité à tout ce qui le tente et qu'il peut atteindre (force, ruse). L'indépendance naturelle, l'absence de toute règle supérieure à sa volonté propre.
\item \textbf{Reçoit (liberté civile) :} La liberté limitée par la volonté générale (loi), garantie par la force publique. La \textbf{propriété} (droit légitime sur ses biens, pas simple possession de fait). La \textbf{liberté morale} (obéissance à une loi qu'on s'est prescrite = autonomie). La dignité d'être \emph{maître de soi} par la loi, non esclave de ses appétits.
\end{itemize}

\textbf{3. Volonté générale vs Volonté de tous}
\begin{itemize}
\item \textbf{Volonté de tous :} Somme arithmétique des volontés particulières (intérêts égoïstes, privés). Peut être contradictoire, instable, injuste. « Souvent la volonté de tous n'est que la somme de petites différences. »
\item \textbf{Volonté générale :} Volonté du corps politique considéré comme un tout unitaire. Vise \textbf{l'intérêt commun / le bien commun}. Elle est \textbf{toujours droite} (tend au bien), mais le peuple peut se tromper sur ce qu'elle est (besoin de législateur, de délibération sans factions). Elle résulte de l'annulation des « plus et des moins » (intérêts contraires).
\end{itemize}

\textbf{4. Résolution du paradoxe « n'obéir qu'à soi-même »}
\begin{itemize}
\item Double qualité du contractant : \textbf{Souverain} (comme citoyen, il vote la loi = auteur) et \textbf{Sujet} (comme particulier, il subit la loi = destinataire).
\item La loi n'est pas un ordre extérieur : elle émane de la volonté générale dont je suis partie intégrante. En votant la loi, je veux ce qui est bon pour tous (donc pour moi). En l'appliquant, j'exécute ma propre volonté de citoyen.
\item \emph{« Chacun, s'unissant à tous, n'obéit pourtant qu'à lui-même »} : L'obéissance à la loi = obéissance à sa propre raison universalisée. C'est la définition de l'\textbf{autonomie} (Kant reprendra cette idée).
\end{itemize}

\textbf{5. Discussion : Pertinence au Cameroun aujourd'hui}
\begin{itemize}
\item \textbf{Oui, comme idéal régulateur :} La Constitution camerounaise (Préambule) proclame que « la souveraineté nationale appartient au peuple ». Les élections (présidentielles, législatives, municipales) sont les moments où le peuple exerce sa souveraineté (délégation). L'État de droit vise à faire coïncider loi et volonté générale (séparation des pouvoirs, Cour Constitutionnelle, Tribunal Administratif).
\item \textbf{Non, comme description réelle (limites) :}
\begin{itemize}
\item \textbf{Déficit de participation :} Abstention forte, défiance envers la classe politique $\rightarrow$ sentiment que les lois ne reflètent pas la volonté générale mais des intérêts partisans (volonté de tous / factions).
\item \textbf{Présidence forte / hyper-exécutif :} Concentration des pouvoirs (nomination PM, magistrats, hauts fonctionnaires, dissolution AN) $\rightarrow$ risque de confusion entre volonté du Prince et volonté générale.
\item \textbf{Pluralisme ethnique/régional :} Difficulté à former une « volonté générale » unique au-delà des particularismes (danger des « associations particulières » que Rousseau redoutait).
\item \textbf{Corruption / détournement :} La loi formelle (légalité) ne garantit pas l'intérêt général (légitimité) si les élites la dévient.
\end{itemize}
\item \textbf{Conclusion :} La théorie reste le \textbf{critère normatif} de la légitimité : un État qui ne permet pas au peuple de se faire ses propres lois (élections libres, contrôle juridictionnel, liberté d'expression) perd sa légitimité rousseauiste. Au Cameroun, le chantier de la \textbf{décentralisation} et de la \textbf{participation citoyenne} (budgets participatifs, consultations) tente de rapprocher la décision de la volonté générale locale.
\end{itemize}

\textbf{Barème indicatif (sur 20) :} Thèse/Structure (4) + Explication aliénation (3) + Distinction VG/VdT (4) + Résolution paradoxe (4) + Discussion argumentée Cameroun (5).
\textbf{Points de vigilance :} Ne pas confondre « aliénation » (terme technique : don total) et « aliénation » (sens marxiste/perte de soi). Bien distinguer \emph{volonté générale} (normative, unitaire) et \emph{volonté de tous} (factuelle, multiple). La discussion doit être \emph{philosophique} (critère de légitimité) et non purement politique/journalistique.
\end{corrige}

\begin{corrige}[Exercice 11 — Analyse de document : dette publique et souveraineté]
\textbf{1. Évolution de la courbe (2010–2024)}
\begin{itemize}
\item \textbf{Tendance générale :} Hausse forte et continue de 2010 à 2022, suivie d'une légère baisse en 2024.
\item \textbf{2010–2015 :} Doublement rapide (15\% $\rightarrow$ 31\%, +16 pts). Début de l'endettement massif.
\item \textbf{2015–2020 :} Poursuite de la hausse (31\% $\rightarrow$ 45\%, +14 pts). Pic structurel.
\item \textbf{2020–2022 :} Stabilisation au niveau très élevé (45\% $\rightarrow$ 46\%, plateau).
\item \textbf{2022–2024 :} Amorce de décrue (46\% $\rightarrow$ 42\%, -4 pts), mais niveau reste supérieur au seuil de vigilance CEMAC (35\% ou 70\% selon critères).
\end{itemize}

\textbf{2. Calculs des taux de variation}
\begin{itemize}
\item \textbf{2010 $\rightarrow$ 2020 :} $\frac{45 - 15}{15} \times 100 = \frac{30}{15} \times 100 = \mathbf{+200\%}$. (L'endettement a triplé).
\item \textbf{2020 $\rightarrow$ 2024 :} $\frac{42 - 45}{45} \times 100 = \frac{-3}{45} \times 100 = \mathbf{-6,67\%}$. (Baisse modérée).
\end{itemize}

\textbf{3. Deux causes possibles}
\begin{enumerate}
\item \textbf{Chocs extérieurs conjugués :} Pandémie de Covid-19 (2020–2021) : chute des recettes pétrolières/cacao, hausse dépenses santé/social, plans de relance. Guerre en Ukraine (2022–) : inflation importée (carburant, engrais, blé), hausse des taux d'intérêt mondiaux (renchérissement du service de la dette), dépréciation du FCFA vs Dollar/Euro (dette libellée en devises).
\item \textbf{Investissements publics « structurants » (Doc 2) :} Stratégie « Vision 2035 » / PND30 : grands chantiers (barrages Nachtigal, Memve'ele ; autoroutes Yaoundé-Douala, Kribi-Lolabé ; port de Kribi ; stades CAN 2021/2025). Financés par emprunts extérieurs (Chine, France, Eurobonds) et intérieurs (Bons/Obass). Effet d'annonce vs rentabilité différée.
\end{enumerate}
\emph{(Autre cause : Insécurité — Boko Haram Extrême-Nord, crise NOSO Nord-Ouest/Sud-Ouest $\rightarrow$ dépenses militaires/sécuritaires non productives).}

\textbf{4. Risques sur la souveraineté économique}
\begin{itemize}
\item \textbf{Conditionnalité des bailleurs (FMI / BM) :} Programmes d'ajustement (FEC, MEC) imposant : gel salaires fonctionnaires, suppression subventions (carburant), hausse taxes, privatisations (SONARA, CAMTEL, CAMWATER), réduction dépenses sociales $\rightarrow$ \textbf{perte de décision budgétaire souveraine}.
\item \textbf{Effet d'éviction (crowding out) :} Service de la dette (intérêts + principal) absorbe \textgreater 30\% des recettes budgétaires $\rightarrow$ moins de marge pour éducation, santé, investissement productif.
\item \textbf{Dépendance aux créanciers :} Dette chinoise (opacité, clauses de garantie sur actifs stratégiques : port, mines, télécoms), Eurobonds (taux élevés, marchés). Risque de \textbf{saisie d'actifs stratégiques} en cas de défaut (ex: port de Mombasa Kenya, aéroport Hambantota Sri Lanka).
\item \textbf{Risque de défaut souverain :} Notation dégradée (Moody's, S\&P, Fitch) $\rightarrow$ fermeture marchés financiers, coût d'emprunt prohibitif, crise de change.
\end{itemize}

\textbf{5. Deux mesures de redressement et conception de l'État}
\begin{enumerate}
\item \textbf{Mesure : Rigueur budgétaire / Consolidation budgétaire (réduction déficit, plafonnement dette, revue dépenses, élargissement assiette fiscale sans hausse taux).}
\textbf{Conception : \textbf{État libéral (État minimal / État régulateur).}} L'État doit se recentrer sur ses fonctions régaliennes (sécurité, justice, monnaie), laisser le marché allouer les ressources, garantir la stabilité macroéconomique (inflation basse, comptes publics sains) pour attirer l'investissement privé. La dette est un fardeau pour les générations futures ; la discipline budgétaire est une vertu morale et économique (école de Fribourg, Hayek, Friedman).
\item \textbf{Mesure : Maintien/Accroissement des dépenses sociales et d'investissement productif financé par fiscalité progressive (impôt sur richesse, lutte évasion fiscale, renégociation dette) + politique industrielle active.}
\textbf{Conception : \textbf{État-providence / État keynésien / État développeur.}} L'État a la responsabilité de réduire les inégalités (services publics universels : santé, éducation, eau, électricité) et de piloter la transformation structurelle (industrialisation, souveraineté alimentaire/énergétique). La dette est un levier si elle finance des actifs productifs (infrastructures, capital humain). La souveraineté économique prime sur les ratios comptables (école de la régulation, Chang Ha-Joon, Stiglitz).
\end{enumerate}

\textbf{Barème indicatif (sur 20) :} Description courbe (3) + Calculs exacts (3) + Causes argumentées (4) + Risques souveraineté (4) + Mesures + Conceptions (6).
\textbf{Points de vigilance :} Taux de variation = $\frac{V_{final}-V_{init}}{V_{init}}\times 100$ (bien mettre 15 au dénominateur pour 2010-2020). Ne pas confondre dette publique \% PIB et dette en valeur absolue. Lier explicitement chaque mesure à une \emph{conception philosophique} de l'État (libéral vs providence/développeur), pas seulement à une technique comptable.
\end{corrige}

\newpage

\coursec{Fiche bilan}

\begin{popfigure}
\begin{tikzpicture}[
  mindmap,
  concept color=popblue,
  level 1/.style={level distance=3.5cm, sibling angle=360/6},
  level 2/.style={level distance=2.8cm, sibling angle=360/6},
  every node/.style={font=\small, text=white, align=center},
  branch1/.style={concept color=poporange},
  branch2/.style={concept color=popgreen},
  branch3/.style={concept color=poppurple},
  branch4/.style={concept color=poppink},
  branch5/.style={concept color=popgold},
  branch6/.style={concept color=popblue!70!black}
]
\node[concept, text width=2.8cm] {L'État\\ et le pouvoir}
  child[branch1] { node[concept, text width=2.2cm] {1. Définition\\ \& Éléments}
    child { node[concept, text width=1.8cm] {Population} }
    child { node[concept, text width=1.8cm] {Territoire} }
    child { node[concept, text width=1.8cm] {Souveraineté} }
    child { node[concept, text width=1.8cm] {Gouvernement} }
  }
  child[branch2] { node[concept, text width=2.2cm] {2. Fondements\\ \& Légitimité}
    child { node[concept, text width=1.8cm] {Force / Fait} }
    child { node[concept, text width=1.8cm] {Tradition} }
    child { node[concept, text width=1.8cm] {Droit / Contrat} }
    child { node[concept, text width=1.8cm] {Charisme (Weber)} }
    child { node[concept, text width=1.8cm] {Rationnel-légal} }
  }
  child[branch3] { node[concept, text width=2.2cm] {3. Formes de\\ gouvernement}
    child { node[concept, text width=1.8cm] {Monarchie / Aristocratie / Démocratie (Aristote)} }
    child { node[concept, text width=1.8cm] {Présidentiel / Parlementaire / Semi-présid.} }
    child { node[concept, text width=1.8cm] {Démocratique / Autoritaire / Totalitaire} }
    child { node[concept, text width=1.8cm] {Régime camerounais} }
  }
  child[branch4] { node[concept, text width=2.2cm] {4. Rôle \&\\ Finalité}
    child { node[concept, text width=1.8cm] {Fonctions régaliennes} }
    child { node[concept, text width=1.8cm] {État-providence} }
    child { node[concept, text width=1.8cm] {Économique / Social} }
    child { node[concept, text width=1.8cm] {Crise / Néolibéralisme} }
  }
  child[branch5] { node[concept, text width=2.2cm] {5. TP Intégré}
    child { node[concept, text width=1.8cm] {Cartographie institutionnelle} }
    child { node[concept, text width=1.8cm] {Analyse action publique locale} }
  }
  child[branch6] { node[concept, text width=2.2cm] {6. Dissertation\\ Bac}
    child { node[concept, text width=1.8cm] {Analyse sujet} }
    child { node[concept, text width=1.8cm] {Problématique} }
    child { node[concept, text width=1.8cm] {Plan dialectique} }
    child { node[concept, text width=1.8cm] {Exemples camerounais} }
  };
\end{tikzpicture}
\legende{Carte mentale de synthèse : Leçon 19 -- L'État et le pouvoir (Terminale Technique).}
\end{popfigure}

\begin{bilan}

\courssub{Définitions clés}
\begin{itemize}
  \item \cle{État} : Institution politique souveraine s'exerçant sur une population, un territoire, dotée d'un gouvernement et d'un monopole de la violence physique légitime (Weber).
  \item \cle{Souveraineté} : Pouvoir suprême, exclusif, inaliénable, imprescriptible (Bodin). Interne (autorité sur les sujets) et Externe (indépendance vis-à-vis des autres États).
  \item \cle{Légitimité} : Croyance en la validité de l'obéissance (Weber). Trois types pur : \emph{traditionnelle}, \emph{charismatique}, \cle{rationnelle-légale}.
  \item \cle{État de droit} : État soumis à la loi, séparation des pouvoirs, garantie des libertés fondamentales, justice indépendante.
  \item \cle{État-providence} : État intervenant dans le social (santé, éducation, retraite, chômage) pour corriger les inégalités du marché.
\end{itemize}

\courssub{Repères théoriques (Auteurs / Thèses)}
\begin{center}
\begin{tabularx}{\linewidth}{|l|X|l|}
\hline
\rowcolor{popblueL}
\textbf{Auteur / École} & \textbf{Thèse centrale} & \textbf{Notion clé} \\
\hline
Aristote & Classification régimes : Monarchie/Tyrannie, Aristocratie/Oligarchie, Démocratie/Démagogie. Fin de l'État = vie bonne. & \cle{Finalité politique} \\
Machiavel & Primat de l'efficacité : « La fin justifie les moyens ». Séparation politique / morale. & \cle{Raison d'État} \\
Bodin / Hobbes & Souveraineté absolue, indivisible. Contrat social pour sortir de l'état de nature (guerre de tous contre tous). & \cle{Souveraineté / Léviathan} \\
Locke / Montesquieu & Limitation du pouvoir : Séparation des pouvoirs (Législatif, Exécutif, Judiciaire), droit de résistance. & \cle{État de droit} \\
Rousseau & Souveraineté populaire, inaliénable, indivisible. Volonté générale $\neq$ volonté de tous. & \cle{Démocratie / Contrat} \\
Marx / Engels & État = instrument de domination de classe (superstructure). Dépérissement de l'État en communisme. & \cle{Lutte des classes} \\
Weber & Monopole de la violence physique légitime. Typologie de la domination (légitimité). Bureaucratie. & \cle{Légitimité / Bureaucratie} \\
Rawls & Justice comme équité. Principes de liberté et de différence. Légitimité par le consentement rationnel. & \cle{Justice sociale} \\
\hline
\end{tabularx}
\end{center}

\courssub{Formes de gouvernement \& Régimes (Résumé)}
\begin{itemize}
  \item \textbf{Critère aristotélicien} : Nombre de gouvernants (Un / Quelques-uns / Tous) $\times$ Finalité (Intérêt général / Intérêt privé).
  \item \textbf{Critère institutionnel moderne} : Relations Exécutif / Législatif.
    \begin{itemize}
      \item \cle{Présidentiel} : Élection au suffrage universel, séparation stricte, irresponsabilité politique du Président (ex: USA, Cameroun).
      \item \cle{Parlementaire} : Exécutif bicéphale (Chef d'État / Chef de gouvernement), responsabilité politique du gouvernement devant le parlement (ex: UK, Allemagne).
      \item \cle{Semi-présidentiel} : Président élu + Premier ministre responsable devant le parlement (ex: France).
    \end{itemize}
  \item \textbf{Critère libertés} : Démocratique (pluralisme, élections libres, DDH) / Autoritaire (pouvoir personnel, libertés limitées) / Totalitaire (contrôle total société, idéologie unique, parti unique, terreur).
\end{itemize}

\courssub{Focus Cameroun : Régime politique}
\begin{itemize}
  \item République unitaire décentralisée (Constitution 1996, révisée 2008).
  \item Régime \cle{présidentiel} : Président de la République élu au suffrage universel direct (mandat 7 ans, renouvelable une fois depuis 2008), Chef de l'État, Chef de l'Exécutif, nomme le PM et le gouvernement.
  \item Parlement bicaméral : Assemblée nationale (députés) + Sénat (sénateurs).
  \item Décentralisation : Collectivités territoriales décentralisées (CTD) -- Communes, Régions.
  \item Multipartisme (depuis 1990), mais dominance historique du RDPC.
\end{itemize}

\courssub{Méthode Express : Dissertation Philo (Bac Tech)}
\begin{enumerate}
  \item \textbf{Analyser} : Définir les termes, identifier le type de sujet (Vrai/Faux, « Peut-on », « Faut-il », Citation).
  \item \textbf{Problématiser} : Trouver la tension (ex: Autorité vs Liberté, Efficacité vs Justice, Fait vs Droit). Formuler une question précise.
  \item \textbf{Plan dialectique (Thèse / Antithèse / Synthèse)} ou \textbf{analytique} (Concept / Auteurs / Application).
  \item \textbf{Argumenter} : 1 idée = 1 paragraphe (Thèse + Argument + Exemple/Ref auteur + Lien). \textbf{Exemples camerounais obligatoires} (Constitution, lois, actualité locale, institutions).
  \item \textbf{Conclure} : Réponse nuancée, ouverture.
\end{enumerate}

\end{bilan}

\begin{aretenir}[Checklist avant le Bac -- Leçon 19]
$\square$ Je sais définir l'État par ses 4 éléments constitutifs (Population, Territoire, Souveraineté, Gouvernement).\\
$\square$ Je distingue \emph{souveraineté} (pouvoir juridique) et \emph{légitimité} (croyance/acceptation).\\
$\square$ Je connais les 3 types de légitimité weberiens (Traditionnelle, Charismatique, Rationnelle-légale) et donne un exemple chacun.\\
$\square$ Je sais classer les régimes selon Aristote (6 formes) et selon le critère moderne (Présidentiel, Parlementaire, Semi-présidentiel).\\
$\square$ Je caractérise le régime politique du Cameroun (Présidentiel, Unitaire, Décentralisé, Bicaméral) et cite les articles clés de la Constitution.\\
$\square$ Je distingue Fonctions régaliennes (Police, Justice, Armée, Diplomatie) et Fonctions sociales/économiques (État-providence).\\
$\square$ J'explique la crise de l'État-providence (coût, dépendance, efficacité) et les réponses néolibérales (dérégulation, privatisation).\\
$\square$ Je sais mobiliser au moins 3 auteurs (ex: Hobbes, Locke, Rousseau, Marx, Weber, Rawls) avec leur thèse précise sur l'État.\\
$\square$ Je sais construire une problématique à partir d'un sujet type Bac (ex: « L'État est-il au service de la justice ? »).\\
$\square$ Je dispose d'exemples concrets camerounais (Décentralisation, PNDP, MINAS, ELECAM, Conseil Constitutionnel, Grèves, Élections) pour illustrer mes arguments.\\
$\square$ Je maîtrise la structure de la dissertation (Intro : Amorce, Définitions, Problématique, Annonce du plan / Développement équilibré / Conclusion : Bilan, Réponse, Ouverture).\\
$\square$ Je relis ma copie pour corriger la syntaxe, l'orthographe des concepts clés et la cohérence logique.
\end{aretenir}

\findecours

\end{document}
